\documentclass[letterpaper, journal]{IEEEtran}
%DIF LATEXDIFF DIFFERENCE FILE
%DIF DEL tmp2.tex   Wed May 27 18:53:56 2026
%DIF ADD tmp1.tex   Wed May 27 18:53:56 2026

%DIF 3-5c3-5
%DIF < \def\cmt{false} % true for comments, false for no comments
%DIF < \def\pub{true} % true for publication, false for draft
%DIF < \newcommand*{\template}{../template} 
%DIF -------
\def\cmt{true} % true for comments, false for no comments %DIF > 
\def\pub{false} % true for publication, false for draft %DIF > 
\newcommand*{\template}{template}  %DIF > 
%DIF -------
\input{\template/preamble/preamble_journal.tex}
%DIF 7c7
%DIF < \newcommand*{\src}{../src} 
%DIF -------
\allowdisplaybreaks %DIF > 
%DIF -------

\newcommand*{\figSizeTwoCol}{.49}
\newcommand*{\figSizeOneCol}{0.98}
 %DIF > 
\newcommand{\NNin}{\mv{x}_{nn}} %DIF > 
 %DIF > 
\allowdisplaybreaks %DIF > 
 %DIF > 
\newcommand{\src}{./src} %DIF > 
%DIF PREAMBLE EXTENSION ADDED BY LATEXDIFF
%DIF UNDERLINE PREAMBLE %DIF PREAMBLE
\RequirePackage[normalem]{ulem} %DIF PREAMBLE
\RequirePackage{color}\definecolor{RED}{rgb}{1,0,0}\definecolor{BLUE}{rgb}{0,0,1} %DIF PREAMBLE
\providecommand{\DIFadd}[1]{{\protect\color{blue}\uwave{#1}}} %DIF PREAMBLE
\providecommand{\DIFdel}[1]{{\protect\color{red}\sout{#1}}} %DIF PREAMBLE
%DIF SAFE PREAMBLE %DIF PREAMBLE
\providecommand{\DIFaddbegin}{} %DIF PREAMBLE
\providecommand{\DIFaddend}{} %DIF PREAMBLE
\providecommand{\DIFdelbegin}{} %DIF PREAMBLE
\providecommand{\DIFdelend}{} %DIF PREAMBLE
\providecommand{\DIFmodbegin}{} %DIF PREAMBLE
\providecommand{\DIFmodend}{} %DIF PREAMBLE
%DIF FLOATSAFE PREAMBLE %DIF PREAMBLE
\providecommand{\DIFaddFL}[1]{\DIFadd{#1}} %DIF PREAMBLE
\providecommand{\DIFdelFL}[1]{\DIFdel{#1}} %DIF PREAMBLE
\providecommand{\DIFaddbeginFL}{} %DIF PREAMBLE
\providecommand{\DIFaddendFL}{} %DIF PREAMBLE
\providecommand{\DIFdelbeginFL}{} %DIF PREAMBLE
\providecommand{\DIFdelendFL}{} %DIF PREAMBLE
%DIF COLORLISTINGS PREAMBLE %DIF PREAMBLE
\RequirePackage{listings} %DIF PREAMBLE
\RequirePackage{color} %DIF PREAMBLE
\lstdefinelanguage{DIFcode}{ %DIF PREAMBLE
%DIF DIFCODE_UNDERLINE %DIF PREAMBLE
  moredelim=[il][\color{red}\sout]{\%DIF\ <\ }, %DIF PREAMBLE
  moredelim=[il][\color{blue}\uwave]{\%DIF\ >\ } %DIF PREAMBLE
} %DIF PREAMBLE
\lstdefinestyle{DIFverbatimstyle}{ %DIF PREAMBLE
	language=DIFcode, %DIF PREAMBLE
	basicstyle=\ttfamily, %DIF PREAMBLE
	columns=fullflexible, %DIF PREAMBLE
	keepspaces=true %DIF PREAMBLE
} %DIF PREAMBLE
\lstnewenvironment{DIFverbatim}[1][]{\lstset{style=DIFverbatimstyle}}{} %DIF PREAMBLE
\lstnewenvironment{DIFverbatim*}[1][]{\lstset{style=DIFverbatimstyle,showspaces=true}}{} %DIF PREAMBLE
\lstset{extendedchars=\true,inputencoding=utf8}

%DIF END PREAMBLE EXTENSION ADDED BY LATEXDIFF

\begin{document}
\bstctlcite{IEEEexample:BSTcontrol}

\title{
    Constrained Optimization-Based Neuro-Adaptive Control (CONAC) for \DIFdelbegin \DIFdel{Euler-Lagrange }\DIFdelend \DIFaddbegin \DIFadd{Unknown }\DIFaddend Systems Under \DIFdelbegin \DIFdel{Weight and }\DIFdelend \DIFaddbegin \DIFadd{Multiple Convex }\DIFaddend Input Constraints
} %% Article title

\author{
    Myeongseok Ryu, Donghwa Hong, and Kyunghwan Choi
\DIFdelbegin %DIFDELCMD < \thanks{
%DIFDELCMD <     Myeongseok Ryu and Kyunghwan Choi are with the Cho Chun Shik Graduate School of Mobility, Korea Advanced Institute of Science and Technology (KAIST), Daejeon 34051, Korea (e-mail: dding\_98@kaist.ac.kr and kh.choi@kaist.ac.kr).
%DIFDELCMD < 

%DIFDELCMD <     Donghwa Hong is with the Department of Mechanical and Robotics Engineering, Gwangju Institute of Science and Technology (GIST), Gwangju 61005, Korea (e-mail: fairytalef00@gm.gist.ac.kr).
%DIFDELCMD < 

%DIFDELCMD <     This work was supported in part by a Korea Research Institute for Defense Technology planning and advancement (KRIT) grant funded by Korea government DAPA (Defense Acquisition Program Administration) (No. KRIT-CT-22-087, Synthetic Battlefield Environment based Integrated Combat Training Platform) and in part by the National Research Foundation of Korea (NRF) grant funded by the Korea government (MSIT) (RS-2025- 00554087). (Corresponding author: Kyunghwan Choi)
%DIFDELCMD <     }%%%
\DIFdelend \DIFaddbegin \thanks{
    Myeongseok Ryu, Donghwa Hong and Kyunghwan Choi are with the Cho Chun Shik Graduate School of Mobility, Korea Advanced Institute of Science and Technology (KAIST), Daejeon 34051, Korea (e-mail: \{dding\_98, fairytale, kh.choi\}@kaist.ac.kr).

    This work was supported in part by a Korea Research Institute for Defense Technology planning and advancement (KRIT) grant funded by Korea government DAPA (Defense Acquisition Program Administration) (No. KRIT-CT-22-087, Synthetic Battlefield Environment based Integrated Combat Training Platform) and in part by the National Research Foundation of Korea (NRF) grant funded by the Korea government (MSIT) (RS-2025- 00554087). (Corresponding author: Kyunghwan Choi)
    }\DIFaddend }

\markboth{Journal of \LaTeX\ Class Files,~Vol.~18, No.~9, September~2020}%
{How to Use the IEEEtran \LaTeX \ Templates}

\maketitle

%% Abstract
\begin{abstract}
  This study presents a constrained optimization-based neuro-adaptive control (CONAC) for unknown \DIFdelbegin \DIFdel{Euler-Lagrange }\DIFdelend systems subject to \DIFdelbegin \DIFdel{weight and }\DIFdelend \DIFaddbegin \DIFadd{multiple }\DIFaddend convex input constraints. 
  A deep neural network (DNN) is employed to approximate the ideal stabilizing control law while \DIFdelbegin \DIFdel{simultaneously learning the unknown system dynamics and addressing both types of }\DIFdelend \DIFaddbegin \DIFadd{handling multiple convex input }\DIFaddend constraints within a unified constrained optimization framework.
  The \DIFdelbegin \DIFdel{adaptation law of DNN weights is formulated to solve the defined }\DIFdelend \DIFaddbegin \DIFadd{adaptive variables including the DNN weights and Lagrange multipliers are updated through adaptation laws derived from the formulated }\DIFaddend constrained optimization problem, \DIFdelbegin \DIFdel{ensuring satisfaction of }\DIFdelend \DIFaddbegin \DIFadd{yielding KKT-like }\DIFaddend first-order optimality conditions at \DIFdelbegin \DIFdel{steady state}\DIFdelend \DIFaddbegin \DIFadd{equilibrium}\DIFaddend .
  The controller's stability is rigorously analyzed using Lyapunov theory, guaranteeing \DIFaddbegin \DIFadd{uniformly ultimately }\DIFaddend bounded tracking errors\DIFdelbegin \DIFdel{and bounded DNN weights}\DIFdelend \DIFaddbegin \DIFadd{, and adaptive variables}\DIFaddend .
  The proposed controller is \DIFaddbegin \DIFadd{compared with existing methods in numerical simulations, highlighting its capability to handle multiple convex input constraints.
  Additionally, the feasibility and effectiveness of CONAC in real-world applications are }\DIFaddend validated through a real-time implementation on a 2-DOF robotic manipulator\DIFdelbegin \DIFdel{, demonstrating its effectiveness in achieving accurate trajectory tracking while satisfying all imposed constraints}\DIFdelend .
\end{abstract}

\begin{IEEEkeywords}
    Neuro-adaptive control, constrained optimization, deep neural network,  \DIFdelbegin \DIFdel{input constraint, weight constraint}\DIFdelend \DIFaddbegin \DIFadd{multiple convex input constraints}\DIFaddend .
\end{IEEEkeywords}

\section*{Notation}

%DIF <  In this study, the following notation is used:
\DIFdelbegin %DIFDELCMD < 

%DIFDELCMD < %%%
%DIF <  \begin{itemize}
%DIF <      \item := denotes ``is defined as" or ``is equal to".
%DIF <      \item $\R$, $\R^n$ and $\R^{n\times m}$ denote the set of real numbers, $n\times 1$ dimensional vectors and $n\times m$ dimensional matrices, respectively, and $\R_{\ge0}:=[0,\infty)$ and $\R_{>0}:=(0,\infty)$.
%DIF <      \item $(\cdot)^\top$ denotes the transpose of a vector or matrix.
%DIF <      \item $\norm{\cdot}$ and $\norm{\cdot}_F$ denote the Euclidean and Frobenius norm, respectively.
%DIF <      \item $\otimes$ denotes the Kronecker product \cite[Chap. 7, Def. 7.1.2]{Bernstein:2009aa}.
%DIF <      \item A vector and a matrix are denoted by $\mv{x}=[x_i]_{i\in\{1,\cdots,n\}}\in\R^n$ and $
%DIF <          \mm A
%DIF <          := 
%DIF <          [a_{ij}]
%DIF <          _{
%DIF <              i\in\{1,\cdots,n\},j\in\{1,\cdots ,m\}
%DIF <          }\allowbreak\in\R^{n\times m}
%DIF <          $, respectively.
%DIF <      \item $\myrow_i(\mm A)$ denotes the $i$th row of the matrix $\mm A\in\R^{n\times m}$. 
%DIF <      \item For the matrix $\mm{A}\in\R^{n\times m}$, $\myvec(\mm A):=(\myrow_1(\mm A^\top),\cdots,\myrow_m(\mm A^\top) )^\top\in\R^{nm}$ denotes the vectorization of $\mm A$.
%DIF <      % \item $\mycol_i(A)$ and $\myrow_j(A)$ denote the $i$th column and $j$th row of $A\in\R^{n\times m}$, respectively.
%DIF <      % \item $\myvec(A):= [\mycol_1(A)^\top  ,\cdots,\mycol_m(A)^\top  ]^\top   $ for $A\in\R^{n\times m}$.
%DIF <      \item $\lambda_{\min}(\mm A)$ denotes the minimum eigenvalue of the matrix $\mm A\in\R^{n\times n}$.
%DIF <      \item $\mm I_n$ denotes the $n\times n$ identity matrix, and $\mm 0_{n\times m}$ denotes the $n\times m$ zero matrix.
%DIF <  \end{itemize}
%DIFDELCMD < 

%DIFDELCMD < %%%
\DIFdel{The Kronecker product is denoted by }\DIFdelend \DIFaddbegin \DIFadd{In this study, the following notation is used: 
}\DIFaddend $\otimes$ \DIFaddbegin \DIFadd{denotes the Kronecker product }\DIFaddend \cite[Chap. 7, Def. 7.1.2]{Bernstein:2009aa}.
A vector and a matrix are denoted by $\mv{x}=[x_i]_{i\in\{1,\cdots,n\}}\in\R^n$ and 
$
  \mm A
  := 
  [a_{ij}]
  _{
      i\in\{1,\cdots,n\},j\in\{1,\cdots ,m\}
  }\allowbreak\in\R^{n\times m}
$, respectively.
\DIFdelbegin \DIFdel{Given a matrix $\mm A\in\R^{n\times m}$, $\myrow_i(\mm A)$ denotes the }\DIFdelend \DIFaddbegin \DIFadd{The }\DIFaddend $i$th row of \DIFdelbegin \DIFdel{$\mm A$}\DIFdelend \DIFaddbegin \DIFadd{the matrix $\mm{A}\in\R^{n\times m}$ is denoted by $\myrow_i(\mm A)$}\DIFaddend .
For the matrix $\mm{A}\in\R^{n\times m}$, $\myvec(\mm A):=(\myrow_1(\mm A^\top),\cdots,\myrow_m(\mm A^\top) )^\top\in\R^{nm}$ denotes the vectorization of $\mm A$.
\DIFdelbegin \DIFdel{$\lambda_{\min}(\mm A)$ denotes the minimum eigenvalue of the matrix $\mm A\in\R^{n\times n}$.
The }\DIFdelend \DIFaddbegin \DIFadd{Finally, $\mm I_n$ denotes the }\DIFaddend $n\times n$ identity \DIFdelbegin \DIFdel{and }\DIFdelend \DIFaddbegin \DIFadd{matrix, and $\mm 0_{n\times m}$ denotes the }\DIFaddend $n\times m$ zero \DIFdelbegin \DIFdel{matrices are denoted by $\mm I_n$ and $\mm 0_{n\times m}$, respectively}\DIFdelend \DIFaddbegin \DIFadd{matrix}\DIFaddend .

%DIF <   SECTION INTRODUCTION ===================================
%DIF >  In this study, the following notation is used:
\DIFaddbegin 

%DIF >  \begin{itemize}
%DIF >      % \item := denotes ``is defined as" or ``is equal to".
%DIF >      % \item $\R$, $\R^n$ and $\R^{n\times m}$ denote the set of real numbers, $n\times 1$ dimensional vectors and $n\times m$ dimensional matrices, respectively, and $\R_{\ge0}:=[0,\infty)$ and $\R_{>0}:=(0,\infty)$.
%DIF >      % \item $(\cdot)^\top$ denotes the transpose of a vector or matrix.
%DIF >      % \item $\norm{\cdot}$ and $\norm{\cdot}_F$ denote the Euclidean and Frobenius norm, respectively.
%DIF >      \item $\otimes$ denotes the Kronecker product \cite[Chap. 7, Def. 7.1.2]{Bernstein:2009aa}.
%DIF >      \item A vector and a matrix are denoted by $\mv{x}=[x_i]_{i\in\{1,\cdots,n\}}\in\R^n$ and $
%DIF >          \mm A
%DIF >          := 
%DIF >          [a_{ij}]
%DIF >          _{
%DIF >              i\in\{1,\cdots,n\},j\in\{1,\cdots ,m\}
%DIF >          }\allowbreak\in\R^{n\times m}
%DIF >          $, respectively.
%DIF >      \item $\myrow_i(\mm A)$ denotes the $i$th row of the matrix $\mm A\in\R^{n\times m}$. 
%DIF >      \item For the matrix $\mm{A}\in\R^{n\times m}$, $\myvec(\mm A):=(\myrow_1(\mm A^\top),\cdots,\myrow_m(\mm A^\top) )^\top\in\R^{nm}$ denotes the vectorization of $\mm A$.
%DIF >      % \item $\mycol_i(A)$ and $\myrow_j(A)$ denote the $i$th column and $j$th row of $A\in\R^{n\times m}$, respectively.
%DIF >      % \item $\myvec(A):= [\mycol_1(A)^\top  ,\cdots,\mycol_m(A)^\top  ]^\top   $ for $A\in\R^{n\times m}$.
%DIF >      % \item $\lambda_{\min}(\mm A)$ denotes the minimum eigenvalue of the matrix $\mm A\in\R^{n\times n}$.
%DIF >      \item $\mm I_n$ denotes the $n\times n$ identity matrix, and $\mm 0_{n\times m}$ denotes the $n\times m$ zero matrix.
%DIF >  \end{itemize}

%DIF >  ------------------------------------
%DIF >        SECTION INTRODUCTION 
%DIF >  ------------------------------------
%DIF >  \vspace{100mm}                
\DIFaddend \section{Introduction}

\subsection{Background}

\DIFdelbegin %DIFDELCMD < \IEEEPARstart{M}{any} %%%
\DIFdel{engineering systems , including those in aerospace, robotics, and automotive applications, can }\DIFdelend \DIFaddbegin \IEEEPARstart{F}{rom} \DIFadd{a practical perspective, many engineering systems are constrained by multiple safety requirements, actuator limitations, and physical restrictions \mbox{%DIFAUXCMD
\cite{Ji:2025aa}}\hskip0pt%DIFAUXCMD
.
These constraints often manifest as convex input constraints, which can be in the form of nonlinear constraints rather than simple box constraints that limit each input independently.
For example, under road friction limits, the admissible global force and moment of vehicles can be characterized by weighted norm bound, which yields convex ellipsoidal constraint \mbox{%DIFAUXCMD
\cite{Chen:2024aa}}\hskip0pt%DIFAUXCMD
.
%DIF >  \MSRY{전력원의 한계로 인해서 토크 놈에 대해서 제약조건을 주는게 명시적으로 논문에 나온게 없더라구요.. 그래서 아래 문장은 삭제되어야 할 것 같습니다.}
%DIF >  \cred{
%DIF >  For instance, due to power limitations for actuators, robotic manipulators may have constraints on the upper bounds on the norm of the joint torques.
%DIF >  }
In motor drive systems, the feasible voltage region is often represented by circular \mbox{%DIFAUXCMD
\cite{Ryu:2025aa} }\hskip0pt%DIFAUXCMD
or hexagonal sets \mbox{%DIFAUXCMD
\cite{Monzen:2026aa}}\hskip0pt%DIFAUXCMD
.
%DIF >  More recently, input constraints induced by state constraints using control barrier functions (CBFs) \cite{Ames:2019aa} are inherently nonlinear and complex.
Furthermore, admissible input sets may }\DIFaddend be \DIFdelbegin \DIFdel{modeled using Euler-Lagrange systems. 
These systems are governed by dynamic equations derived from energy principles and describe the motion of mechanical systems with constraints }\DIFdelend \DIFaddbegin \DIFadd{determined by multiple constraints that must hold simultaneously, such as component-wise bounds combined with norm constraints, which can be represented as the intersection of multiple convex sets.
Therefore, addressing these multiple convex input constraints is crucial for ensuring the safety and reliability of control systems in real-world applications}\DIFaddend .
\DIFdelbegin \DIFdel{In practice, however, such systems often exhibit uncertainties due to unmodeled dynamics, parameter variations, or external disturbances. 
These uncertainties can significantly degrade control performance and , in some cases, lead to instability.
In the worst cases, any prior knowledge of the system dynamics may not be available, making it challenging to design effective controllers.
%DIF <  To address these challenges, adaptive control methods have been widely employed to ensure robust performance in the presence of system uncertainties \cite{Ioannou:2006aa, Tao:2003aa}.
To address these challenges, adaptive control methods have been widely employed to ensure robust performance for uncertain or unknown systems \mbox{%DIFAUXCMD
\cite{Ioannou:2006aa, Tao:2003aa}}\hskip0pt%DIFAUXCMD
.
}\DIFdelend 

\DIFdelbegin \DIFdel{More recently, neuro-adaptive control (NAC) approaches have been introduced to approximate uncertain or unknown system dynamicsor entire control laws using neural networks (NNs) \mbox{%DIFAUXCMD
\cite{Lewis:1998aa,Farrell:2006aa}}\hskip0pt%DIFAUXCMD
. 
NNs are well-known for their universal approximation property, which allows them to approximate any smooth function over a compact set with minimal error. 
Various types of NNs have been utilized in NAC, including simpler architectures like single-hidden layer (SHL) NNs \mbox{%DIFAUXCMD
\cite{Lewis:1996aa,Ge:2010aa, Yesildirek:1995aa,Ryu:2024aa,Ryu:2025aa} }\hskip0pt%DIFAUXCMD
and radial basis function (RBF) NNs \mbox{%DIFAUXCMD
\cite{Liu:2013ab,Ge:2002aa}}\hskip0pt%DIFAUXCMD
, as well as more complex models like deep NNs (DNNs) \mbox{%DIFAUXCMD
\cite{Patil:2022aa} }\hskip0pt%DIFAUXCMD
and their variations.
Conventionally, SHL and RBF NNs are often employed to approximate }\DIFdelend \DIFaddbegin \DIFadd{In parallel, another challenge in engineering systems is uncertain dynamics or even completely unknown dynamics, which can arise from modeling errors, unmodeled dynamics, and external disturbances.
Although adaptive control provides a systematic mechanism for online adaptation of controllers \mbox{%DIFAUXCMD
\cite{Ioannou:2006aa}}\hskip0pt%DIFAUXCMD
, conventional formulations often rely on structural assumptions on the }\DIFaddend uncertain or unknown \DIFdelbegin \DIFdel{system dynamicsor controllers due to their simplicity \mbox{%DIFAUXCMD
\cite{Esfandiari:2014aa,Esfandiari:2015aa,Yesildirek:1995aa,Gao:2006aa}}\hskip0pt%DIFAUXCMD
.
However, since DNNs offer greater expressive power, making them more effective for complex system approximations \mbox{%DIFAUXCMD
\cite{Rolnick:2018aa}}\hskip0pt%DIFAUXCMD
, variations of DNNs, such as long short-term memory (LSTM) networks for time-varying dynamics \mbox{%DIFAUXCMD
\cite{Griffis:2023aa}}\hskip0pt%DIFAUXCMD
, physics-informed NNs (PINNs) for leveraging physical system knowledge \mbox{%DIFAUXCMD
\cite{Hart:2024aa}}\hskip0pt%DIFAUXCMD
, and convolutional NNs(CNNs) for learning from historical sensor data \mbox{%DIFAUXCMD
\cite{Ryu:2024ac}}\hskip0pt%DIFAUXCMD
, have further extended the capabilities of NAC systems.
}\DIFdelend \DIFaddbegin \DIFadd{dynamics, which may not be available in practice.
Neuro-adaptive control (NAC) alleviates this limitation by using neural networks (NNs) to approximate the dynamics in a structured and parameterized representation \mbox{%DIFAUXCMD
\cite{Farrell:2006aa,Patil:2022aa}}\hskip0pt%DIFAUXCMD
.
Like adaptive control, NAC methods ensure closed-loop stability using Lyapunov theory, and the NN weights are adapted online without prior knowledge of the system dynamics.
The NAC methods have shown promising results in various real-world applications, including flexible robotic manipulators \mbox{%DIFAUXCMD
\cite{Gao:2026aa}}\hskip0pt%DIFAUXCMD
, active suspensions \mbox{%DIFAUXCMD
\cite{Zhao:2016aa}}\hskip0pt%DIFAUXCMD
, synchronous machine drives \mbox{%DIFAUXCMD
\cite{Ryu:2025aa}}\hskip0pt%DIFAUXCMD
, and swarm control of mobile robots \mbox{%DIFAUXCMD
\cite{Yu:2022aa}}\hskip0pt%DIFAUXCMD
.
}\DIFaddend 

\DIFdelbegin \DIFdel{A critical aspect of NAC is the NN weight adaptation law, which governs how NN weights are updated. 
Most studies derived these laws using Lyapunov-based methods, ensuring the boundedness of the tracking error and weight estimation error, thus maintaining system stability under uncertainty.
}%DIFDELCMD < 

%DIFDELCMD < %%%
\DIFdel{However, two significant challenges persist in using NNs for adaptive control. 
First, the boundedness of NN weights is not inherently guaranteed, which can result in unbounded outputs. 
When NN outputs are used directly in the control law, this may lead to excessive control inputs, violating input constraints. 
Such constraints are commonly encountered in industrial systems, where actuators are limited by physical and safety requirements in terms of amplitude, rate, or energy \mbox{%DIFAUXCMD
\cite{Esfandiari:2021aa}}\hskip0pt%DIFAUXCMD
.
Failing to address these constraints can degrade control performance or even destabilize the system.
}%DIFDELCMD < 

%DIFDELCMD < %%%
\DIFdel{Therefore, addressing these two key issues—ensuring weight boundedness and satisfying input constraints—is essentialfor the reliable design of NAC.
The following section provides a detailed }\DIFdelend \DIFaddbegin \DIFadd{Because multiple convex input constraints and uncertain or unknown dynamics often coexist in real-world applications, incorporating input constraint handling into NAC is essential.
This motivates a }\DIFaddend review of existing \DIFdelbegin \DIFdel{solutions to these challenges}\DIFdelend \DIFaddbegin \DIFadd{input constraint handling methods in NAC and related constrained control frameworks}\DIFaddend .

\subsection{Literature Review}

\DIFdelbegin \subsubsection{\DIFdel{Ensuring Weight Norm Boundedness}}
%DIFAUXCMD
\addtocounter{subsubsection}{-1}%DIFAUXCMD
\DIFdelend %DIF >  ================================
%DIF >  STRUCTURAL METHODS
%DIF >  ================================

\DIFdelbegin \DIFdel{A common challenge in NACis maintaining the boundedness of the NN weights to ensure stability . 
In many studies, projection operators were employed to enforce upper bounds on the weight norms, ensuring that the weights do not grow unboundedly.
For example, in \mbox{%DIFAUXCMD
\cite{Zhou:2023aa,Griffis:2023aa,Patil:2022aa}}\hskip0pt%DIFAUXCMD
, projection operators were used to constrain the weight norms to remain below predetermined constants.
However, these constants were often selected as large as possible due to the lack of theoretical guarantees regarding the global optimality of the weight values.
While this approach ensured that the NN remained stable, it did not necessarily result in optimal performance.
}\DIFdelend \DIFaddbegin \subsubsection{\DIFadd{Structural Methods for Input Constraint Handling in NAC}}
\DIFaddend 

\DIFdelbegin \DIFdel{In addition to projection operators, some studies utilized modification techniques like $\sigma$-modification \mbox{%DIFAUXCMD
\cite{Ge:2002aa} }\hskip0pt%DIFAUXCMD
and $\epsilon$-modification \mbox{%DIFAUXCMD
\cite{Esfandiari:2015aa,Gao:2006aa}}\hskip0pt%DIFAUXCMD
. 
These methods ensured that the NN weights remained within an invariant set by incorporating stabilizing functions into the adaptation law.
Although these techniques were effective in ensuring boundedness and avoiding weight divergence, they similarly lacked a formal analysis of the optimality of the adapted weights by biasing the weights toward the origin \mbox{%DIFAUXCMD
\cite{Ryu:2024aa}}\hskip0pt%DIFAUXCMD
.
This leaves room for improvement in terms of performance optimization.
}\DIFdelend \DIFaddbegin \DIFadd{In NAC, structural methods have been mainly developed for input constraint handling, because they preserve the continuous-time controller structure and are amenable to Lyapunov stability analysis.
Auxiliary-system-based approaches are the most widely used method owing to their simplicity and effectiveness \mbox{%DIFAUXCMD
\cite{Esfandiari:2014aa,Esfandiari:2015aa,Chen:2011aa,He:2016aa,Zhao:2016aa,Gao:2026aa,Yu:2022aa,Liu:2025aa,Ma:2025aa}}\hskip0pt%DIFAUXCMD
.
When input saturation occurs, the auxiliary state is excited and evolves according to the auxiliary system, whose dynamics is linear \mbox{%DIFAUXCMD
\cite{Yu:2022aa}}\hskip0pt%DIFAUXCMD
, nonlinear \mbox{%DIFAUXCMD
\cite{Chen:2011aa,Gao:2026aa,Ma:2025aa} }\hskip0pt%DIFAUXCMD
or fixed-time-inspired \mbox{%DIFAUXCMD
\cite{Liu:2025aa}}\hskip0pt%DIFAUXCMD
.
Typically, the excited auxiliary state is directly fed back into the control law as additional compensation terms.
}\DIFaddend 

\DIFdelbegin %DIFDELCMD < \hfill
%DIFDELCMD < 

%DIFDELCMD < %%%
\subsubsection{\DIFdel{Satisfying Input Constraints}}
%DIFAUXCMD
\addtocounter{subsubsection}{-1}%DIFAUXCMD
%DIFDELCMD < 

%DIFDELCMD < %%%
\DIFdel{The second major issue is satisfying input constraints , particularly in systems where actuators are subject to physical limitations. 
This issue becomes more pronounced in NAC architectures where NNs are used to augment conventional model-based controllers—such as feedback linearization or backstepping—rather than acting as the sole control input }\DIFdelend \DIFaddbegin \DIFadd{However, from a learning perspective, introducing additional compensation terms into the control law can interfere with the NN adaptation process.
This occurs because the adaptation law typically depends on the feedback error, which may become less informative for updating the NN weights when the additional terms dominate the closed-loop response.
To address this issue, some studies incorporated the auxiliary state into the feedback error used in adaptation laws \mbox{%DIFAUXCMD
\cite{Esfandiari:2014aa,Esfandiari:2015aa}}\hskip0pt%DIFAUXCMD
.
This approach allows the NN weights to adapt not only for approximating the ideal control law but also for compensating the effect of input saturation}\DIFaddend .
\DIFdelbegin \DIFdel{In such cases, if the system model is inaccurate or the dynamics targeted by the NN are already contributing to stability, }\DIFdelend \DIFaddbegin 

\DIFadd{Additionally, several other structural methods have been introduced for handling input constraints in NAC. 
For instance, in \mbox{%DIFAUXCMD
\cite{Gao:2006aa}}\hskip0pt%DIFAUXCMD
, the saturated amount of control input is estimated using NNs and compensated in the control law.
In \mbox{%DIFAUXCMD
\cite{Li:2024aa}}\hskip0pt%DIFAUXCMD
, }\DIFaddend the \DIFdelbegin \DIFdel{combined control effort may result in overly aggressive inputs , increasing the risk of input saturation \mbox{%DIFAUXCMD
\cite{Khalil:2002aa}}\hskip0pt%DIFAUXCMD
}\DIFdelend \DIFaddbegin \DIFadd{control law is augmented with a smooth saturation function, such as $\tanh(\cdot)$, to accommodate input saturation effects, while the closed-loop stability is analyzed using Lyapunov theory.
For strict-feedback systems, command filters have also been used to limit the magnitude and rate of virtual control inputs in backstepping-based NAC designs \mbox{%DIFAUXCMD
\cite{Chen:2011aa}}\hskip0pt%DIFAUXCMD
, thereby indirectly handling input constraints}\DIFaddend .

\DIFdelbegin \DIFdel{To address the aforementioned issue, several approaches have been proposed in \mbox{%DIFAUXCMD
\cite{Gao:2006aa,Arefinia:2020aa,He:2016aa,Peng:2020aa,Esfandiari:2014aa,Karason:1994aa,Esfandiari:2015aa} }\hskip0pt%DIFAUXCMD
to handle input constraintsin NAC systems.
Particularly, the auxiliary systems have been widely introduced.
In \mbox{%DIFAUXCMD
\cite{Arefinia:2020aa,He:2016aa,Peng:2020aa}}\hskip0pt%DIFAUXCMD
, the auxiliary states were generated whenever input saturation was detected, and these states were used as feedback terms in the control law to directly compensate for the effects of input saturation constraints.
On the other hand, in \mbox{%DIFAUXCMD
\cite{Esfandiari:2014aa,Karason:1994aa,Esfandiari:2015aa}}\hskip0pt%DIFAUXCMD
}\DIFdelend \DIFaddbegin \DIFadd{Furthermore, barrier functions have been getting attention for handling system constraints.
In NAC literature, Barrier Lyapunov function (BLF)-based approaches have been studied \mbox{%DIFAUXCMD
\cite{Zhao:2025aa,Yang:2021aa}}\hskip0pt%DIFAUXCMD
.
BLF-based methods are particularly effective for enforcing state-, output-, or error-related constraints.
When input saturation is considered, they are often combined with auxiliary systems or additional compensation mechanisms.
For instance, in \mbox{%DIFAUXCMD
\cite{Yang:2021aa}}\hskip0pt%DIFAUXCMD
}\DIFaddend , the auxiliary \DIFdelbegin \DIFdel{systems were }\DIFdelend \DIFaddbegin \DIFadd{system is additionally }\DIFaddend incorporated into the \DIFdelbegin \DIFdel{adaptation law by adding the auxiliary states to the feedback error for learning.
This approach helped the NN reduce input saturationby indirectly regulating the auxiliary states, during the adaptation process}\DIFdelend \DIFaddbegin \DIFadd{BLF-based NAC design to handle input saturation}\DIFaddend .

\DIFdelbegin \DIFdel{However, these approaches typically handle input bound constraintson a per-input basis, }%DIFDELCMD < \ie %%%
\DIFdel{applying bounds to each scalar control variable individually, and may not account for more complex, nonlinear constraints, like input norm constraints }\DIFdelend \DIFaddbegin \DIFadd{Despite their effectiveness, most structural NAC methods that explicitly address input saturation are limited to simple box input constraints}\DIFaddend , \DIFdelbegin \DIFdel{which are commonly found in physical systems such as robotic actuators or motor systems. 
}\DIFdelend \DIFaddbegin \DIFadd{such as symmetric or asymmetric bounds on each input.
Thus, extending them to admissible input sets described by multiple convex constraints is not straightforward.
}\DIFaddend 

\DIFdelbegin \DIFdel{Furthermore, from a learning perspective, augmenting conventional control methods with NNs by simply adding NN outputs to existing control inputs can interfere with the adaptation process.
This occurs because the adaptation law typically depends on the feedback error, which becomes insensitive to the NN output when the conventional controller dominates the system behavior.
As a result, the NN ends up counteracting the conventional control effort rather than directly learning the system dynamics, thereby weakening its ability to adapt.
To address this, input constraints should be handled within a unified adaptive framework —such as those proposed in \mbox{%DIFAUXCMD
\cite{Esfandiari:2014aa,Karason:1994aa,Esfandiari:2015aa}}\hskip0pt%DIFAUXCMD
—where both adaptation and constraint satisfaction are treated jointly without relying on conventional control inputs.
}\DIFdelend %DIF >  ================================
%DIF >  OPTIMIZATION-BASED METHODS
%DIF >  ===============================

\DIFdelbegin %DIFDELCMD < \hfill
%DIFDELCMD < %%%
\DIFdelend \DIFaddbegin \subsubsection{\DIFadd{Optimization-Based Methods for Input Constraint Handling}}
\DIFaddend 

\DIFdelbegin \subsubsection{\DIFdel{Potential of Constrained Optimization}}
%DIFAUXCMD
\addtocounter{subsubsection}{-1}%DIFAUXCMD
%DIF <  \subsubsection{Limitations of Existing Approaches and Potential of Constrained Optimization}
\DIFdelend \DIFaddbegin \DIFadd{Besides structural methods, optimization-based methods have been studied for handling input constraints across a broad range of control design frameworks.
These methods compute the control input by solving constrained optimization problems at each time instant, as in model predictive control (MPC) \mbox{%DIFAUXCMD
\cite{Draeger:1995aa,Chen:2024aa,Cetin:2019aa,Vahidi-Moghaddam:2025aa} }\hskip0pt%DIFAUXCMD
and control barrier function (CBF)-based quadratic programming (QP) filters \mbox{%DIFAUXCMD
\cite{Ames:2019aa,Ko:2025aa,Sweatland:2025aa}}\hskip0pt%DIFAUXCMD
. 
In these methods, the constrained control problem is explicitly formulated as a constrained optimization problem.
Therefore, complex nonlinear (even non-convex) input constraints can be systematically incorporated into the control design, which is a significant advantage over structural methods.
}\DIFaddend 

%DIF <  \MSRY{각 문단에 이미 문제점을 지적해서 요약을 안해도 괜찮을 것 같습니다.}
%DIF <  \ctxt{red}
%DIF <  {
%DIF <      Although both the existing methods for weight boundedness and input constraints have shown effectiveness, they come with significant limitations. 
%DIF <      Projection operators and modification techniques do not guarantee the optimality of the adapted weights. 
%DIF <      Auxiliary systems typically handle only simple forms of input constraints, such as individual input bound constraints, limiting their ability to address more complex, nonlinear constraints.
%DIF <      Moreover, using conventional control methods in conjunction with NNs can disrupt the adaptation process.
%DIF <  }
\DIFaddbegin \DIFadd{Despite this advantage, optimization-based methods require additional conditions to ensure closed-loop stability.
Representatively, in MPC design, the recursive feasibility, computational tractability, and sufficiently long prediction horizons are essential \mbox{%DIFAUXCMD
\cite{Greune:2011aa}}\hskip0pt%DIFAUXCMD
.
Since the model accuracy is most critical across these conditions, the uncertain or unknown dynamics have motivated robust MPC \mbox{%DIFAUXCMD
\cite{Bemporad:1999aa}}\hskip0pt%DIFAUXCMD
, data-driven MPC \mbox{%DIFAUXCMD
\cite{Berberich:2025aa}}\hskip0pt%DIFAUXCMD
, and learning-based MPC \mbox{%DIFAUXCMD
\cite{Draeger:1995aa,Vahidi-Moghaddam:2025aa}}\hskip0pt%DIFAUXCMD
.
For CBF-QP filters, neuro-adaptive system identifiers have been employed to accurately impose the CBF constraint using approximated system dynamics \mbox{%DIFAUXCMD
\cite{Sweatland:2025aa}}\hskip0pt%DIFAUXCMD
.
Additionally, for computational tractability, some studies have proposed learning the optimal control policy offline using NNs \mbox{%DIFAUXCMD
\cite{Vahidi-Moghaddam:2025aa,Mallick:2025aa,Khodaverdian:2025aa}}\hskip0pt%DIFAUXCMD
.
Although these methods can provide stability guarantees under the stated conditions, the guarantees are typically obtained through feasibility, prediction, or filtering mechanisms that depend on model accuracy, uncertainty descriptions, data quality, or offline training coverage.
}\DIFaddend 

\DIFdelbegin \DIFdel{A promising approach to overcoming the limitations of the existing methods lies in constrained optimization.
By formulating the NAC problem as a unified optimization problem with constraints }\DIFdelend \DIFaddbegin \subsection{\DIFadd{Motivation and Proposed Approach}}

\DIFadd{Both structural and optimization-based methods are effective for input-constraint handling, but they have complementary limitations.
Structural methods offer Lyapunov-based stability guarantees and solver-free computation, but they are limited in the class of constraints they can accommodate.
Optimization-based methods can handle complex constraints, but their closed-loop stability guarantees are often tied to additional conditions.
}

\DIFadd{These observations motivate a unified NAC framework in which more complex input constraints are incorporated directly into the online adaptation law, as presented in \mbox{%DIFAUXCMD
\cite{Ryu:2025aa}}\hskip0pt%DIFAUXCMD
.
In the proposed CONAC framework, the NAC problem is formulated as a constrained optimization problem in which multiple convex input constraints are imposed as inequality constraints.
The constrained optimization formulation is not used to solve an online control allocation problem at each time instant; instead}\DIFaddend , it is \DIFdelbegin \DIFdel{possible to adapt the NN weights while minimizing an objective function, }%DIFDELCMD < \eg %%%
\DIFdel{tracking error, subject to both weight and input constraints, simultaneously.
Constrained optimization provides a theoretical framework for defining optimality and presents numerical methods for finding solutions that satisfy the constraints \mbox{%DIFAUXCMD
\cite{Nocedal:2006aa}}\hskip0pt%DIFAUXCMD
.
}\DIFdelend \DIFaddbegin \DIFadd{used to derive continuous-time adaptation laws for the DNN weights and the Lagrange multipliers.
Therefore, the proposed framework retains the Lyapunov-friendly structure of NAC while incorporating the constraint-handling capability of optimization-based methods.
}\DIFaddend 

\DIFdelbegin \DIFdel{In existing literature, constrained optimization techniques, such as the Augmented Lagrangian Method (ALM) \mbox{%DIFAUXCMD
\cite{Evens:2021aa} }\hskip0pt%DIFAUXCMD
and the Alternating Direction Method of Multipliers (ADMM) \mbox{%DIFAUXCMD
\cite{Wang:2019aa,Taylor:2016aa}}\hskip0pt%DIFAUXCMD
, have been used to train NNs offline. 
%DIF <  These methods impose constraints to address issues like gradient vanishing in backpropagation. 
%DIF <  \MSRY{
%DIF <      여기를 수정해야 하지 않나요? 저희의 사전 연구가 있기 때문입니다.
%DIF <  }
%DIF <  \ctxt{red}
%DIF <  {
    %DIF <  However, to the best of the authors' knowledge, no prior work has applied constrained optimization to NAC systems with real-time weight adaptation. 
    %DIF <  }
%DIF <  \ctxt{blue}
%DIF <  {
However, there are few works that have applied constrained optimization to NAC systems with real-time weight adaptation \mbox{%DIFAUXCMD
\cite{Ryu:2024aa,Ryu:2025aa}}\hskip0pt%DIFAUXCMD
.
%DIF <  }
This gap suggests that constrained optimization could be key to addressing both weight boundedness and input constraints in a unified, theoretically grounded framework, particularly in real-time NAC.
}%DIFDELCMD < 

%DIFDELCMD < %%%
\DIFdelend \subsection{Contributions}

The main contributions of this study are\DIFdelbegin \DIFdel{listed as follows}\DIFdelend :
\begin{itemize}
    \item 
    A constrained optimization-based neuro-adaptive control (CONAC) \DIFdelbegin \DIFdel{is proposed, where the control problem is formulated as a constrained optimization problem. In this formulation, weight and }\DIFdelend \DIFaddbegin \DIFadd{framework is proposed for unknown systems under multiple }\DIFaddend convex input constraints\DIFdelbegin \DIFdel{are incorporated as inequality constraints, while minimizing a given objective function.
    }\DIFdelend \DIFaddbegin \DIFadd{;
    %DIF >  The admissible input set is represented by convex inequality constraints and embedded directly into the online adaptation law;
}

    \DIFaddend \item
    \DIFdelbegin \DIFdel{Convex input constraints can be applied to the proposed CONAC, including input bound constraints and norm constraints. Moreover, the proposed CONAC can handle any combination of these constraints , which can be useful in practical applications where multiple constraints need to be satisfied simultaneously.
    }\DIFdelend \DIFaddbegin \DIFadd{Any combination of convex input constraints can be handled, thereby significantly extending the applicability to real-world systems with complex input constraints;
    }

    \DIFaddend \item 
    \DIFdelbegin \DIFdel{Since the proposed CONAC approximates the entire ideal control law without any conventional controller, prior knowledge of the system dynamics or the system identification process for designing a nominal conventional controller can be avoided, which is often difficult and time-consuming to implement in practice. 
    }%DIFDELCMD < \item %%%
\item%DIFAUXCMD
\DIFdelend The adaptation laws for \DIFaddbegin \DIFadd{adaptive variables including the }\DIFaddend DNN weights and \DIFaddbegin \DIFadd{the }\DIFaddend Lagrange multipliers are systematically derived \DIFdelbegin \DIFdel{to solve the defined problem, using constrained optimization theory. 
    These laws guarantee convergence to the }\DIFdelend \DIFaddbegin \DIFadd{from the formulated constrained optimization problem. 
    At an equilibrium of the adaptation dynamics, the KKT-like }\DIFaddend first-order optimality conditions \DIFdelbegin \DIFdel{at steady state, specifically the Karush–Kuhn–Tucker (KKT) conditions, as described in }\DIFdelend \DIFaddbegin \DIFadd{are satisfied }\DIFaddend \cite[Chap. 12, Thm. 12.1]{Nocedal:2006aa}\DIFdelbegin \DIFdel{.
    }\DIFdelend \DIFaddbegin \DIFadd{; and
}

    \DIFaddend \item 
    \DIFdelbegin \DIFdel{The proposed CONAC's stability is rigorously analyzed using Lyapunov theory, ensuring that }\DIFdelend \DIFaddbegin \DIFadd{Lyapunov-based stability analysis is provided to establish boundedness of }\DIFaddend the tracking error\DIFdelbegin \DIFdel{and DNN weights remain bounded. This analysis guarantees that the proposed CONAC can achieve stable online adaptation.
    }%DIFDELCMD < \item %%%
\item%DIFAUXCMD
\DIFdel{The proposed CONAC is implemented in }\DIFdelend \DIFaddbegin \DIFadd{, and the adaptive variables.
    The effectiveness and feasibility are validated through simulations and }\DIFaddend real-time \DIFaddbegin \DIFadd{experiments }\DIFaddend on a 2-DOF robotic manipulator\DIFdelbegin \DIFdel{, demonstrating its effectiveness in achieving accurate trajectory tracking while satisfying all imposed constraints.
The experimental results validate the theoretical findings and confirm the practical applicability of the proposed CONAC.
}\DIFdelend \DIFaddbegin \DIFadd{.
}

\DIFaddend \end{itemize}

\DIFdelbegin %DIFDELCMD < \hfill
%DIFDELCMD < 

%DIFDELCMD < %%%
\DIFdelend This study extends our preliminary work in \DIFdelbegin \DIFdel{\mbox{%DIFAUXCMD
\cite{Ryu:2025aa}}\hskip0pt%DIFAUXCMD
}\DIFdelend \DIFaddbegin \DIFadd{\mbox{%DIFAUXCMD
\cite{Ryu:2025ab,Ryu:2025aa}}\hskip0pt%DIFAUXCMD
}\DIFaddend , which introduced inequality constraints on the weight and input norms within a constrained optimization framework for \DIFdelbegin \DIFdel{neuro-adaptive control (NAC ) systems using a SHLNN}\DIFdelend \DIFaddbegin \DIFadd{NAC using single hidden layer NNs}\DIFaddend . 
In the present work, the same framework is generalized to accommodate \DIFdelbegin \DIFdel{arbitrary }\DIFdelend \DIFaddbegin \DIFadd{multiple }\DIFaddend convex input constraints and extended to DNNs, while still \DIFdelbegin \DIFdel{retaining the capability to regulate the weight norm}\DIFdelend \DIFaddbegin \DIFadd{ensuring Lyapunov-based stability guarantees}\DIFaddend .
Furthermore, unlike the previous study, which was limited to numerical simulations, the proposed CONAC is experimentally validated through real-time implementation on a 2-DOF robotic manipulator.

\subsection{Organization}

The remainder of this paper is organized as follows. 
Section\DIFdelbegin \DIFdel{\ref{sec:Problem Formulation} }\DIFdelend \DIFaddbegin \DIFadd{~\ref{sec:problem:formulation} }\DIFaddend presents the target system and control objective.
Section\DIFdelbegin \DIFdel{\ref{sec:ctrl design} }\DIFdelend \DIFaddbegin \DIFadd{~\ref{sec:main:result} }\DIFaddend introduces the proposed CONAC and the \DIFdelbegin \DIFdel{architecture of the DNN used in the controller. 
In Section \ref{sec:adap_laws}, the adaptation law is developed. 
Section\ref{sec:stability} analyzes the stability of CONAC.
Section \ref{sec:sim} presents a }\DIFdelend \DIFaddbegin \DIFadd{adaptation laws. 
Section~\ref{sec:stability:analysis} analyzes the closed-loop stability.
%DIF >  Section~\ref{sec:tuning} provides guidelines for tuning the parameters of CONAC, including the adaptation gain, update rates of Lagrange multipliers, and initialization of NN weights.
Section~\ref{sec:val} presents a numerical simulation comparing the proposed CONAC with existing methods, and a }\DIFaddend real-time implementation of CONAC on a 2-DOF robotic manipulator.
Finally, Section\DIFaddbegin \DIFadd{~}\DIFaddend \ref{sec:conclusion} concludes the paper and discusses potential future work.
\DIFdelbegin \DIFdel{Candidates for convex input constraints are provided in Appendix \ref{sec:appen:cstr}. 
}\DIFdelend 

%DIF <  ====================================
%DIF >  ------------------------------------
%       SECTION PROBLEM FORMULATION
%DIF <  ====================================
%DIF >  ------------------------------------
\section{Problem Formulation} \DIFdelbegin %DIFDELCMD < \label{sec:Problem Formulation}
%DIFDELCMD < %%%
\DIFdelend \DIFaddbegin \label{sec:problem:formulation}
\DIFaddend 

\DIFdelbegin \subsection{\DIFdel{Model Dynamics and Control Objective}}
%DIFAUXCMD
\addtocounter{subsection}{-1}%DIFAUXCMD
%DIFDELCMD < 

%DIFDELCMD < %%%
\DIFdelend We consider an unknown \DIFdelbegin \DIFdel{Euler-Lagrange }\DIFdelend \DIFaddbegin \DIFadd{nonlinear }\DIFaddend system modeled as 
\DIFdelbegin %DIFDELCMD < \begin{equation}
%DIFDELCMD <     \rbM\ddq +\rbVm \dq+ \rbF + \rbG + \rbd
%DIFDELCMD <     =
%DIFDELCMD <     \mysat(\rbu)
%DIFDELCMD <     ,
%DIFDELCMD <     \label{eq:sys1}
%DIFDELCMD < \end{equation}%%%
\DIFdelend \DIFaddbegin \begin{equation} \label{eq:sys}
    \ddtt\mv{x} = \mv{f}(\mv{x}) + \mm{B}(\mv{x})\mysat(\mv{u})
    ,
\end{equation}\DIFaddend 
where \DIFdelbegin \DIFdel{$\q\in \R^n$ denotes the generalized coordinate, $\rbd$ represents the disturbance and $\rbu\in\R^n$ denotes the generalized control input . 
The terms $\rbM:=\rbM(\q)\in\R^{n\times n}$, $\rbVm:=\rbVm(\q,\dq)\in\R^{n\times n}$, $\rbF:=\rbF(\dq)\in\R^{n}$, and $\rbG:=\rbG(\q)\in\R^{n}$ represent the unknown inertia matrix, Coriolis/centripetal matrix, friction vector, and gravity vector, respectively.
The function $\mysat(\cdot):\R^n\to\R^n$ represents the inherent physical limitations of the actuators such that $\norm{\mysat(\rbu)}\le \overline\tau$, }\DIFdelend \DIFaddbegin \DIFadd{$\mv{x},\mv{u}\in\R^n$ denote state and control input vectors, respectively, and $\mv{f}:\R^n\to\R^n$ and $\mm{B}:\R^n\to\R^{n\times n}$ denote unknown system function and control effectiveness matrix, which are locally Lipschitz continuous, respectively, and $\mysat:\R^n\to\R^n$ denotes a saturation function.
%DIF >  The compact set $\Omega_x$ denotes a approximation region where NN will approximate with minimum bounded approximation error, later.
%DIF >  The saturation function $\mysat:\R^n\to\Omega_u$ radially clips the input onto the admissible input set $\Omega_u$, as illustrated in Fig.~\ref{fig:saturation:prop:multi}.
}

\DIFadd{The system \eqref{eq:sys} is subject to the following assumptions.
}

\begin{assum} \label{assum:B:bound}
    \DIFadd{The control effectiveness matrix $\mm{B}(\mv{x})$ is nonsingular and bounded as   
    $
        \underline{b}\mm{I}_n
        \preceq 
        \mm{B}(\mv{x})
        \preceq 
        \bar{b}\mm{I}_n
        ,
    $
    }\DIFaddend where \DIFdelbegin \DIFdel{$\overline\tau\in\R_{>0}$ denotes the maximum norm of control input }\DIFdelend \DIFaddbegin \DIFadd{$\underline{b},\bar{b}\in\R_{>0}$ are known positive constants.\
}\end{assum}

\begin{assum} \label{assum:B:bound:grad}
    \DIFadd{\mbox{%DIFAUXCMD
\cite[Assumption 1(b)]{Boulkroune:2010aa}
    }\hskip0pt%DIFAUXCMD
There exists an unknown positive function $\beta(\mv{x}):\R^n\to\R_{>0}$ such that $\tfrac{1}{2}\mv{v}^\top\left(\ddtt\mm{B}(\mv{x})^{-1}\right)\mv{v} \le \beta(\mv{x})\norm{\mv{v}}^2$, $\forall \mv{v}\in\R^n$ and $\forall \mv{x}\in\R^n$}\DIFaddend .
\DIFaddbegin \end{assum}
\DIFaddend 

\DIFdelbegin \DIFdel{In this study, the }\DIFdelend \DIFaddbegin \begin{assum} \label{assum:sat}
    \DIFadd{The admissible input set $\Omega_u$ is defined by smooth convex constraint functions $c_j^u(\mv u)$, $\forall j\in\mathcal I$, where $\mathcal{I}$ denotes the input constraint index set, as
    $
        \Omega_u
        =
        \bigcap_{j\in\mathcal I}
        \{\mv u\in\R^n:c_j^u(\mv u)\le0\}
        .
    $
    The origin is contained in its interior, }\ie \DIFadd{$-\underline{c}_j^u \le c_j^u(\mv 0)<0$, $\forall j\in\mathcal I$, for some known $\underline{c}_j^u\in\R_{>0}$.
  The }\DIFaddend saturation function $\mysat(\cdot)$ \DIFdelbegin \DIFdel{is assumed to be convex, which is practically common in many engineering systems.This is because actuator limits are typically convex —for example, minimum and maximum torque constraints —which define box- or ball-shaped bounds in the input space.
To account for these limitations, it is essential to incorporate physically motivated constraints into the controller design.
Appendix \ref{sec:appen:cstr} introduces candidate constraints that can be applied to ensure compliance with these physical limitations.
}\DIFdelend \DIFaddbegin \DIFadd{satisfies $\mysat(\mv u)=\mv u$ for $\mv u\in\Omega_u$ and, for $\mv u\notin\Omega_u$, $\mysat(\mv u)$ lies on the ray from the origin through $\mv u$, as illustrated in Fig.~\ref{fig:saturation}.
    Moreover, there exists $\underline u>0$ such that
    $
        \inf_{\mv{u}\notin\Omega_u}
        \left(
            \norm{\mysat(\mv u)}
        \right)
        =
        \underline u
        .
    $
}\end{assum}
\DIFaddend 

\DIFdelbegin \DIFdel{The Euler-Lagrange system \eqref{eq:sys1} satisfies several important physical properties, as presented in \mbox{%DIFAUXCMD
\cite[Chap. 3, Tab. 3.2.1]{Lewis:1998aa}}\hskip0pt%DIFAUXCMD
.The key properties are introduced below:
}%DIFDELCMD < \begin{prop} 
%DIFDELCMD <     %%%
\DIFdel{The inertia matrix $\rbM$ is symmetric, positive definite, and bounded.
}%DIFDELCMD < \label{prop:M}
%DIFDELCMD < \end{prop}
%DIFDELCMD < %%%
\DIFdelend %DIF >  \begin{assum} \label{assum:sat}
%DIF >      The admissible input set $\Omega_u$ is convex and defined by multiple convex functions $c_j^{u}(\mv{u}),\ \forall j\in\mathcal{I},$ where each $c_j^{u}(\mv{u})$ is $\mathcal{C}^1$ and $\mathcal{I}$ is an input constraint index set.
%DIF >      Thus, the admissible input set is defined as
%DIF >      $
%DIF >          \Omega_u 
%DIF >          = 
%DIF >          \bigcap_{j \in \mathcal{I}} \left\{ \mv{u} \in \mathbb{R}^n : c_j^{u}(\mv{u}) \leq 0 \right\}
%DIF >          .
%DIF >      $ 
%DIF >      In addition, the admissible input set contains the origin in its interior, \ie $c_j^{u}(\mv{0})<0,\ \forall j\in\mathcal{I}$.
%DIF >      Moreover, the norm of any boundary point of $\Omega_u$ is bounded, such that $0<\underline{u}\le\norm{\mv{u}}\le\overline{u},\ \forall\mv{u}\in\partial\Omega_u$, where $\partial \Omega_u$ denotes the boundary of $\Omega_u$.
%DIF >  \end{assum}

\DIFdelbegin %DIFDELCMD < \begin{prop} 
%DIFDELCMD <     %%%
\DIFdel{The Coriolis/centripetal matrix $\rbVm$ can always be selected, so that the matrix $\ddtt{\rbM}-2\rbVm$ is skew-symmetric, }%DIFDELCMD < \ie %%%
\DIFdel{$\mv{x}^\top(\ddtt{\rbM}-2\rbVm)\mv{x}=0,\ \forall \mv{x}\in\R^n$.
}%DIFDELCMD < \label{prop:skew}
%DIFDELCMD < \end{prop}
%DIFDELCMD < %%%
\DIFdelend %DIF >  \begin{assum} \label{assum:sat}
%DIF >      The admissible input set $\Omega_u$ is convex and defined by multiple convex functions $c_j^{u}(\mv{u}),\ \forall j\in\mathcal{I},$ where each $c_j^{u}(\mv{u})$ is $\mathcal{C}^1$ and $\mathcal{I}$ is an input constraint index set.
%DIF >      Thus, the admissible input set is defined as
%DIF >      $
%DIF >          \Omega_u 
%DIF >          = 
%DIF >          \bigcap_{j \in \mathcal{I}} \left\{ \mv{u} \in \mathbb{R}^n : c_j^{u}(\mv{u}) \leq 0 \right\}
%DIF >          .
%DIF >      $ 
%DIF >      In addition, the admissible input set contains the origin in its interior, \ie $c_j^{u}(\mv{0})<0,\ \forall j\in\mathcal{I}$.
%DIF >      Moreover, the norm of any boundary point of $\Omega_u$ is bounded, such that $0<\underline{u}\le\norm{\mv{u}}\le\overline{u},\ \forall\mv{u}\in\partial\Omega_u$, where $\partial \Omega_u$ denotes the boundary of $\Omega_u$.
%DIF >  \end{assum}

\DIFdelbegin %DIFDELCMD < \begin{prop}
%DIFDELCMD <     %%%
\DIFdel{The disturbance $\rbd$ is bounded, }%DIFDELCMD < \ie %%%
\DIFdel{$\norm{\rbd}\le \overline\tau_d$, where $\overline\tau_d\in\R_{\ge0}$.
    }%DIFDELCMD < \label{prop:dis_bound}
%DIFDELCMD < \end{prop}
%DIFDELCMD < %%%
\DIFdelend \DIFaddbegin \DIFadd{By Assumption~\ref{assum:sat}, the saturation function $\mysat(\mv{u})$ can represent any combination of convex input constraints, rather than simple box constraints which are typically used in the NAC literature, as illustrated in Fig.~\ref{fig:saturation}.
}\DIFaddend 

\DIFaddbegin \begin{figure}[t]
    \centering
    %DIF >  \subfloat[]{
    \subfloat[Conventional input constraints.]{
        \includegraphics[width=0.45\linewidth]{\src/figures/inputFuncConv.drawio.pdf}
        \label{fig:saturation:conv:box}
    }
    %DIF >  \hfill
    %DIF >  \subfloat[]{
    \subfloat[Multiple convex input constraints.]{
        \includegraphics[width=0.45\linewidth]{\src/figures/inputFuncProp.drawio.pdf}
        \label{fig:saturation:prop:multi}
    }
    \caption{
        \DIFaddFL{Comparison of conventional box input constraints and multiple convex input constraints.
        In Fig.~\ref{fig:saturation:conv:box}, symmetric and asymmetric box constraints are illustrated.
        In Fig.~\ref{fig:saturation:prop:multi}, a combination of symmetric and norm constraint are illustrated.
    }}
    \label{fig:saturation}
\end{figure}


%DIF >  According to locally Lipschitz continuity of $\mv{f}(\cdot)$ and $\mm{B}(\cdot)$, and the boundedness of the saturation function $\mysat(\cdot)$ as defined in Assumption \ref{assum:sat}, the system \eqref{eq:sys} is locally Lipschitz continuous.

\DIFadd{Additionally, the desired trajectory $\mv{x}_d(t):\R_{\ge0}\to\R^n$ is assumed to be given and satisfies the following assumption.
}

\begin{assum} \label{assum:xd}
    \DIFadd{The desired trajectory $\mv{x}_d(t)$ is sufficiently smooth and bounded, such that $\mv{x}_d,\ddtt\mv{x}_d$ are bounded.
    %DIF >  The desired trajectory $\mv{x}_d(t)$ is sufficiently smooth and bounded, such that $\mv{x}_d\in\Omega_{x_d}:=\left\{\mv{x}_d\in\R^n : \norm{\mv{x}_d}\le \overline{x}_d\right\}$ and $\ddtt\mv{x}_d$ is bounded in a compact set $\Omega_{dx_d}$, as well.
}\end{assum}

\DIFaddend The control \DIFdelbegin \DIFdel{objective is to develop a NAC that enables $\q$ to track a continuously differentiable desired trajectory $\qd:=\qd(t): \R_{\ge0} \to \R^n$, compensating }\DIFdelend \DIFaddbegin \DIFadd{objectives are to design a NAC }\DIFaddend for the unknown system \DIFdelbegin \DIFdel{dynamics while addressing the imposed constraints .
The desired trajectory $\qd(t)$ is supposed to satisfy the following assumption:
}%DIFDELCMD < \begin{assum}
%DIFDELCMD <     %%%
\DIFdel{The desired trajectory $\qd(t)$ is assumed to be bounded, }%DIFDELCMD < \ie %%%
\DIFdel{$\norm{\qd(t)}\le \overline{q}_d\in\R_{>0}$, and available for designing a bounded control input.
}%DIFDELCMD < \label{assum:feasible}
%DIFDELCMD < \end{assum}
%DIFDELCMD < %%%
\DIFdelend \DIFaddbegin \DIFadd{\eqref{eq:sys} such that $\mv{x}(t)$ tracks $\mv{x}_d(t)$, while satisfying the multiple convex input constraints defined by $\mysat(\cdot)$, and ensuring closed-loop stability.
}\DIFaddend 

%DIF <   SECTION CONTROLLER DESIGN ==============================
%DIF >  ------------------------------------
%DIF >        SECTION MAIN RESULT
%DIF >  ------------------------------------
\section{\DIFdelbegin \DIFdel{Control Law Development}\DIFdelend \DIFaddbegin \DIFadd{Proposed CONAC Framework}\DIFaddend } \DIFdelbegin %DIFDELCMD < \label{sec:ctrl design}
%DIFDELCMD < %%%
\DIFdelend \DIFaddbegin \label{sec:main:result}
\DIFaddend 

\DIFaddbegin \begin{figure}[!t]
    \centering
    \includegraphics[width=.79\linewidth]{\src/figures/Controller.drawio.pdf}
    \caption{
        \cred{(draft figure)}
        \DIFaddFL{Architecture of the proposed constrained optimization-based neuro-adaptive controller (CONAC).
    }}
    \label{fig:ctrl:diagram}
\end{figure}

\DIFaddend The architecture of the proposed CONAC consists of a DNN that functions as a NAC and a \DIFdelbegin \DIFdel{weight optimizer for the DNN}\DIFdelend \DIFaddbegin \DIFadd{adaptation process for the adaptive variables}\DIFaddend , as illustrated in Fig.\DIFaddbegin \DIFadd{~}\DIFaddend \ref{fig:ctrl:diagram}.
Before \DIFdelbegin \DIFdel{proceeding to the controller design, some mathematical preliminaries are revisited in Section\ref{sec:sub:math preliminaries}.
Section \ref{sec:sub:NAC} introduces the NAC, followed by the DNN model in Section\ref{sec:sub:NN definition}. 
The constrained optimization-based weight adaptation law is presented in Section\ref{sec:adap_laws}}\DIFdelend \DIFaddbegin \DIFadd{presenting the CONAC, we first review some preliminaries in Section~\ref{sec:main:result:sub:prelim}.
Then, the CONAC is presented in Section~\ref{sec:main:result:sub:ctrl:design} along with the adaptation laws in Section~\ref{sec:main:result:sub:adap:laws}}\DIFaddend .

\DIFdelbegin %DIFDELCMD < \begin{figure}[!t]
%DIFDELCMD <     \centering
%DIFDELCMD <     \includegraphics[width=.99\linewidth]{\src/figures/Controller.drawio.pdf}
%DIFDELCMD <     %%%
%DIFDELCMD < \caption{%
{%DIFAUXCMD
\DIFdelFL{Architecture of the proposed constrained optimization-based neuro-adaptive controller (CONAC).}}
    %DIFAUXCMD
%DIFDELCMD < \label{fig:ctrl:diagram}
%DIFDELCMD < \end{figure}
%DIFDELCMD < %%%
\DIFdelend \DIFaddbegin \subsection{\DIFadd{Preliminaries}} \label{sec:main:result:sub:prelim}
\DIFaddend 

\DIFdelbegin \subsection{\DIFdel{Mathematical Preliminaries}}%DIFAUXCMD
\addtocounter{subsection}{-1}%DIFAUXCMD
%DIFDELCMD < \label{sec:sub:math preliminaries}
%DIFDELCMD < 

%DIFDELCMD < %%%
\DIFdelend The following proposition \DIFdelbegin \DIFdel{and lemma }\DIFdelend will be used in \DIFaddbegin \DIFadd{the }\DIFaddend subsequent sections:
%DIF >  The following propositions and lemmas will be used in the subsequent sections:

\DIFdelbegin %DIFDELCMD < \begin{propsit}[see {\cite[Chap. 7, Prop. 7.1.9]{Bernstein:2009aa}}] %%%
\DIFdelend \DIFaddbegin \begin{propsit}[see {\cite[Prop. 7.1.9]{Bernstein:2009aa}}] \DIFaddend \label{propsit:kron}
    For a matrix $\mm A\in\R^{n\times m}$ and a vector $\mv b\in\mathbb R^n$, the following property holds:
    \begin{equation}
        \mm A^\top \mv b 
        = 
        \myvec(\mm A^\top\mv b)
        =
        \myvec(\mv b^\top\mm A)
        = 
        (\mm{I}_{m}\otimes \mv b^\top) \myvec(\mm A)
        .
    \end{equation}
\end{propsit}

\DIFdelbegin %DIFDELCMD < \begin{lem} \label{lem:stable:set}
%DIFDELCMD <     %%%
\DIFdel{For a Lyapunov function $V:=V(\mv{x})\ge0$ with $\mv{x}\in\R^n$, if the time derivative of $V$ is given by $\ddtt V\le -a_1 \mv{x}^\top \mv{x} + a_2 \mv{x}$ for some positive constants $a_1\in\R_{>0}$ and $a_2\in\R_{>0}$, then $\mv{x}$ is bounded as $\mv{x}\in\{\mv{x}\mid\norm{\mv{x}}\le\tfrac{a_2}{a_1}\}$.
}%DIFDELCMD < \end{lem}
%DIFDELCMD < %%%
\DIFdelend %DIF >  \begin{lem}[Comparison Lemma {\cite[pp. 102--103, Lemma 3.4]{Khalil:2002aa}}] \label{lem:comp}
%DIF >      Consider a scalar continuous function $V(t)\ge0,\ \forall t\in\R_{\ge0}$.
%DIF >      If $V(t)$ satisfies the differential inequality    
%DIF >      % \begin{equation}
%DIF >      $
%DIF >          \ddtt V(t) \le -a V(t) + b
%DIF >          ,
%DIF >      % \end{equation}
%DIF >      $
%DIF >      with $a,b\in\R_{>0}$, then $V(t)$ is bounded as 
%DIF >      % \begin{equation}
%DIF >      $
%DIF >          V(t) \le \left(V(0) - \tfrac{b}{a}\right)\exp^{-at} + \tfrac{b}{a}
%DIF >          ,
%DIF >          \ \forall t\in\R_{\ge0}
%DIF >          .
%DIF >      % \end{equation}
%DIF >      $
%DIF >  \end{lem}

\DIFdelbegin %DIFDELCMD < \begin{proof}
%DIFDELCMD <     %%%
\DIFdel{The time derivative of $V$ can be rewritten as
    }%DIFDELCMD < \begin{equation}
%DIFDELCMD <         \ddtt V
%DIFDELCMD <         \le
%DIFDELCMD <         -a_1 \norm{\mv{x}}^2 + a_2 \norm{\mv{x}}
%DIFDELCMD <         =
%DIFDELCMD <         -a_1
%DIFDELCMD <         \left(
%DIFDELCMD <             \norm{\mv{x}}-\tfrac{a_2}{2a_1}
%DIFDELCMD <         \right)^2
%DIFDELCMD <         +
%DIFDELCMD <         \tfrac{a_2^2}{4a_1}
%DIFDELCMD <         .
%DIFDELCMD <         \label{eq:lem:lya:dot}
%DIFDELCMD <     \end{equation}
%DIFDELCMD <     %%%
\DIFdel{According to \eqref{eq:lem:lya:dot}, it can be concluded that $\ddtt V$ is negative definite if $\norm{\mv{x}}>\tfrac{a_2}{a_1}$ holds.
    In other words, $\mv{x}$ remains within the bounded set $\{\mv{x}\mid\norm{\mv{x}}\le\tfrac{a_2}{a_1}\}$, since $\ddtt V$ is negative definite when $\mv{x}$ is outside the set.
}%DIFDELCMD < \end{proof}
%DIFDELCMD < %%%
\DIFdelend %DIF >  \begin{proof}
%DIF >      See \cite[pp. 689--660]{Khalil:2002aa}.
%DIF >  \end{proof}

\DIFdelbegin \subsection{\DIFdel{Neuro-Adaptive Controller Design}}%DIFAUXCMD
\addtocounter{subsection}{-1}%DIFAUXCMD
%DIFDELCMD < \label{sec:sub:NAC}
%DIFDELCMD < %%%
\DIFdelend %DIF >  Based on the Comparison Lemma (Lemma \ref{lem:comp}), the following lemma establishes the uniform ultimate boundedness (UUB) of the state vector $\mv{z}(t):\R_{\ge0}\to\R^n$ within the constrained set $\mathcal{Z}:=\left\{\mv{z}\in\R^n : \norm{\mv{z}}\le\overline{z}\in\R_{>0}\right\}$.

\DIFdelbegin \DIFdel{Given that the Euler-Lagrange system \eqref{eq:sys1} exhibits second-order dynamics, a filtered error $\fe\in\R^n$ is introduced to convert the system into a first-order system, as follows:
}%DIFDELCMD < \begin{equation}
%DIFDELCMD <     \fe := \ddtt{\mv e}+\mm\Lambda \mv e
%DIFDELCMD <     ,
%DIFDELCMD <     \label{eq:err:filtered}
%DIFDELCMD < \end{equation}
%DIFDELCMD < %%%
\DIFdel{where $\mv e := \q - \qd$ denotes the tracking error, $\ddtt{\mv e} := \dq - \dqd$ represents the time derivative of the tracking error, and $\mm\Lambda\in\R^{n\times n}_{>0}$ is a user-designed filtering matrix.
Since \eqref{eq:err:filtered} is stable system, it implies that $\mv e$ is bounded if $\fe$ is bounded.
}\DIFdelend %DIF >  \begin{lem} \label{lem:comp:compact}
%DIF >      Consider a scalar continuous function $V(\mv{z})$, such that $\gamma_1\norm{\mv{z}}^2 \le V(\mv{z}) \le \gamma_2\norm{\mv{z}}^2$, where $\gamma_1,\gamma_2\in\R_{>0}$.
%DIF >      If $V(\mv{z})$ satisfies    
%DIF >      \begin{equation}
%DIF >          \ddtt V(\mv{z}) 
%DIF >          \le 
%DIF >          -\gamma_3\norm{\mv{z}}^2 + c
%DIF >          ,\ 
%DIF >          \forall t\in\R_{\ge0}
%DIF >          ,\ 
%DIF >          \forall \mv{z}(t)\in\mathcal{Z}
%DIF >          ,
%DIF >      \end{equation}
%DIF >      where $\gamma_3,c\in\R_{>0}$, then $\mv{z}(t),\ \forall \mv{z}\in\mathcal{Z}$, is bounded as 
%DIF >      \begin{equation}
%DIF >          \norm{\mv{z}(t)} 
%DIF >          \le 
%DIF >          \sqrt{
%DIF >              \left(
%DIF >                  \tfrac{\gamma_2}{\gamma_1}    
%DIF >                  \norm{\mv{z}(0)}^2
%DIF >                  - 
%DIF >                  \tfrac{\gamma_2 c}{\gamma_1\gamma_3}
%DIF >              \right)
%DIF >              \exp(-\tfrac{\gamma_3}{\gamma_2}t) 
%DIF >              + 
%DIF >              \tfrac{\gamma_2 c}{\gamma_1\gamma_3}
%DIF >          }
%DIF >          .
%DIF >      \end{equation}
%DIF >      Conversely, to ensure that $\mv{z}(t)\in\mathcal{Z},\ \forall t\in\R_{\ge0}$, it is sufficient to satisfy
%DIF >      % $
%DIF >      %     \sqrt{
%DIF >      %         \tfrac{\gamma_2}{\gamma_1}    
%DIF >      %         \norm{\mv{z}(0)}^2 
%DIF >      %         +
%DIF >      %         \tfrac{\gamma_2 c}{\gamma_1\gamma_3}
%DIF >      %     }
%DIF >      %     \le
%DIF >      %     \overline{z}
%DIF >      %     ,
%DIF >      % $
%DIF >      $
%DIF >          \max
%DIF >          \left(
%DIF >              \sqrt{
%DIF >                  \tfrac{\gamma_2}{\gamma_1}
%DIF >              }
%DIF >              \norm{\mv{z}(0)}
%DIF >              ,
%DIF >              \sqrt{
%DIF >                  \tfrac{\gamma_2 c}
%DIF >                  {\gamma_1\gamma_3}
%DIF >              }
%DIF >          \right)
%DIF >          \le
%DIF >          \overline{z}
%DIF >      $
%DIF >      which holds if 
%DIF >      % $
%DIF >      %     \norm{\mv{z}(0)} 
%DIF >      %     \le 
%DIF >      %     \sqrt{
%DIF >      %         \tfrac{\gamma_1}{\gamma_2}\overline{z}^2
%DIF >      %         -
%DIF >      %         \tfrac{c}{\gamma_3}
%DIF >      %     }
%DIF >      % $
%DIF >      $
%DIF >          \norm{\mv{z}(0)}
%DIF >          \le
%DIF >          \overline{z}
%DIF >          \sqrt{
%DIF >              \tfrac{\gamma_1}{\gamma_2}
%DIF >          }
%DIF >      $
%DIF >      .
%DIF >  \end{lem}

\DIFdelbegin \DIFdel{Using $\fe$, the system dynamics \eqref{eq:sys1} can be rewritten as
}%DIFDELCMD < \begin{equation}
%DIFDELCMD <     \rbM \ddtt\fe
%DIFDELCMD <     =
%DIFDELCMD <     -\rbVm \fe
%DIFDELCMD <     -\mm K \fe
%DIFDELCMD <     + \mv f
%DIFDELCMD <     -\rbd + \mysat(\rbu)
%DIFDELCMD <     ,
%DIFDELCMD <     \label{eq:sys:filtered}
%DIFDELCMD < \end{equation}
%DIFDELCMD < %%%
\DIFdel{where $\mm{K}=\mm{K}^\top\in\R^{n\times n}_{>0}$ denotes an arbitrary unknown matrix and $
    \mv f:= \mv f
    \left(
        \q,\dq,\qd,\dqd,\ddqd
    \right)
    =
    \mm K \fe
    +\rbM
    \left(
        -\ddqd+\mm\Lambda\ddtt{\mv e}
    \right)
    +
    \rbVm
    \left(
        -\dqd+\mm\Lambda\mv e
    \right)
    -
    \rbF
    -
    \rbG
    \in\R^n
$ denotes the lumped unknown system.
%DIF <  \color{red}
%DIF <  The unknown matrix $\mm{K}$ is not available to design in practice, 
%DIF <  \color{black}
}\DIFdelend %DIF >  \begin{proof}
%DIF >      % Using Comparison Lemma \cite[pp. 689--660, Lemma 3.4]{Khalil:2002aa}, the proof is straightforward.
%DIF >      Using Comparison Lemma (Lemma \ref{lem:comp}), the proof is straightforward, and thus omitted for brevity.
%DIF >  \end{proof}

\DIFdelbegin \DIFdel{Consider the Lyapunov function $V_1:= \tfrac{1}{2}\fe^\top \rbM\fe$. 
Invoking Property \ref{prop:skew}, the time derivative of $V_1$ is
}%DIFDELCMD < \begin{equation}
%DIFDELCMD <     \begin{aligned}
%DIFDELCMD <         \ddtt{V}_1
%DIFDELCMD <         =&
%DIFDELCMD <         \fe^\top \rbM \ddtt \fe
%DIFDELCMD <         +
%DIFDELCMD <         \tfrac{1}{2}\fe^\top \ddtt{\rbM}\fe
%DIFDELCMD <         \\
%DIFDELCMD <         =&
%DIFDELCMD <         \fe^\top 
%DIFDELCMD <         \left(
%DIFDELCMD <             -\rbVm \fe -\mm K \fe + \mv f
%DIFDELCMD <             -\rbd + \mysat(\rbu)
%DIFDELCMD <         \right)
%DIFDELCMD <         \\
%DIFDELCMD <         &
%DIFDELCMD <         +
%DIFDELCMD <         \tfrac{1}{2}
%DIFDELCMD <         \fe^\top \ddtt{\rbM}\fe
%DIFDELCMD <         \\
%DIFDELCMD <         =&
%DIFDELCMD <         -
%DIFDELCMD <         \fe^\top \mm K \fe 
%DIFDELCMD <         +
%DIFDELCMD <         \fe^\top 
%DIFDELCMD <         \left(
%DIFDELCMD <             \mv f
%DIFDELCMD <             -\rbd
%DIFDELCMD <             +\mysat(\rbu)
%DIFDELCMD <         \right)
%DIFDELCMD <         \\
%DIFDELCMD <         &
%DIFDELCMD <         +
%DIFDELCMD <         \tfrac{1}{2}
%DIFDELCMD <         \fe^\top(\ddtt{\rbM}- 2\rbVm)\fe
%DIFDELCMD <         \\
%DIFDELCMD <         \le&
%DIFDELCMD <         -\lambda_{\min}(\mm K)\norm{\fe}^2
%DIFDELCMD <         +
%DIFDELCMD <         \overline\tau_d\norm{\fe}
%DIFDELCMD <         +
%DIFDELCMD <         \fe^\top (\mv f+\mysat(\rbu))
%DIFDELCMD <         % \\
%DIFDELCMD <         % \le&
%DIFDELCMD <         % -\lambda_{\min}(\mm K)\norm{\fe}^2
%DIFDELCMD <         % +
%DIFDELCMD <         % \overline\tau_d\norm{\fe}
%DIFDELCMD <         % +
%DIFDELCMD <         % \fe^\top (\mysat(\rbu)-\rbu^*)
%DIFDELCMD <         .
%DIFDELCMD <     \end{aligned}
%DIFDELCMD <     \label{eq:lya1:dot1}
%DIFDELCMD < \end{equation} 
%DIFDELCMD < %%%
\DIFdel{Therefore, invoking Lemma \ref{lem:stable:set}, it can be concluded that the filtered error $\fe$ is uniformly ultimately bounded with $\lim_{t\to\infty}\norm{\fe}\le\tfrac{\overline\tau_d}{\lambda_{\min}(\mm K)}$, provided that the lumped unknown system $\mv{f}$ can be perfectly cancelled by the control input $\rbu$, neglecting the effect of input saturation.
Hence, the ideal control input $\rbu^*$ can be designed as $-\mv f$.
However, $\rbu^*$ is not available in practice, since the system dynamics encapsulated in $\mv f$ are unknown.
}\DIFdelend \DIFaddbegin \begin{figure}[t]
    \centering
    \includegraphics[width=0.99\linewidth]{\src/figures/DNN.drawio.pdf}
    \caption{
        \DIFaddFL{Architecture of the DNN $\NN(\NNin;\wth)$ with $k=2,l_0=2,l_1=l_2=3$, and $l_3=2$.
    }}
    \label{fig:DNN}
\end{figure}
\DIFaddend 

\DIFdelbegin \DIFdel{To overcome this issue, a DNN is employed to approximate $\rbu^*$.
Let $\NN:=\NN({\q}_n;\wth): \R^{l_0+1}\times\R^{\Xi}\to\R^{n}$ represent the DNN , where ${\q}_n\in\R^{l_0+1}$ is the DNN input vector, and $\wth\in\R^{\Xi}$ is the vector of trainable weights.
The architecture of DNN $\NN({\q}_n;\wth)$ will be presented in Section \ref{sec:sub:NN definition}, later.
}\DIFdelend %DIF >  ————————————————————————————————
%DIF >  DNN Preliminaries
%DIF >  ————————————————————————————————
\DIFaddbegin 

\DIFadd{We review the DNN structure and its approximation capability, which are essential for the subsequent controller design and stability analysis.
}\DIFaddend According to the universal approximation \DIFdelbegin \DIFdel{theorem for DNNs \mbox{%DIFAUXCMD
\cite{Kidger:2020aa}}\hskip0pt%DIFAUXCMD
, $\NN({\q}_n;\wth)$ can approximate a nonlinear function, denoted as $\mv g(\cdot)$, with an ideal weight vector $\idealwth\in\R^\Xi$ over a compact subset $\Omega_{n}\in\R^{l_0+1}$ to within $\epsilon$-accuracy, such that 
$\sup_{{\q}_n\in\Omega_{n}}\norm{ \NN({\q}_n;\idealwth) - \mv g(\cdot)} = \epsilon < \infty$.
The compact set $\Omega_{n}$ is defined as the set of all feasible states $\q,\dq$ and filtered errors $\fe$.
%DIF <  Furthermore, the theorem states that the norm of $\idealwth$ is bounded such that $\norm{\idealwth}\le \overline\wth<\infty$.
In this study, the ideal weight vector$\idealwth$ is assumed to be bounded by Assumption \ref{assum:feasible} and \mbox{%DIFAUXCMD
\cite[Assum. 1]{Lewis:1996aa}}\hskip0pt%DIFAUXCMD
, and considered a locally optimal solution, rather than a global one.
}%DIFDELCMD < 

%DIFDELCMD < %%%
\DIFdel{Accordingly, the ideal control input $\rbu^*$ can be approximated by the DNN with the ideal weight vector $\idealNN:=\NN({\q}_n;\idealwth)$ as follows:
}%DIFDELCMD < \begin{align}
%DIFDELCMD <     \rbu^*=& \idealNN+\mv\epsilon
%DIFDELCMD <     ,
%DIFDELCMD <     \label{eq:control:ideal}
%DIFDELCMD < \end{align}%%%
\DIFdelend \DIFaddbegin \DIFadd{property \mbox{%DIFAUXCMD
\cite{Patil:2022aa}}\hskip0pt%DIFAUXCMD
, DNNs $\NN: \mathbb{R}^{l_0+1} \times \mathbb{R}^{\Xi} \to \mathbb{R}^{l_k+1}$ can be used to approximate any smooth unknown nonlinear functions $\mv{g}:\R^{l_0+1}\to\R^{l_{k+1}}$.
Specifically, there exists a DNN such that 
}\begin{equation}
    \sup_{\NNin\in\Omega_{\NNin}}
    \norm{
        \mv{g}(\NNin) - \NN(\NNin;\wth^*)
    }
    =
    \overline{\epsilon}
    ,
\end{equation}\DIFaddend 
where \DIFdelbegin \DIFdel{$\mv\epsilon\in\R^n$ is the approximation error vector, bounded by $\norm{\mv\epsilon}\le \overline{\epsilon}$ for some $\overline{\epsilon}\in\R_{>0}$}\DIFdelend \DIFaddbegin \DIFadd{$\NNin \in \mathbb{R}^{l_0+1}$ is the DNN input vector, $\Omega_{\NNin} \subset \mathbb{R}^{l_0+1}$ is a compact set, and $\idealwth \in \mathbb{R}^{\Xi}$ denotes the ideal DNN weights that minimize the approximation error $\overline{\epsilon} > 0$}\DIFaddend .
The ideal \DIFdelbegin \DIFdel{control input $\rbu^*$ is estimated online by
}%DIFDELCMD < \begin{align}
%DIFDELCMD <     \rbu =& \estNN,
%DIFDELCMD <     \label{eq:control:est}
%DIFDELCMD < \end{align}
%DIFDELCMD < %%%
\DIFdel{where $\hat\NN:=\NN({\q}_n;\estwth)$, and $\estwth\in\R^\Xi$ is the estimated weight vector for the ideal weight vector }\DIFdelend \DIFaddbegin \DIFadd{DNN weights }\DIFaddend $\idealwth$ \DIFdelbegin \DIFdel{.
}%DIFDELCMD < 

%DIFDELCMD < %%%
\DIFdel{Substituting \eqref{eq:control:ideal} and \eqref{eq:control:est} into \eqref{eq:lya1:dot1}, the time derivative of $V_1$ can be rewritten as
}%DIFDELCMD < \begin{equation}
%DIFDELCMD <     \ddtt{V}_1
%DIFDELCMD <     \le 
%DIFDELCMD <     -\lambda_{\min}(\mm K)\norm{\fe}^2
%DIFDELCMD <     +
%DIFDELCMD <     \overline\tau_d\norm{\fe}
%DIFDELCMD <     +
%DIFDELCMD <     \fe^\top 
%DIFDELCMD <     \left(
%DIFDELCMD <         -\idealNN
%DIFDELCMD <         -\mv\epsilon
%DIFDELCMD <         +\mysat(\hat\NN)
%DIFDELCMD <     \right)
%DIFDELCMD <     .
%DIFDELCMD <     \label{eq:lya1:dot2}
%DIFDELCMD < \end{equation}
%DIFDELCMD < %%%
\DIFdel{Therefore, it is concluded that the filtered error $\fe$ can be stabilized by adapting $\estwth$ to $\idealwth$, }%DIFDELCMD < \ie %%%
\DIFdel{$\hat\NN\to\idealNN$, neglecting the effect of input saturation.
Note that the control input saturation has not yet been considered in the derivation of the ideal control input; this will be addressed in the subsequent section.
}%DIFDELCMD < 

%DIFDELCMD < %%%
\subsection{\DIFdel{Deep Neural Network (DNN) Model}}%DIFAUXCMD
\addtocounter{subsection}{-1}%DIFAUXCMD
%DIFDELCMD < \label{sec:sub:NN definition}
%DIFDELCMD < 

%DIFDELCMD < \begin{figure}[t]
%DIFDELCMD <     \centering
%DIFDELCMD <     \includegraphics[width=0.99\linewidth]{\src/figures/DNN.drawio.pdf}
%DIFDELCMD <     %%%
%DIFDELCMD < \caption{%
{%DIFAUXCMD
\DIFdelFL{Architecture of the DNN $\NN({\q}_n;\wth)$ with $k=2,l_0=2,l_1=l_2=3$, and $l_3=2$.
    }}
    %DIFAUXCMD
%DIFDELCMD < \label{fig:DNN}
%DIFDELCMD < \end{figure}
%DIFDELCMD < 

%DIFDELCMD < %%%
\DIFdel{The DNN architecture $\NN({\q}_n;\wth) = \NN_k$ }\DIFdelend \DIFaddbegin \DIFadd{are assumed to be bounded.
The DNN }\DIFaddend can be recursively \DIFdelbegin \DIFdel{represented, as
follows:
}%DIFDELCMD < \begin{equation}
%DIFDELCMD <     \NN_i :=
%DIFDELCMD <     \begin{cases}
%DIFDELCMD <         \wV_i^\top \act_i(\NN_{i-1}), 
%DIFDELCMD <         &
%DIFDELCMD <         i\in\{1,\dots ,k\},
%DIFDELCMD <         \\
%DIFDELCMD <         \wV_0^\top {\q}_n,
%DIFDELCMD <         &
%DIFDELCMD <         i=0
%DIFDELCMD <         ,
%DIFDELCMD <     \end{cases}
%DIFDELCMD <     \label{eq:DNN:architecture}
%DIFDELCMD < \end{equation}%%%
\DIFdelend \DIFaddbegin \DIFadd{defined as
}\begin{equation} \label{eq:dnn:def}
    \NN_i :=
    \begin{cases}
        \wV_i^\top \act_i(\NN_{i-1}), 
        &
        i\in\{1,\dots ,k\},
        \\
        \wV_0^\top \NNin,
        &
        i=0
        ,
    \end{cases}
\end{equation}\DIFaddend 
where \DIFdelbegin \DIFdel{$\wV_i=[w_{ij}]_{i\in\{1,\cdots,l_i+1\},j\in\{1,\cdots,l_{i+1}\}}\in\R^{(l_i+1)\times l_{i+1}}$ }\DIFdelend \DIFaddbegin \DIFadd{$\wV_i\in\R^{(l_i+1)\times l_{i+1}}$ }\DIFaddend and $\act_i: \R^{l_i}\to\R^{l_i+1}$ represent the weight matrix and the activation function of the $i$th layer \DIFdelbegin \DIFdel{, respectively}\DIFdelend \DIFaddbegin \DIFadd{for all $i\in\mathcal{K}$, respectively, and $\mathcal{K}:=\{0,1,\cdots,k\}$ denotes the index set of layers}\DIFaddend .
The DNN \DIFdelbegin \DIFdel{$\NN({\q}_n;\wth)$ }\DIFdelend \DIFaddbegin \DIFadd{$\NN(\NNin;\wth)$ }\DIFaddend has $k+1$ layers, where the input \DIFdelbegin \DIFdel{layer is indexed by }\DIFdelend \DIFaddbegin \DIFadd{and output layers are indexed by }\DIFaddend $i=0$ \DIFdelbegin \DIFdel{, and the output layer is indexed by }\DIFdelend \DIFaddbegin \DIFadd{and }\DIFaddend $i=k$\DIFdelbegin \DIFdel{.
For instance, }\DIFdelend \DIFaddbegin \DIFadd{, respectively, as illustrated }\DIFaddend in Fig.~\ref{fig:DNN} \DIFdelbegin \DIFdel{, the DNN $\NN({\q}_n;\wth)$ with $k=2,l_0=2,l_1=l_2=3$, and $l_3=2$ is illustrated.
Notice that the output size of $\NN(\cdot)$ is the same as that of the control input $\rbu$, }%DIFDELCMD < \ie %%%
\DIFdel{$l_{k+1}=n$.
}%DIFDELCMD < 

%DIFDELCMD < %%%
\DIFdel{The DNN input vector ${\q}_n$ is defined as ${\q}_n=(\q^\top,\dq^\top,\fe^\top,1)^\top\in\R^{3n+1}$ to include the current state of the system, the filtered error, and the augmentation term $1$ to account for the bias term in the input weight matrix $\wV_0$.
The }\DIFdelend \DIFaddbegin \DIFadd{as an example.
The }\DIFaddend activation function is defined as \DIFdelbegin \DIFdel{$\act_i(\mv x)=\left(\sigma(x_1),\sigma(x_2),\cdots, \sigma(x_{l_{i}}), 1\right)^\top,\ \forall\mv x\in\R^{l_i}$, where $\sigma: \R\to\R$ is a nonlinear function, and the augmentation of $1$ is used to account for bias terms in the weight matrices, similar to the input vector ${\q}_n$.
The weights $\wV_i,\ \forall i\in\{0,\cdots,k\}$, are initialized randomly using a uniform distribution at the beginning of the control process.
}%DIFDELCMD < 

%DIFDELCMD < %%%
%DIF <  \MSRY{생략 가능할 것 같습니다.}
%DIF <  \ctxt{red}
%DIF <  {
    \DIFdel{In selecting the nonlinear function $\sigma(\cdot)$, the ReLU function \mbox{%DIFAUXCMD
\cite{Maas:2013aa} }\hskip0pt%DIFAUXCMD
is one of the most widely used choices for large DNNs, since it effectively mitigates the gradient vanishing problem during error backpropagation. 
    However, for control applications, where relatively shallow DNNs are typically sufficient and the gradient vanishing issue is less severe, the sigmoid function or the hyperbolic tangent function is commonly used as the nonlinear function $\sigma(\cdot)$. 
    These functions simplify stability analysis, since they are continuously differentiable andtheir outputs and gradients are bounded. 
%DIF <  }
In this study, the hyperbolic tangent function $\tanh(\cdot)$ was selected as the nonlinear function, }%DIFDELCMD < \ie %%%
\DIFdel{$\sigma(\cdot) = \tanh(\cdot)$, which provides desirable boundedness with $\norm{\sigma(x)}<1$ }\DIFdelend \DIFaddbegin \DIFadd{$\act_i(\mv{z}) = \left(\tanh(z_1),\cdots,\tanh(z_{l_i}),1\right)^\top,$ for some $\mv{z}\in\R^{l_i}$.
We simplify the notations by vectorizing all weight matrices as $\wth_i := \myvec(\wV_i)\in\R^{(l_i+1)l_{i+1}},\ \forall i\in\mathcal{K}$, }\DIFaddend and\DIFdelbegin \DIFdel{$0<\norm{\ddtfrac{\sigma(x)}{x}}\le 1$}\DIFdelend \DIFaddbegin \DIFadd{, finally, stacking them as $\wth := \left[\wth_i\right]_{i\in\mathcal{K}}$, where $\Xi := \sum_{i\in\mathcal{K}} (l_i+1)l_{i+1}$.
Accordingly, the boundedness of the ideal DNN weights $\idealwth$ implies the boundedness of each $\idealwth_i$, such that $\norm{\idealwth_i}\le\overline{\theta}^*_i$, $\forall i\in\mathcal{K}$, for some $\overline{\theta}^*_i\in\R_{>0}$}\DIFaddend .

\DIFdelbegin \DIFdel{For simplicity, each layer's weights are vectorized as $\wth_i:=\myvec(\wV_i)\in\R^{\Xi_i},\ \forall i\in\{0,\cdots,k\}$, where $\Xi_i:= (l_i+1)l_{i+1}$ is the number of weights in the $i$th layer.
The total weight vector $\wth\in\R^{\Xi}$ is defined by augmenting $\wth_i, \forall i\in \{0,\cdots,k\}$, as 
}%DIFDELCMD < \begin{equation}
%DIFDELCMD <     \wth := 
%DIFDELCMD <     \begin{pmatrix}
%DIFDELCMD <         \wth_k\\
%DIFDELCMD <         \wth_{k-1}\\
%DIFDELCMD <         \vdots\\
%DIFDELCMD <         \wth_0
%DIFDELCMD <     \end{pmatrix}
%DIFDELCMD <     =
%DIFDELCMD <     \begin{pmatrix}
%DIFDELCMD <         \myvec(\wV_k)\\
%DIFDELCMD <         \myvec(\wV_{k-1})\\
%DIFDELCMD <         \vdots\\
%DIFDELCMD <         \myvec(\wV_0)
%DIFDELCMD <     \end{pmatrix},
%DIFDELCMD < \end{equation}
%DIFDELCMD < %%%
\DIFdel{where $\Xi:={\sum_{i=0}^{k} \Xi_i}$ represents the total number of weights in the DNN.
}%DIFDELCMD < 

%DIFDELCMD < %%%
\DIFdelend The gradient of \DIFdelbegin \DIFdel{$ \NN({\q}_n;\wth)$ }\DIFdelend \DIFaddbegin \DIFadd{$ \NN(\NNin;\wth)$ }\DIFaddend with respect to $\wth$ can be obtained using Proposition\DIFaddbegin \DIFadd{~}\DIFaddend \ref{propsit:kron} and the chain rule as follows:
\DIFdelbegin %DIFDELCMD < \begin{equation}
%DIFDELCMD <     \pptfrac{\NN}{\wth}=
%DIFDELCMD <     \begin{bmatrix}
%DIFDELCMD <         \pptfrac{\NN}{\wth_k}&
%DIFDELCMD <         \pptfrac{\NN}{\wth_{k-1}}&
%DIFDELCMD <     \cdots &
%DIFDELCMD <         \pptfrac{\NN}{\wth_0}
%DIFDELCMD <     \end{bmatrix}
%DIFDELCMD <     \in\R^{n \times \Xi}
%DIFDELCMD <     ,
%DIFDELCMD <     \label{eq:gradient:phi}
%DIFDELCMD < \end{equation}%%%
\DIFdelend \DIFaddbegin \begin{equation}
    \ppfrac{\NN}{\wth}=
    \begin{bmatrix}
        \ppdfrac{\NN}{\wth_k}&
        \ppdfrac{\NN}{\wth_{k-1}}&
    \cdots &
        \ppdfrac{\NN}{\wth_0}
    \end{bmatrix}
    \in\R^{l_{k+1} \times \Xi}
    ,
    \label{eq:gradient:phi}
\end{equation}\DIFaddend 
where
\DIFdelbegin %DIFDELCMD < \begin{equation}
%DIFDELCMD <     \pptfrac{\NN}{\wth_i} = 
%DIFDELCMD <     \begin{cases}
%DIFDELCMD <         (
%DIFDELCMD <             \mm I_{l_{k+1}}
%DIFDELCMD <             \otimes 
%DIFDELCMD <             \act_{k}^\top  
%DIFDELCMD <         )
%DIFDELCMD <         , 
%DIFDELCMD <         &
%DIFDELCMD <         i=k
%DIFDELCMD <         ,
%DIFDELCMD <         \\
%DIFDELCMD <         \wV_k^\top\act_{k}' 
%DIFDELCMD <         (
%DIFDELCMD <             \mm I_{l_{k}}
%DIFDELCMD <             \otimes  
%DIFDELCMD <             \act_{k-1}^\top  
%DIFDELCMD <         )
%DIFDELCMD <         , 
%DIFDELCMD <         & 
%DIFDELCMD <         i=k-1
%DIFDELCMD <         ,
%DIFDELCMD <         \\
%DIFDELCMD <         \vdots 
%DIFDELCMD <         &
%DIFDELCMD <         \vdots 
%DIFDELCMD <         \\
%DIFDELCMD <         \wV_k^\top\act'_{k} 
%DIFDELCMD <         \cdots 
%DIFDELCMD <         \wV_1^\top\act_1' 
%DIFDELCMD <         (
%DIFDELCMD <             \mm I_{l_1}
%DIFDELCMD <             \otimes 
%DIFDELCMD <             {\q}_n^\top  
%DIFDELCMD <         )
%DIFDELCMD <         , 
%DIFDELCMD <         &
%DIFDELCMD <         i = 0
%DIFDELCMD <         ,
%DIFDELCMD <     \end{cases}
%DIFDELCMD <     \label{eq:NN:grad}
%DIFDELCMD < \end{equation}%%%
\DIFdelend \DIFaddbegin \begin{equation}
    \ppfrac{\NN}{\wth_i} = 
    \begin{cases}
        % (
            \mm I_{l_{k+1}}
            \otimes 
            \act_{k}^\top  
        % )
        , 
        &
        i=k
        ,
        \\
        \wV_k^\top\act_{k}' 
        (
            \mm I_{l_{k}}
            \otimes  
            \act_{k-1}^\top  
        )
        , 
        & 
        i=k-1
        ,
        \\
        \vdots 
        &
        \vdots 
        \\
        \wV_k^\top\act'_{k} 
        \cdots 
        \wV_c^\top\act_1' 
        (
            \mm I_{l_1}
            \otimes 
            \NNin^\top  
        )
        , 
        &
        i = 0
        ,
    \end{cases}
    \label{eq:NN:grad}
\end{equation}\DIFaddend 
and $\act_i:= \act_i(\NN_{i-1})$ and $\act_i':= \ddtfrac{\act_i}{\NN_{i-1}}$\DIFaddbegin \DIFadd{, for all $i\in\mathcal{K}\setminus\{0\}$}\DIFaddend .

\DIFdelbegin \DIFdel{In the following sections, let $\idealNN_i$ represent the output of }\DIFdelend \DIFaddbegin \subsection{\DIFadd{Neuro-Adaptive Control Design}} \label{sec:main:result:sub:ctrl:design}

\DIFadd{We define the tracking error as $\mv{e} = \mv{x} - \mv{x}_d$, and derive the error dynamics from \eqref{eq:sys} as
}\begin{equation} \label{eq:error:dyn}
    \ddtt\mv{e}
    =
    \mv{f}(\mv{x}) + \mm{B}(\mv{x})\mysat(\mv{u}) - \ddtt\mv{x}_d
    .
\end{equation}
\DIFadd{Let us consider Lyapunov function $V_c:=\tfrac{1}{2}\mv{e}^\top\mm{B}^{-1}(\mv{x})\mv{e}$.
Taking time derivative of $V_c$ yields
}\begin{align} \label{eq:Vc:dot:1} \nonumber
    \ddtt V_c
    =&
    \tfrac{1}{2}\mv{e}^\top
    \left(
        \ddtt\mm{B}^{-1}
    \right)
    \mv{e}
    +
    \mv{e}^\top\mm{B}^{-1}
    \left(
        \mv{f} + \mm{B}\mysat(\mv{u}) 
        - \ddtt\mv{x}_d
    \right)
    \\
    \le&
    \beta(\mv{x})\norm{\mv{e}}^2
    \!+\!
    \mv{e}^\top\mv{u}
    \!+\!
    \mv{e}^\top\mm{B}^{-1}
    \left(
        \mv{f} \!-\! \ddtt\mv{x}_d
    \right)
    \!+\!
    \mv{e}^\top
    \Delta\mv{u}
    ,
\end{align}
\DIFadd{where arguments of $\mv{f}(\mv{x})$ and $\mm{B}(\mv{x})$ are omitted for brevity and $\Delta\mv{u} := \mysat(\mv{u}) - \mv{u}$ denotes amount of input saturation.
}

\DIFadd{A desired control law $\mv{u}_d$ such that the tracking error $\mv{e}$ is stable in the absence of the input saturation, is selected as 
}\begin{equation} \label{eq:ctrl:ideal}
    \mv{u}_d
    :=
    -
    \mm{P}
    \mv{e}
    - 
    \beta(\mv{x})\mv{e}
    +
    \mm{B}^{-1}
    \left(
        - \mv{f} 
        + \ddtt\mv{x}_d
    \right)
    % +
    % h(\mv{x},\mv{x}_d)
    ,
\end{equation}
\DIFadd{where user-defined matrix $\mm{P}=\mm{P}^\top\in\R^{n\times n}$ is a symmetric positive definite matrix satisfying $\underline{p}\mm{I}_n \preceq \mm{P} \preceq \overline{p}\mm{I}_n$ for known positive constants $\underline{p},\overline{p}\in\R_{>0}$.
%DIF >  The function $h:\R^n\times\R^n\to\R^n$ is correction term to satisfy the ideal control law $\mv{u}^*$ within the input constraints, \ie $\mv{u}^*\in\Omega_u$.
}

\DIFadd{By substituting $\mv{u}_d$ \eqref{eq:ctrl:ideal} into \eqref{eq:Vc:dot:1}, we obtain
}\begin{align} \label{eq:Vc:dot:2} 
    \ddtt V_c
    \le &
    -\mv{e}^\top
    \mm{P}
    \mv{e}
    +
    \mv{e}^\top
    \left(
        \mv{u} - \mv{u}_d
    \right)
    +
    \mv{e}^\top
    \Delta\mv{u}
    \\ \nonumber
    \le & 
    -
    \underline{p}\norm{\mv{e}}^2
    +
    \mv{e}^\top
    \left(
        \mv{u} - \mv{u}_d       
    \right)
    +
    \mv{e}^\top
    % \left(
        \Delta\mv{u}
    %     +
    %     h(\mv{x},\mv{x}_d)
    % \right)
    .
\end{align}
\DIFadd{Notably, in ideal case, }\ie \DIFadd{$\mv{u} = \mv{u}_d$ and $\Delta\mv{u} = \mv{0}$, }\DIFaddend the \DIFdelbegin \DIFdel{$i$th layer with the ideal weight vector $\idealwth$.
Define $\idealact_i=\act_i(\idealNN_{i-1})$ and ${\idealact}'_i= \ddtfrac{\idealact_i}{\idealNN_{i-1}}$. 
Similarly, let $\estNN_i$ denote the output of the $i$th layer with the estimated weight vector $\estwth$, }\DIFdelend \DIFaddbegin \DIFadd{tracking error $\mv{e}$ is asymptotically exponentially stable with decay rate $\underline{p}\overline{b}$, since 
$
    \tfrac{1}{\overline{b}}
    \ddtt
    \left(
        \norm{\mv{e}}^2
    \right)
    \le
    \ddtt V_c 
    \le 
    -\underline{p}\norm{\mv{e}}^2
$.
}

%DIF >  It is notable that \eqref{eq:Vc:dot:2} can be rewritten as 
%DIF >  $
%DIF >      \ddtt V_c 
%DIF >      \le 
%DIF >      -\underline{p}\norm{\mv{e}}^2 
%DIF >      + 
%DIF >      \norm{
%DIF >          \mysat(\mv{u})
%DIF >          -
%DIF >          \mysat(\mv{u}_d)
%DIF >      }
%DIF >      \norm{\mv{e}}
%DIF >      \le 
%DIF >      -
%DIF >      \left(
%DIF >          \underline{p} - 1
%DIF >      \right)
%DIF >      \norm{\mv{e}}^2 
%DIF >      +
%DIF >      \overline{u}^2
%DIF >      ,
%DIF >  $
%DIF >  which concludes that, according to Lemma \ref{lem:comp:compact}, if $\underline{p} > 1$, the tracking error $\mv{e}$ is bounded within a compact set $\Omega_e := \left\{\mv{e} : \norm{\mv{e}}\le\overline{e}\in\R_{>0}\right\}$, where 
%DIF >  $
%DIF >      \sqrt{
%DIF >          \tfrac{\overline{b}}{\underline{b}}
%DIF >          \norm{\mv{e}(0)}^2
%DIF >          +
%DIF >          \tfrac
%DIF >          {\overline{u}^2}
%DIF >          {\underline{p} - 1}
%DIF >          \tfrac
%DIF >          {\overline{b}}
%DIF >          {\underline{b}}
%DIF >      }
%DIF >      \le
%DIF >      \overline{e}
%DIF >      ,
%DIF >  $
%DIF >  which only holds if $\mv{e}$ is initialized within certain bounds, $\norm{e}(0)\le\sqrt{\tfrac{\underline{b}}{\overline{b}}\overline{e}^2 - \tfrac{\overline{u}^2}{\underline{p}-1}}$.
%DIF >  This implies that $\mv{x}$ remains in a compact set $\Omega_x := \left\{\mv{x} : \norm{\mv{x}}\le \overline{e} + \overline{x}_d\right\}$, since $\mv{x}_d$ is bounded in $\Omega_{x_d}$, and the compact set $\Omega_x$ can be interpreted as physically limited domain of the system states in practice.

\DIFadd{Since the ideal control law \eqref{eq:ctrl:ideal} contains unknown functions $\beta(\mv{x})$, $\mv{f}$ }\DIFaddend and \DIFdelbegin \DIFdel{define $\estact_i=\act_i(\estNN_{i-1})$ and $\estact_i'= \ddtfrac{\estact_i}{\estNN_{i-1}}$}\DIFdelend \DIFaddbegin \DIFadd{$\mm{B}$, we employ a DNN $\NN(\NNin;\wth)$ in \eqref{eq:dnn:def} to approximate $\mv{u}_d$, where the DNN input vector is defined as $\NNin := \left(\mv{x}^\top , \mv{x}_d^\top , \ddtt{\mv{x}}_d^\top, 1 \right)^\top\in\R^{3n+1}$.
Then, we have 
}\begin{equation} \label{eq:ctrl:NN:ideal}
    \mv{u}_d
    =
    \mm{P}
    \big(
        \NN(\NNin;\idealwth)
        +
        \mv{\epsilon}
    \big)
    ,
\end{equation}
\DIFadd{where $\mv{\epsilon}\in\R^n$ denotes an approximation error vector, such that 
$
    \inf_{\NNin\in\Omega_{\NNin}}   
    \left( 
        \norm{\mv{\epsilon}}
    \right)
    =
    \overline{\epsilon}
    .
$
The actual control law is designed as 
}\begin{equation} \label{eq:ctrl:NN:actual}
    \mv{u}
    :=
    \mm{P}
    \NN(\NNin;\estwth)
    ,
\end{equation}
\DIFadd{where $\estwth\in\R^{\Xi}$ denotes the estimated DNN weights.
}

\DIFadd{By substituting \eqref{eq:ctrl:NN:ideal} and \eqref{eq:ctrl:NN:actual} into \eqref{eq:Vc:dot:2}, we obtain
%DIF >  \begin{align} \label{eq:Vc:dot:3} 
}\begin{equation} \label{eq:Vc:dot:3}
    \ddtt V_c
    \le
    - \underline{p}
    \norm{\mv{e}}^2
    +
    \mv{e}^\top\mm{P}
    \left(
        \estNN - \idealNN
        - \mv{\epsilon}
    \right)
    +
    \mv{e}^\top
    \Delta\mv{u}
    ,
% \end{align}
\end{equation}
\DIFadd{where $\estNN\!:=\!\NN(\NNin;\estwth)$ and $\idealNN\!:=\!\NN(\NNin;\wth^*)$ are abbreviated}\DIFaddend .

%DIF <   SECTION ADAPTATION LAW DERIVATION =======================
\DIFdelbegin \section{\DIFdel{Weight Adaptation Laws}}%DIFAUXCMD
\addtocounter{section}{-1}%DIFAUXCMD
%DIFDELCMD < \label{sec:adap_laws}
%DIFDELCMD < %%%
\DIFdelend \DIFaddbegin \subsection{\DIFadd{Adaptation Law Derivation}} \label{sec:main:result:sub:adap:laws}
\DIFaddend 

%DIF <  \subsection{Adaptation Law using Lagrangian Function}
\DIFdelbegin \subsection{\DIFdel{Weight Optimizer Design}}%DIFAUXCMD
\addtocounter{subsection}{-1}%DIFAUXCMD
%DIFDELCMD < \label{sec:sub:weight optimizer}
%DIFDELCMD < 

%DIFDELCMD < %%%
\DIFdel{The control objective can be represented }\DIFdelend %DIF >  The control objective which minimizes the tracking error $\mv{e}$ while satisfying both weight boundedness and input constraints can be represented as follows:
\DIFaddbegin \DIFadd{The control objectives defined in Section~\ref{sec:problem:formulation} can be formulated in a constrained optimization problem }\DIFaddend as follows:
\DIFdelbegin %DIFDELCMD < \begin{equation}
%DIFDELCMD <     \begin{matrix}
%DIFDELCMD <         \min_{\estwth} \ J(\fe;\estwth):= 
%DIFDELCMD <         \tfrac{1}{2} \fe^\top \fe
%DIFDELCMD <         ,
%DIFDELCMD <         \vspace{.5em}
%DIFDELCMD <         \\
%DIFDELCMD <         \begin{aligned}
%DIFDELCMD <         \text{subject to }&c_{j}(\estwth) 
%DIFDELCMD <         \le0, \quad \forall j\in\mathcal{I},
%DIFDELCMD <         \end{aligned}
%DIFDELCMD <     \end{matrix}
%DIFDELCMD <     \label{eq:opt:problem1}
%DIFDELCMD < \end{equation}%%%
\DIFdelend \DIFaddbegin \begin{subequations} \label{eq:opt:prob} 
    \begin{align} 
        \min_{\estwth} \quad 
        & 
        J(\mv{e};\estwth)= \tfrac{1}{2} \mv{e}^\top \mv{e} \label{eq:opt:prob:obj} 
        \\ 
        \text{subject to} \quad 
        & 
        c^{\theta}_i(\estwth_i)=\tfrac{1}{2}\norm{\estwth_i}^2 - \tfrac{1}{2}\overline{\theta}_i^2 \le 0, 
        \  \forall i\in\mathcal{K} \label{eq:opt:prob:const1} 
        \\ 
        & 
        c^{u}_j(\estwth)
        % c_j^u
        % \left(
        %     \mm{P}\NN(\NNin;\estwth)
        % \right) 
        \le 0, 
        \quad \forall j\in\mathcal{I} \label{eq:opt:prob:const2} 
        ,
    \end{align}\DIFaddend \DIFdelbegin \DIFdel{where $J(\fe;\estwth)$ }\DIFdelend %DIF > 
\DIFaddbegin \end{subequations}
\DIFadd{where $J(\mv{e};\estwth)$ }\DIFaddend is the objective function.
Inequality constraints \DIFdelbegin \DIFdel{$c_j:=c_j(\estwth),\ \forall j\in\mathcal{I}$, are imposed during the weight adaptation process, where $\mathcal I$ denotes the set of }\DIFdelend \DIFaddbegin \DIFadd{of weight boundedness $c^{\theta}_i(\estwth_i),\ \forall i\in\mathcal{K},$ and input saturation $c^{u}_j(\estwth),\ \forall j\in\mathcal{I},$ are considered.
}

\DIFadd{Specifically, the input constraints used in \eqref{eq:opt:prob:const2} are }\DIFaddend the \DIFdelbegin \DIFdel{imposed inequality constraints . 
%DIF <  Here, $\fe$ is considered a pre-defined data or static parameter for this optimization problem, so that the adaptation process does not consider the dynamics of $\fe$.
}\DIFdelend \DIFaddbegin \DIFadd{same convex constraint functions that characterize the admissible input set in Assumption~\ref{assum:sat}.
With a slight abuse of notation, the composed input constraint is denoted by
$c_j^u(\estwth):=c_j^u(\mv u(\estwth))=c_j^u(\mm P\NN(\NNin;\estwth)),\ \forall j\in\mathcal I$, to explicitly represent the dependence of the input constraints on the DNN weights $\estwth$.
Additionally, all active constraints are assumed to satisfy the linear independence constraint qualification (LICQ) to avoid degeneracy \mbox{%DIFAUXCMD
\cite[pp. 465--467]{Nocedal:2006aa}}\hskip0pt%DIFAUXCMD
.
}

\DIFaddend The Lagrangian function \DIFdelbegin \DIFdel{$L(\cdot)$ }\DIFdelend \DIFaddbegin \DIFadd{$L:=L(\mv{e},\estwth,[\sigma_i]_{i\in\mathcal{K}},[\lambda_j]_{j\in\mathcal I})$ }\DIFaddend is defined as
\DIFdelbegin %DIFDELCMD < \begin{equation}
%DIFDELCMD <     L(
%DIFDELCMD <         \fe,\estwth,[\lambda_j]_{j\in\mathcal I}
%DIFDELCMD <     ) 
%DIFDELCMD <     = 
%DIFDELCMD <     J(\fe;\estwth) 
%DIFDELCMD <     + 
%DIFDELCMD <     \textstyle\sum_{j\in\mathcal I}
%DIFDELCMD <     \lambda_{j}
%DIFDELCMD <     c_{j}(\estwth)
%DIFDELCMD <     ,
%DIFDELCMD <     \label{eq:lagrangian}
%DIFDELCMD < \end{equation}%%%
\DIFdelend \DIFaddbegin \begin{equation}
    L
    = 
    J(\mv{e};\estwth) 
    + 
    \textstyle\sum_{i\in\mathcal{K}}
    \sigma_{i}
    c^{\theta}_{i}(\estwth_i)
    +
    \textstyle\sum_{j\in\mathcal I}
    \lambda_{j}
    c^{u}_{j}(\estwth)
    ,
    \label{eq:lagrangian}
\end{equation}\DIFaddend 
where \DIFdelbegin \DIFdel{$\lambda_j,\ \forall j\in\mathcal{I}$, denotes the Lagrange multiplier for each constraint.
}\DIFdelend \DIFaddbegin \DIFadd{$[\sigma_i]_{i\in\mathcal{K}}$ and $[\lambda_j]_{j\in\mathcal{I}}$ denote the Lagrange multipliers for each weight \eqref{eq:opt:prob:const1} and input constraints \eqref{eq:opt:prob:const2}, respectively.
Then, the primal-dual saddle-point problem can be formulated as follows:
}\begin{equation} \label{eq:opt:prob:saddle}
    \min_{\estwth}
    \max_{
        [\sigma_i]_{i\in\mathcal{K}}
        ,
        [\lambda_j]_{j\in\mathcal{I}}
    } 
    L(\mv{e},\estwth,[\sigma_i]_{i\in\mathcal{K}},[\lambda_j]_{j\in\mathcal I})
    .
\end{equation}
\DIFaddend 

The adaptation laws for \DIFaddbegin \DIFadd{the adaptive variables, }\DIFaddend $\estwth$\DIFaddbegin \DIFadd{, and $[\sigma_i]_{i\in\mathcal{K}}$ }\DIFaddend and \DIFdelbegin \DIFdel{$[\lambda_j]_{j\in\mathcal I}$ }\DIFdelend \DIFaddbegin \DIFadd{$[\lambda_j]_{j\in\mathcal{I}}$, }\DIFaddend are derived to solve \DIFdelbegin \DIFdel{the dual problem of \eqref{eq:opt:problem1}}\DIFdelend \DIFaddbegin \DIFadd{\eqref{eq:opt:prob:saddle} by using a primal-dual gradient update law}\DIFaddend , \ie 
\DIFdelbegin \DIFdel{$\min_{\estwth} \max_{[\lambda_j]_{j\in\mathcal I}}L(\fe,\estwth,[\lambda_j]_{j\in\mathcal I})$, as follows:
}\DIFdelend \DIFaddbegin \DIFadd{$
    \ddtt\estwth 
    = 
    -\alpha \pptfrac{L}{\estwth}
$, $
    \ddtt\sigma_i 
    = 
    \myproj_{\sigma_i\ge0}
    \left(
        \beta_i^\theta 
        \pptfrac{L}{\sigma_i}
    \right)
$, and $
    \ddtt\lambda_j 
    = 
    \myproj_{\lambda_j\ge0}
    \left(
        \beta_j^u 
        \pptfrac{L}{\lambda_j}
    \right)
$, 
}\DIFaddend \begin{subequations} \DIFdelbegin %DIFDELCMD < \begin{align}
%DIFDELCMD <             \ddtt{\estwth}
%DIFDELCMD <             &
%DIFDELCMD <             =
%DIFDELCMD <             -\alpha \pptfrac{L}{\estwth}
%DIFDELCMD <             =
%DIFDELCMD <             -\alpha 
%DIFDELCMD <             \left(
%DIFDELCMD <                 \pptfrac{J}{\estwth}
%DIFDELCMD <                 +
%DIFDELCMD <                 \textstyle\sum_{j\in\mathcal{I}}
%DIFDELCMD <                 \lambda_j 
%DIFDELCMD <                 \pptfrac{c_j}{\estwth}
%DIFDELCMD <             \right),
%DIFDELCMD <         \label{eq:adap:th}
%DIFDELCMD <             \\
%DIFDELCMD <             \ddtt\lambda_j
%DIFDELCMD <             & 
%DIFDELCMD <             = 
%DIFDELCMD <             \beta_j\pptfrac{L}{\lambda_j} 
%DIFDELCMD <             = 
%DIFDELCMD <             \beta_j c_j ,
%DIFDELCMD <             \quad \forall j\in\mathcal I,
%DIFDELCMD <         \label{eq:adap:lbd}
%DIFDELCMD <             \\
%DIFDELCMD <             \lambda_j & \leftarrow \max(\lambda_j,0) ,
%DIFDELCMD <         \label{eq:adap:lbd:max}
%DIFDELCMD <     \end{align}%%%
\DIFdelend \DIFaddbegin \label{eq:adap}
    \begin{align}
            \ddtt{\estwth}
            &
            % =
            % -\alpha \pptfrac{L}{\estwth}
            =
            -\alpha 
            \left(
                \pptfrac{J}{\estwth}
                +
                \textstyle\sum_{i\in\mathcal{K}}
                \sigma_i 
                \estwth_i
                +
                \textstyle\sum_{j\in\mathcal{I}}
                \lambda_j 
                \pptfrac{c^{u}_{j}}{\estwth}
            \right),
        \label{eq:adap:th}
            \\
            \ddtt\sigma_i
            & 
            % = 
            % \myproj_{\sigma_i\ge0}
            % \left(
            %     \beta_i \pptfrac{L}{\sigma_i} 
            % \right)
            = 
            \myproj_{\sigma_i\ge0}
                \left(
                \beta_i^{\theta}
                c^{\theta}_{i} 
            \right)
            ,
            \quad \forall i\in\mathcal{K},
        \label{eq:adap:sig}
            \\
            \ddtt\lambda_j
            & 
            % = 
            % \myproj_{\lambda_j\ge0}
            % \left(
            %     \beta_j \pptfrac{L}{\lambda_j} 
            % \right)
            = 
            \myproj_{\lambda_j\ge0}
                \left(
                \beta_j^{u}
                c^{u}_{j} 
            \right)
            ,
            \quad \forall j\in\mathcal{I},
        \label{eq:adap:lbd}
    \end{align}\DIFaddend \DIFdelbegin %DIFDELCMD < \label{eq:adaptation law}%%%
%DIF <              %%
\DIFdelend %DIF > 
\end{subequations}
where $\alpha\in\R_{>0}$ denotes the adaptation gain (learning rate) and \DIFdelbegin \DIFdel{$\beta_j\in\R_{>0}$ denotes }\DIFdelend \DIFaddbegin \DIFadd{$\beta_i^{\theta},\beta_j^{u}\in\R_{>0}$ denote }\DIFaddend the update rate of the Lagrange multipliers\DIFdelbegin \DIFdel{in $\mathcal I$, and the arguments of $L(\cdot)$ and $J(\cdot)$ are suppressed for brevity}\DIFdelend . 
The Lagrange multipliers associated with inequality constraints are \DIFaddbegin \DIFadd{kept }\DIFaddend non-negative, \ie \DIFdelbegin \DIFdel{$\lambda_j\ge 0$}\DIFdelend \DIFaddbegin \DIFadd{$\sigma_i,\lambda_j\ge 0$}\DIFaddend , due to \DIFdelbegin \DIFdel{\eqref{eq:adap:lbd:max}.
}\DIFdelend \DIFaddbegin \DIFadd{the projection operator $\myproj_{(\cdot)\ge0}(\cdot)$ in \eqref{eq:adap:sig} and \eqref{eq:adap:lbd} \mbox{%DIFAUXCMD
\cite[Appendix E]{Krstic:1995aa}}\hskip0pt%DIFAUXCMD
.
%DIF >  To ensure Lipschitz continuity, smooth projection operators are employed, which are presented in \cite[Appendix E, eq. (E.4)]{Krstic:1995aa}.
}

\DIFaddend When a constraint \DIFdelbegin \DIFdel{$c_j$ becomes active}\DIFdelend \DIFaddbegin \DIFadd{becomes violated}\DIFaddend , \ie \DIFdelbegin \DIFdel{$c_j > 0$}\DIFdelend \DIFaddbegin \DIFadd{positive constraint violation $c^{\theta}_{i},c^{u}_{j}\ge0$}\DIFaddend , the corresponding Lagrange multiplier \DIFdelbegin \DIFdel{$\lambda_j$ }\DIFdelend increases, see \DIFaddbegin \DIFadd{\eqref{eq:adap:sig} and }\DIFaddend \eqref{eq:adap:lbd}, to address the violation \DIFdelbegin \DIFdel{, see \eqref{eq:adap:th}.
Once the violation is resolved and the constraint is no longer active, }%DIFDELCMD < \ie %%%
\DIFdel{$c_j \le 0$, }\DIFdelend \DIFaddbegin \DIFadd{by updating weight $\estwth$, as shown in \eqref{eq:adap:th}.
Ultimately, }\DIFaddend the Lagrange multiplier \DIFdelbegin \DIFdel{$\lambda_j$ decreases gradually until it returns to zero , see \eqref{eq:adap:lbd:max}.
%DIF <  \MSRY{생략 가능할 것 같습니다.}
%DIF <  \ctxt{red}
%DIF <  {
%DIF <      Note that this adaption law is similar to the ALM in \cite{Nocedal:2006aa}, where the adaptation law for Lagrange multipliers is given by $\lambda_j\leftarrow \max\left(\lambda_j+\tfrac{c_j}{\mu},0\right)$, with $\mu\in\R_{>0}$ being the penalty parameter. 
%DIF <  }
}\DIFdelend \DIFaddbegin \DIFadd{exhibits two distinct behaviors depending on the constraint status: it either decreases to zero (}\eg \DIFadd{via a projection operator) when the constraint becomes inactive, $c^{\theta}_i,c^{u}_j < 0$, or converges to a positive value when the constraint remains binding, $c^{\theta}_i,c^{u}_j = 0$, to enforce the limit.
}\DIFaddend 

\DIFdelbegin \DIFdel{At steady state, where $\ddtt{\estwth}=0$ and $\ddtt\lambda_j=0$, the KKT conditions are satisfied, }%DIFDELCMD < \ie %%%
\DIFdel{$\pptfrac{L}{\estwth}=0$, $c_j \le 0$, $\lambda_j \ge 0$, and $\lambda_j c_j=0$.
In other words, the proposed weight optimizer updates }\DIFdelend \DIFaddbegin \DIFadd{On the other hand, the adaptation law of }\DIFaddend $\estwth$ \DIFdelbegin \DIFdel{and $\lambda_j$ in a way that satisfies the KKT conditions.
These conditions represent the first-order necessary conditions for optimality, guiding the updates toward candidates for a locally optimal point.
%DIF <  The KKT condition is satisfied, when $\estwth$ converges to $\idealwth$, since $\partial L/\partial \estwth=0$ (i.e. the first-order necessary condition of the optimality is satisfied.).
}%DIFDELCMD < 

%DIFDELCMD < %%%
\subsection{\DIFdel{Approximation of the Gradient of Objective Function}}
%DIFAUXCMD
\addtocounter{subsection}{-1}%DIFAUXCMD
%DIFDELCMD < 

%DIFDELCMD < %%%
\DIFdel{The adaptation law for $\estwth$ }\DIFdelend defined in \eqref{eq:adap:th} involves the partial derivative of the \DIFdelbegin \DIFdel{filtered }\DIFdelend error with respect to control input \DIFdelbegin \DIFdel{$\pptfrac{\fe}{\rbu}$}\DIFdelend \DIFaddbegin \DIFadd{$\pptfrac{\mv{e}}{\mv{u}}$}\DIFaddend , \ie 
\DIFdelbegin \DIFdel{$
    \pptfrac{J}{\estwth}
    =
    \left(
        (\pptfrac{\fe}{\rbu})
        (\pptfrac{\rbu}{\estwth})
    \right)
    ^\top \fe
    =
    \left(
        (\pptfrac{\fe}{\rbu})
        (\pptfrac{\estNN}{\estwth})
    \right)
    ^\top \fe
$}\DIFdelend \DIFaddbegin \DIFadd{$
    \pptfrac{J}{\estwth}
    =
    \big(
        (\pptfrac{\mv{e}}{\mv{u}})
        (\pptfrac{\mv{u}}{\estwth})
    \big)
    ^\top \mv{e}
    =
    % \left(
    \big(
        (\pptfrac{\mv{e}}{\mv{u}})
        \mm{P}
        (\pptfrac{\estNN}{\estwth})
    \big)
    % \right)
    ^\top \mv{e}
$}\DIFaddend . 
Since the objective function depends on \DIFdelbegin \DIFdel{$\fe$ }\DIFdelend \DIFaddbegin \DIFadd{$\mv{e}$ }\DIFaddend of a dynamic system, obtaining \DIFdelbegin \DIFdel{$\pptfrac{\fe}{\rbu}$ }\DIFdelend \DIFaddbegin \DIFadd{$\pptfrac{\mv{e}}{\mv{u}}$ }\DIFaddend is not straightforward. 
\DIFdelbegin \DIFdel{The recommended method to calculate the exact value of $\pptfrac{\fe}{\rbu}$ is to use the forward sensitivity method \mbox{%DIFAUXCMD
\cite{Sengupta:2014aa,Ryu:2024aa} }\hskip0pt%DIFAUXCMD
by simulating the sensitivity equation as follows: 
}%DIFDELCMD < \begin{equation}
%DIFDELCMD <     \ddtt 
%DIFDELCMD <     \left(
%DIFDELCMD <         \pptfrac{\fe}{\rbu}
%DIFDELCMD <     \right)
%DIFDELCMD <     =
%DIFDELCMD <     \pptfrac{}{
%DIFDELCMD <         \rbu
%DIFDELCMD <     }
%DIFDELCMD <     \left(
%DIFDELCMD <         \rbM^{-1}(
%DIFDELCMD <             -\rbVm \fe
%DIFDELCMD <             -\mm K \fe
%DIFDELCMD <             + \mv f
%DIFDELCMD <             -\rbd + \mysat(\rbu)
%DIFDELCMD <         )
%DIFDELCMD <     \right)
%DIFDELCMD <     .
%DIFDELCMD < \end{equation}
%DIFDELCMD < %%%
\DIFdel{However, this cannot be realized, since we do not know the exact system dynamics, }%DIFDELCMD < \ie %%%
\DIFdel{$\rbM,\rbVm, \rbF$, and $\rbG$.
Additionally, the computational cost of the }\DIFdelend \DIFaddbegin \DIFadd{Exact sensitivity computation methods, such as the }\DIFaddend forward sensitivity method \DIFdelbegin \DIFdel{is high, as the number of DNN weightsis generally large.
}%DIFDELCMD < 

%DIFDELCMD < %%%
\DIFdel{In \mbox{%DIFAUXCMD
\cite{Douratsos:2007aa,Saerens:1991aa}}\hskip0pt%DIFAUXCMD
}\DIFdelend \DIFaddbegin \DIFadd{\mbox{%DIFAUXCMD
\cite{Khalil:2002aa,Ryu:2025ab}}\hskip0pt%DIFAUXCMD
, can be used in principle, but they are often computationally demanding for real-time applications with numerous DNN weights.
Therefore, several approximation methods have been proposed \mbox{%DIFAUXCMD
\cite{Douratsos:2007aa,Saerens:1991aa}}\hskip0pt%DIFAUXCMD
.
In the literature}\DIFaddend , the authors approximate \DIFdelbegin \DIFdel{$\pptfrac{\fe}{\rbu}$ }\DIFdelend \DIFaddbegin \DIFadd{$\pptfrac{\mv{e}}{\mv{u}}$ }\DIFaddend by taking the sign of each entry, \ie
\DIFdelbegin \DIFdel{$
    \pptfrac{\fe}{\rbu}
    \approx
    \left[
        \mysign
        \left(
            \pptfrac{r_i}{\tau_j}
        \right)
    \right]_{i,j\in\{1,\cdots,n\}}
$, 
assuming known controldirections}\DIFdelend \DIFaddbegin \DIFadd{$
    \pptfrac{\mv{e}}{\mv{u}}
    \approx
    \big[
        \mysign
        % \left(
        (
            \pptfrac{e_i}{u_j}
        % \right)
        )
    \big]_{i,j\in\{1,\cdots,n\}}
    ,
$
which corresponds to MIT rule in the context of adaptive control}\DIFaddend .
However, \DIFaddbegin \DIFadd{since }\DIFaddend this approach is not applicable to \DIFdelbegin \DIFdel{\eqref{eq:sys:filtered}, where control directions are unknown.
Instead, since }\DIFdelend \DIFaddbegin \DIFadd{\eqref{eq:sys}, we instead use the approximation $\pptfrac{\mv{e}}{\mv{u}}\approx \mm{I}_n$ under Assumption~\ref{assum:B:bound}, given that }\DIFaddend the sign of the control \DIFdelbegin \DIFdel{input }\DIFdelend \DIFaddbegin \DIFadd{effectiveness }\DIFaddend matrix is known\DIFdelbegin \DIFdel{—owing to the positive definiteness of $\rbM^{-1}$, see Property \ref{prop:M}—we approximate $\pptfrac{\fe}{\rbu}\approx \mm I_n$}\DIFdelend .
Consequently, the adaptation law in \eqref{eq:adap:th} becomes:
\DIFdelbegin %DIFDELCMD < \begin{equation}
%DIFDELCMD <     \ddtt{\estwth}
%DIFDELCMD <     \approx
%DIFDELCMD <     -\alpha 
%DIFDELCMD <     \left(
%DIFDELCMD <         {\pptfrac{\estNN}
%DIFDELCMD <         {\estwth}}^\top
%DIFDELCMD <         \fe
%DIFDELCMD <         +
%DIFDELCMD <         \textstyle\sum_{j\in\mathcal{I}}
%DIFDELCMD <         \lambda_j 
%DIFDELCMD <         \pptfrac{c_j}{\estwth}
%DIFDELCMD <     \right)
%DIFDELCMD <     .
%DIFDELCMD <     \label{eq:adap:th:approx}
%DIFDELCMD < \end{equation}%%%
\DIFdelend \DIFaddbegin \begin{equation}
    \ddtt{\estwth}
    \approx
    -\alpha 
    \left(
        {\pptfrac{\estNN}
        {\estwth}}^\top
        \mm{P}
        \mv{e}
        +
        \textstyle\sum_{i\in\mathcal{K}}
        \sigma_i 
        \estwth_i
        +
        \textstyle\sum_{j\in\mathcal{I}}
        \lambda_j 
        \pptfrac{c^{u}_{j}}{\estwth}
    \right)
    .
    \label{eq:adap:th:approx}
\end{equation}\DIFaddend 

\DIFdelbegin \subsection{\DIFdel{Essential and Potential Constraint Candidates}}%DIFAUXCMD
\addtocounter{subsection}{-1}%DIFAUXCMD
%DIFDELCMD < \label{sec:sub:cstr} 
%DIFDELCMD < 

%DIFDELCMD < %%%
\DIFdel{This section introduces essential constraints required to ensure the stability of the proposed CONAC, as well as conditions for potential input constraints.
}%DIFDELCMD < 

%DIFDELCMD < %%%
\DIFdel{First, weight norm constraints $c_{\theta_i}(\estwth),\ \forall i\in\{0,\cdots ,k\}$, defined in \eqref{eq:cstr:weight:ball}, are essential to prevent the weights from diverging during the adaptation process:
}%DIFDELCMD < \begin{equation}
%DIFDELCMD <     c_{\theta_i}
%DIFDELCMD <     =
%DIFDELCMD <     \tfrac{1}{2}
%DIFDELCMD <     \left(
%DIFDELCMD <         \Vert \estwth_i\Vert^2 
%DIFDELCMD <         -
%DIFDELCMD <         \overline\theta_i^2 
%DIFDELCMD <     \right)    
%DIFDELCMD <     % \le 0
%DIFDELCMD <     ,
%DIFDELCMD <     \label{eq:cstr:weight:ball}
%DIFDELCMD < \end{equation}
%DIFDELCMD < %%%
\DIFdel{with $\overline\theta_i\in\R_{>0}$ denoting the maximum allowable norm for $\estwth_i$, and the arguments of $c_{\theta_i}(\cdot)$ are omitted for brevity.
The gradient of $c_{\theta_i}$ with respect to $\estwth$ can be obtained by accumulating the gradients of each layer as follows:
}%DIFDELCMD < \begin{equation}
%DIFDELCMD <     \pptfrac{c_{\theta_i}}{\estwth_j} 
%DIFDELCMD <     =
%DIFDELCMD <     \begin{cases}
%DIFDELCMD <         \estwth_i,
%DIFDELCMD <         &
%DIFDELCMD <         \text{if } i=j,
%DIFDELCMD <         \\
%DIFDELCMD <         \mv{0},
%DIFDELCMD <         &
%DIFDELCMD <         \text{if } i\ne j
%DIFDELCMD <         .
%DIFDELCMD <     \end{cases} 
%DIFDELCMD <     \label{eq:cstr:weight:ball:grad}
%DIFDELCMD < \end{equation}
%DIFDELCMD < 

%DIFDELCMD < %%%
\DIFdel{Next, we present the following assumptions, which specify the conditions for the potential input constraints introduced in Appendix \ref{sec:appen:cstr}:
}%DIFDELCMD < 

%DIFDELCMD < \begin{assum}
%DIFDELCMD <     %%%
\DIFdel{The constraint function $c_j(\estwth)$ is convex in the (input-) $\tau$-space and satisfies $c_j(0) \le 0$ }\DIFdelend \DIFaddbegin \DIFadd{Additionally, the adaptation law \eqref{eq:adap:th} provides an optimality interpretation at equilibrium.
At an equilibrium of the adaptation law \eqref{eq:adap}, where $\ddtt{\estwth}=0$, $\ddtt\sigma_i=0$ }\DIFaddend and \DIFdelbegin \DIFdel{$c_j(\idealwth)\le 0$.
}%DIFDELCMD < \label{assum:convex}
%DIFDELCMD < \end{assum}
%DIFDELCMD < %%%
\DIFdelend \DIFaddbegin \DIFadd{$\ddtt\lambda_j=0$, the resulting equilibrium conditions correspond to KKT-like conditions, }\ie \DIFadd{$\pptfrac{L}{\estwth}\approx0$, $c^{\theta}_i,c^{u}_j \le 0$, $\sigma_i,\lambda_j \ge 0$, and $\sigma_i c^{\theta}_i=0$, and $\lambda_j c^{u}_j=0$, for all $i\in\mathcal{K}$ and $j\in\mathcal{I}$.
Thus, the equilibria characterizes candidates for locally optimal solutions of the formulated constrained adaptation problem.
}\DIFaddend 

\DIFdelbegin %DIFDELCMD < \begin{assum}
%DIFDELCMD <     %%%
\DIFdel{The selected constraints satisfy the Linear Independence Constraint Qualification (LICQ) \mbox{%DIFAUXCMD
\cite[Chap.~12, Def.~12.1]{Nocedal:2006aa}}\hskip0pt%DIFAUXCMD
.
    }%DIFDELCMD < \label{assum:LICQ}
%DIFDELCMD < \end{assum}
%DIFDELCMD < %%%
\DIFdelend %DIF >  \begin{remark}[Verification of Approximated Adaptation Law]
%DIF >      In traditional adaptive control literature, by assuming known sign 

\DIFdelbegin %DIFDELCMD < \begin{remark}
%DIFDELCMD <     %%%
\DIFdel{Assumption \ref{assum:convex} is not restrictive, since, as aforementioned, the control input saturation is practically convex.
    Moreover, Assumption \ref{assum:LICQ} is a standard assumption in optimization problems, ensuring that the gradients of the active constraints are linearly independent.
This assumption will be used in the stability analysis (see Lemma \ref{lem:convex:angle}).
}%DIFDELCMD < \end{remark}
%DIFDELCMD < %%%
\DIFdelend %DIF >      Furthermore, the approximated adaptation law \eqref{eq:adap:th:approx} coincides with the various existing adaptation laws in the literature when there are no constraints, \ie $\sigma_i=0,\ \forall i\in\mathcal{K}$ and $\lambda_j=0,\ \forall j\in\mathcal{I}$, derived using Lyapunov stability theory \cite{} and static approximation methods \cite{}.

\DIFdelbegin \DIFdel{The algorithm of the proposed CONAC is presented by Algorithm \ref{alg:CONAC}, implemented in a discrete-time controller with a sampling time of~$T_s$, where $(\cdot^{(k)})$ denotes the value at the $k$th time step.
}\DIFdelend %DIF >  \begin{remark}[Locally Lipschitz Continuity of the Adaptation Law] \label{rem:adap:lip}
%DIF >      The approximated adaptation law \eqref{eq:adap:th:approx} is locally Lipschitz continuous under the stated regularity assumptions.
%DIF >      Specifically, the error dynamics are locally Lipschitz, the DNN Jacobian $\pptfrac{\estNN}{\estwth}$ is locally bounded due to the use of $\tanh(\cdot)$ activation functions, and the input constraints $c_j^u(\mv{u})$ are $\mathcal{C}^1$ with respect to $\estwth$, while $\mv u=\mv P\estNN(\cdot;\estwth)$ is smooth in $\estwth$.
%DIF >      This property, in conjunction with the boundedness of signals established via subsequent Lyapunov analysis in Section \ref{sec:stability:analysis}, precludes the finite escape time of the closed-loop system under the proposed CONAC
%DIF >  \end{remark}

\begin{algorithm}[!t]
    \caption{Proposed CONAC}\label{alg:CONAC}
    \begin{algorithmic}[1]
        \STATE \textbf{Initialize:} \DIFdelbegin \DIFdel{DNN weights }\DIFdelend $\estwth^{(0)}$, \DIFdelbegin \DIFdel{Lagrange multipliers $\lambda_j^{(0)}$
        }%DIFDELCMD < \FOR{$k = 0 $ to $\infty$}
%DIFDELCMD <             %%%
\DIFdelend \DIFaddbegin \DIFadd{$\sigma_i^{(0)},\ \forall i\in\mathcal{K}$, $\lambda_j^{(0)},\ \forall j\in\mathcal{I}$
        }\FOR{$r = 0 $ to $\infty$}
            \DIFaddend \STATE \DIFdelbegin \DIFdel{Measure $\q^{(k)},\dq^{(k)}$
            }%DIFDELCMD < \STATE %%%
\DIFdel{Obtain $\qd^{(k)},\dqd^{(k)}$ from given $\qd(t)$
            }\DIFdelend \DIFaddbegin \DIFadd{Obtain $\mv{x}^{(r)}$ and $\mv{x}_d^{(r)}$
            }\DIFaddend \STATE \DIFdelbegin \DIFdel{$\mv{e}^{(k)}\leftarrow\q^{(k)}-\qd^{(k)}$ and $\ddtt\mv{e}^{(k)}\leftarrow\ddtt\q^{(k)}-\dqd^{(k)}$
            }%DIFDELCMD < \STATE %%%
\DIFdel{$\fe^{(k)} \leftarrow \ddtt{\mv e^{(k)}}+\mm\Lambda \mv e^{(k)}$ based on \eqref{eq:err:filtered}
}\DIFdelend \DIFaddbegin \DIFadd{$\mv{e}^{(r)}\leftarrow\mv{x}^{(r)}-\mv{x}_d^{(r)}$
}\DIFaddend 

            \STATE \textit{// Neuro-adaptive control law (Section\DIFdelbegin \DIFdel{\ref{sec:ctrl design}}\DIFdelend \DIFaddbegin \DIFadd{~\ref{sec:main:result:sub:ctrl:design}}\DIFaddend )}

                \STATE \DIFdelbegin \DIFdel{${\q}_n^{(k)}\leftarrow
                    \left(
                        {\q^{(k)}}^\top,{\dq^{(k)}}^\top,{\fe^{(k)}}^\top,1
                    \right)^\top$
                }\DIFdelend \DIFaddbegin \DIFadd{$\NNin^{(r)}\leftarrow
                    \left(
                        {\mv{x}^{(r)}}^\top,
                        {\mv{x}_d^{(r)}}^\top,
                        {\ddtt{\mv{x}}_d^{(r)}}^\top,
                        1
                    \right)^\top$
                }\DIFaddend 

                \STATE \DIFdelbegin \DIFdel{$\rbu^{(k)}\leftarrow\NN({\q}_n^{(k)};\estwth^{(k)}) $ based on \eqref{eq:control:est}
                }\DIFdelend \DIFaddbegin \DIFadd{$\mv{u}^{(r)}\leftarrow\NN(\NNin^{(r)};\estwth^{(r)}) $ based on \eqref{eq:ctrl:NN:actual}
                }\DIFaddend \STATE Apply \DIFdelbegin \DIFdel{$\rbu^{(k)}$ }\DIFdelend \DIFaddbegin \DIFadd{$\mv{u}^{(r)}$ }\DIFaddend to the system in \DIFdelbegin \DIFdel{\eqref{eq:sys1}
}\DIFdelend \DIFaddbegin \DIFadd{\eqref{eq:sys}
}\DIFaddend 

            \STATE \textit{// \DIFdelbegin \DIFdel{Train }\DIFdelend \DIFaddbegin \DIFadd{Update }\DIFaddend DNN \DIFaddbegin \DIFadd{weights }\DIFaddend (Section\DIFdelbegin \DIFdel{\ref{sec:adap_laws}}\DIFdelend \DIFaddbegin \DIFadd{~\ref{sec:main:result:sub:adap:laws}}\DIFaddend )}

            \STATE Compute \DIFdelbegin \DIFdel{$\pptfrac{\NN}{\estwth^{(k)}}$ }\DIFdelend \DIFaddbegin \DIFadd{$\pptfrac{\NN}{\estwth^{(r)}}$ }\DIFaddend as defined in \eqref{eq:gradient:phi}
            \DIFdelbegin \DIFdel{, $c_j$ and $\pptfrac{c_j}{\estwth^{(k)}}$ for all $j\in\mathcal{I}$          
            }%DIFDELCMD < 

%DIFDELCMD <             %%%
\DIFdelend \STATE \DIFdelbegin \DIFdel{$\estwth^{(k+1)}
                    \leftarrow
                    \estwth^{(k)}
                    -\alpha 
                    \left(
                        {\pptfrac{\estNN}
                        {\estwth^{(k)}}}^\top
                        \fe^{(k)}
                        +
                        \textstyle\sum_{j\in\mathcal{I}}
                        \lambda_j^{(k)} 
                        \pptfrac{c_j}{\estwth^{(k)}}
                    \right)\cdot T_s$ based on \eqref{eq:adap:th}
}\DIFdelend \DIFaddbegin \DIFadd{Compute $c^{\theta}_i,c^{u}_j$ and $\pptfrac{c^{u}_{j}}{\estwth^{(r)}}$ for all $i\in\mathcal{K}, j\in\mathcal{I}$
            }\DIFaddend 

            \STATE \DIFdelbegin \DIFdel{$\lambda_j^{(k+1)}
            \leftarrow
            \lambda_j^{(k)}
            +
            \beta_j c_j\cdot T_s$ for all $j\in\mathcal{I}$\,based\,on\,\eqref{eq:adap:lbd}}%DIFDELCMD < \STATE %%%
\DIFdel{$\lambda_j^{(k+1)} \leftarrow \max(\lambda_j^{(k+1)},0)$ for all $j\in\mathcal{I}$ from \eqref{eq:adap:lbd:max} 
}\DIFdelend \DIFaddbegin \DIFadd{Update $\estwth^{(r+1)}$, $\sigma_i^{(r+1)},\ \forall i\in\mathcal{K}$, and $\lambda_j^{(r+1)},\ \forall j\in\mathcal{I}$ using \eqref{eq:adap:th:approx} \eqref{eq:adap:sig} and \eqref{eq:adap:lbd}, respectively.
}\DIFaddend 

            %DIF >  \STATE $\estwth^{(r+1)}
            %DIF >          \leftarrow
            %DIF >          \estwth^{(r)}
            %DIF >          -\alpha 
            %DIF >          \bigg(
            %DIF >              {\pptfrac{\estNN}
            %DIF >              {\estwth^{(r)}}}^\top
            %DIF >              \fe^{(r)}
            %DIF >              +
            %DIF >              \textstyle\sum_{i\in\mathcal{K}}
            %DIF >              \sigma_i^{(r)} 
            %DIF >              \estwth_i^{(r)}
            %DIF >              +
            %DIF >              \textstyle\sum_{j\in\mathcal{I}}
            %DIF >              \lambda_j^{(r)} 
            %DIF >              \pptfrac{c^{u}_{j}}{\estwth^{(r)}}
            %DIF >          \bigg)\cdot T_s$ \eqref{eq:adap:th}
            %DIF >  \STATE $\sigma_i^{(r+1)} \leftarrow \max( \sigma_i^{(r)} + \beta_i c^{\theta}_i \cdot T_s, 0), \forall i \in \mathcal{K}$ (\eqref{eq:adap:sig}, \eqref{eq:adap:max})
            %DIF >  \STATE $\lambda_j^{(r+1)} \leftarrow \max( \lambda_j^{(r)} + \beta_j c^{u}_j \cdot T_s, 0), \forall j \in \mathcal{I}$ (\eqref{eq:adap:lbd}, \eqref{eq:adap:max})
        \ENDFOR
    \end{algorithmic}
\end{algorithm}

%DIF <   SECTION STABILITY ANALYSIS ==============================
\DIFaddbegin \DIFadd{The algorithm of the proposed CONAC is presented by Algorithm~\ref{alg:CONAC}, implemented in a discrete-time controller with a sampling time of~$T_s$, where $(\cdot^{(r)})$ denotes the value at the $r$th time step.
}

%DIF >  ------------------------------------
%DIF >        SECTION STABILITY ANALYSIS
%DIF >  ------------------------------------
\DIFaddend \section{Stability Analysis} \DIFdelbegin %DIFDELCMD < \label{sec:stability}
%DIFDELCMD < %%%
\DIFdelend \DIFaddbegin \label{sec:stability:analysis}
\DIFaddend 

\DIFdelbegin \DIFdel{Before conducting the stability analysis, let us define the weight estimation error as $\tilde\wth:= [\tilde\wth_i]_{i\in\{0,\cdots,k\}}$, where $\tilde\wth_i:=\estwth_i-\idealwth_i, \forall i\in\{0,\cdots,k\}$.
The following lemmas are introduced for }\DIFdelend \DIFaddbegin \DIFadd{In this section, the stability of the closed-loop system under the proposed CONAC is analyzed.
To facilitate }\DIFaddend the stability analysis\DIFaddbegin \DIFadd{, we introduce the following lemmas}\DIFaddend .

\DIFdelbegin %DIFDELCMD < \begin{lem}
%DIFDELCMD <     %%%
\DIFdel{If Assumptions \ref{assum:convex} and \ref{assum:LICQ} hold, then for each active constraint $c_j$, the angle between the output layer 's weight vector }\DIFdelend \DIFaddbegin \begin{lem}[Convexity of Input Constraints with respect to $\estwth_k$] \label{lem:convexity:input:constraint}
    \DIFadd{The input constraint function $c_j^{u}(\estwth_k),\ \forall j \in \mathcal{I},$ is convex with respect to the output layer weights }\DIFaddend $\estwth_k$\DIFdelbegin \DIFdel{and $\pptfrac{c_j}{\estwth_k}$ is positive; that is $\pptfrac{c_j}{\estwth_k}^\top\estwth_k>0$.
}%DIFDELCMD < \label{lem:convex:angle}
%DIFDELCMD < %%%
\DIFdelend \DIFaddbegin \DIFadd{, resulting in inequality $c_j^{u}(\estwth_k) - \pptfrac{c_j^{u}}{\estwth_k}^\top \estwth_k \le c_j^{u}(\estwth)\big|_{\estwth_k=\mv{0}}$.
}\DIFaddend \end{lem}

\begin{proof}
    \DIFdelbegin %DIFDELCMD < 

%DIFDELCMD < %%%
\DIFdel{Since $\rbu = \estNN$, a linear mapping $\mm T(\cdot):\estwth_k\to\rbu$, which is independent of $\estwth_k$, can be derived using Proposition \ref{propsit:kron}:
}%DIFDELCMD < \begin{equation}\label{eq:linear:map}
%DIFDELCMD <     \begin{aligned}
%DIFDELCMD <     \rbu 
%DIFDELCMD <     = 
%DIFDELCMD <     &
%DIFDELCMD <     \estNN 
%DIFDELCMD <     % = 
%DIFDELCMD <     % \myvec(\estNN)
%DIFDELCMD <     =
%DIFDELCMD <     \myvec(
%DIFDELCMD <         \estwV_k^\top \estact_k
%DIFDELCMD <     ) 
%DIFDELCMD <     = 
%DIFDELCMD <     (
%DIFDELCMD <         \mm I_{l_{k+1}}
%DIFDELCMD <         \otimes 
%DIFDELCMD <         \estact_k^\top
%DIFDELCMD <     )^\top
%DIFDELCMD <     \myvec(\estwV_k)
%DIFDELCMD <     \\
%DIFDELCMD <     = &
%DIFDELCMD <     \underbrace{
%DIFDELCMD <         (
%DIFDELCMD <         \mm I_{l_{k+1}}
%DIFDELCMD <         \otimes 
%DIFDELCMD <         \estact_k^\top
%DIFDELCMD <     )^\top}_{=:\mm T(\estact_k)}
%DIFDELCMD <     \estwth_k 
%DIFDELCMD <     =
%DIFDELCMD <     \mm T(\estact_k) \estwth_k.
%DIFDELCMD <     \end{aligned}
%DIFDELCMD < \end{equation}%%%
\DIFdelend %DIF >  The dependence of the preceding layers and $\NNin$ on the output layer appears through the activation vector $\act_k$.
    \DIFaddbegin \DIFadd{According to the DNN gradient in \eqref{eq:gradient:phi}, we have
    }\begin{equation} \label{eq:NN:linear}
        \estNN
        =
        \frac{
            \partial \estNN
        }{
            \partial\estwth_k
        }
        \estwth_k
        = 
        \left(
            \mm{I}_{l_{k+1}} \otimes \estact_k^\top
        \right)
        \estwth_k
        .
    \end{equation}\DIFaddend 
    \DIFaddbegin \DIFadd{This implies that $\mv{u}=\mm{P}\estNN$ depends linearly on $\estwth_k$.
    }\DIFaddend Therefore, the convexity of \DIFdelbegin \DIFdel{the input constraints in the $\rbu$-space (see Assumption \ref{assum:convex}) is preserved in }\DIFdelend \DIFaddbegin \DIFadd{$c_j^u(\estwth),\ \forall j \in \mathcal{I}$, with respect to $\mv{u}$ is preserved with respect to }\DIFaddend $\estwth_k$\DIFdelbegin \DIFdel{-space, implying
that $\pptfrac{c_j}{\estwth_k}^\top\estwth_k>0$.
    }%DIFDELCMD < 

%DIFDELCMD < %%%
\DIFdelend \DIFaddbegin \DIFadd{, due to the linear dependence of $\mv{u}$ on $\estwth_k$.
    In addition, the convexity of $c_j^u(\estwth)$ with respect to $\estwth_k$ implies the desired inequality, which concludes the proof.
    %DIF >  , where $c_j^u(\mv{0})$ denotes the value at $\estwth_k=\mv{0}$
}\DIFaddend \end{proof}

\DIFdelbegin %DIFDELCMD < \begin{lem} 
%DIFDELCMD <     %%%
\DIFdel{If input constraint $c_j(\estwth),\ \forall j\in\mathcal I \setminus \{\theta_i\}_{i\in\{0,\cdots,k\}}$, satisfies Assumption \ref{assum:convex}, then $\norm{\pptfrac{c_j}{\estwth_i}},\ \forall i\in\{k-1,\cdots, 0\}$, is bounded, provided the norms of $\estwth_i,\ \forall i\in\{k,\cdots, i+1\}$, remain bounded.
    }%DIFDELCMD < \label{lem:cstr:grad:bound}
%DIFDELCMD < %%%
\DIFdelend \DIFaddbegin \begin{lem}[Boundedness of $\Delta\mv{u}$ with respect to $\estwth_k$] \label{lem:delta:u:bound}
    \DIFadd{The saturated amount of control input $\Delta\mv{u}$ is bounded by $\estwth_k$ as $\norm{\Delta\mv{u}}\le \overline{p}\sqrt{l_k+1}\norm{\estwth_k}$.
}\DIFaddend \end{lem}

\DIFdelbegin \DIFdel{For instance, if $k=3$ and $i=1$, $\norm{\pptfrac{c_j}{\estwth_1}}$ is bounded , provided that $\norm{\estwth_3}$ and $\norm{\estwth_2}$ are bounded.
}%DIFDELCMD < 

%DIFDELCMD < %%%
\DIFdelend \begin{proof}
    \DIFdelbegin %DIFDELCMD < 

%DIFDELCMD < %%%
\DIFdel{The gradient of $c_j,\ \forall j\in\mathcal I \setminus \{\theta_i\}_{i\in\{0,\cdots,k\}}$, with respect to $\estwth_i$ is represented as
}%DIFDELCMD < \begin{equation}
%DIFDELCMD <     \pptfrac{c_j}{\estwth_i} 
%DIFDELCMD <     = 
%DIFDELCMD <     \pptfrac{c_j}{\rbu} 
%DIFDELCMD <     \pptfrac{\rbu}{\estNN} 
%DIFDELCMD <     \pptfrac{\estNN}{\estwth_i}
%DIFDELCMD <     = 
%DIFDELCMD <     \pptfrac{c_j}{\rbu} 
%DIFDELCMD <     \mm{I}_n
%DIFDELCMD <     \pptfrac{\estNN}{\estwth_i}
%DIFDELCMD <     .
%DIFDELCMD < \end{equation}%%%
\DIFdelend \DIFaddbegin \DIFadd{Since $\act_k(\cdot)$ employs the bounded activation function $\tanh(\cdot)$, we have 
    }\begin{equation}
        \norm{\act_k(\cdot)} \le \sqrt{l_k+1}
        ,
    \end{equation}\DIFaddend 
    \DIFdelbegin \DIFdel{The boundedness of the first term $\pptfrac{c_j}{\rbu}$ is guaranteed, if the input argument $\rbu$ is bounded owing to the convexity of $c_j$ by Assumption \ref{assum:convex}.
}\DIFdelend \DIFaddbegin \DIFadd{for any output of the preceding layers and $\NNin$.
    Hence, invoking the relation in \eqref{eq:NN:linear}, the control input $\mv{u}=\mm{P}\estNN$ is bounded by $\estwth_k$ as follows:
    }\begin{equation} \label{eq:NN:bound}
        \norm{\mv{u}} 
        \le
        \overline{p}
        \norm{\estNN}
        \le
        \overline{p}
        \sqrt{l_k+1}\norm{\estwth_k}
        .   
    \end{equation}

    \DIFaddend In addition, \DIFdelbegin \DIFdel{the boundedness of $\rbu$ can be ensured by the boundedness of $\estwth_k$, since $\rbu=\estNN=(\mm I_{l_{k+1}}\otimes \estact_k^\top)\estwth_k$ }\DIFdelend \DIFaddbegin \DIFadd{since $\mysat(\cdot)$ radially clips the input onto $\Omega_u$ from the origin, 
    $
        \norm{\Delta\mv u}
        =
        \norm{\mv u}-\norm{\mysat(\mv u)}
    $ holds.
    Moreover, by Assumption~\ref{assum:sat}, $\norm{\mysat(\mv u)}\ge\underline u$ for $\mv{u}\notin\Omega_u$.
    Hence,
    $
        \norm{\Delta\mv u}
        \le
        \norm{\mv u}-\underline u
    $
    when $\mv u\notin\Omega_u$, while $\Delta\mv u=\mv 0$ when $\mv u\in\Omega_u$.
    Thus, we have
    }\begin{equation}
        \norm{\Delta\mv{u}}
        \le
        \max
        \left\{
            \overline{p}\sqrt{l_k+1}\norm{\estwth_k}
            -
            \underline{u}
            ,
            0
        \right\}
        \le
        \overline{p}\sqrt{l_k+1}\norm{\estwth_k}
        ,  
    \end{equation}
    \DIFadd{which concludes the proof.    
}

\end{proof}

\begin{lem} \label{lem:proj:ineq}
  \DIFadd{For any $a,b\in\R_{\ge0}$ }\DIFaddend and \DIFdelbegin \DIFdel{$\norm{\estact_k}$ is bounded due to the properties of the activation functions.
    The third term, $\pptfrac{\estNN}{\estwth_i}$ is bounded, provided that the norms of $\estwth_i,\ \forall i \in\{k,\cdots,i+1\}$,
    are bounded.
    This }\DIFdelend \DIFaddbegin \DIFadd{$c\in\R_{>0}$, the following inequality holds:
  }\begin{equation}
    (a-b)\myproj_{a\ge0}(c) \le (a-b) c.
  \end{equation}
\end{lem}

\begin{proof}
  \DIFadd{The inequality }\DIFaddend can be verified by \DIFdelbegin \DIFdel{using the definition of $\pptfrac{\estNN}{\estwth_i}$ given in \eqref{eq:NN:grad}.
Consequently, the boundedness of $\norm{\pptfrac{c_j}{\estwth_i}},\ \forall j\in\mathcal I \setminus \{\theta_i\}_{i\in\{0,\cdots,k\}}$, depends on the boundedness of $\estwth_i,\ \forall i\in \{k,\cdots,i+1\}$.
}%DIFDELCMD < 

%DIFDELCMD < %%%
\DIFdelend \DIFaddbegin \DIFadd{considering the two cases of $a>0$ and $a=0$, as follows:
  }\begin{equation*}
    (a-b)\myproj_{a\ge0}(c) 
    =
    \begin{cases}
      (a-b) c, & \text{if } a>0,
      \\
      0, & \text{if } a=0,
    \end{cases}
  \end{equation*}
  \DIFadd{This implies that $(a-b)\myproj_{a\ge0}(c) \le (a-b) c$ holds, which concludes the proof.
}\DIFaddend \end{proof}

The following theorem \DIFdelbegin \DIFdel{states that $\fe$ and }\DIFdelend \DIFaddbegin \DIFadd{establishes the uniformly ultimately boundedness (UUB) of the tracking error $\mv{e}$, the DNN weights }\DIFaddend $\estwth$\DIFdelbegin \DIFdel{are bounded}\DIFdelend \DIFaddbegin \DIFadd{, and the Lagrange multipliers $\sigma_i,\ \forall i \in \mathcal{I}$ and $\lambda_j,\ \forall j \in \mathcal{I}$}\DIFaddend .

\begin{theorem} \DIFdelbegin \DIFdel{For the dynamical system described in \eqref{eq:sys1}, the NAC in \eqref{eq:control:est} with the weight adaptation laws in \eqref{eq:adaptation law} ensure the boundedness of the filtered error $\fe$ and }\DIFdelend \DIFaddbegin \label{thm:stability}
  \DIFadd{Consider the closed-loop system consisting of \eqref{eq:sys}, the control law \eqref{eq:ctrl:NN:actual}, and the adaptation laws \eqref{eq:adap}.
  Suppose that Assumptions~\ref{assum:B:bound}--\ref{assum:sat} hold and the initial condition belongs to the compact invariant set 
  $
  \bigcap_{i\in\mathcal K}
  \mathcal S_i(\overline V_i)
  $ 
  constructed in \eqref{eq:Sk:def} and \eqref{eq:Si:def}.
  Then, the tracking error $\mv e$, }\DIFaddend the \DIFdelbegin \DIFdel{weight estimate }\DIFdelend \DIFaddbegin \DIFadd{DNN weights }\DIFaddend $\estwth$, \DIFdelbegin \DIFdel{under the input constraints satisfying Assumption \ref{assum:convex} and \ref{assum:LICQ}, provided that the weight norm constraints \eqref{eq:cstr:weight:ball} are imposed}\DIFdelend \DIFaddbegin \DIFadd{and the Lagrange multipliers $\sigma_i,\ \forall i \in \mathcal{K},$ and $\lambda_j,\ \forall j \in \mathcal{I},$ are uniformly ultimately bounded}\DIFaddend .
\end{theorem}

\begin{proof}

%DIF <  The boundednesses of $\estwth$ and $\fe$ are analyzed recursively from the last $k$th layer to the first layer of $\estNN$. 
\DIFdelbegin \DIFdel{Invoking the fact that the amplitude of the control input depends only on }\DIFdelend \DIFaddbegin \DIFadd{The proof is divided into two steps.
We first analyze the UBB of the tracking error $\mv{e}$, the adaptive variables of output-layer, }\DIFaddend $\estwth_k$, \DIFdelbegin %DIFDELCMD < \ie %%%
\DIFdel{$\rbu=\estNN=(\mm I_{l_{k+1}}\otimes \estact_k^\top)\estwth_k$ and $\norm{\estact_k}$ is bounded, the boundedness of inner layers' weights will be established after proving the boundedness of $\fe$ and $\estwth_k$.
In addition, without loss of generality, weight norm constraints $c_{\theta_i},\ \forall i\in\{0,\cdots,k\}$, are supposed to be active.
This assumption is justified because, even if $c_{\theta_i}$ is inactive, $\estwth$ is adapted to minimize the objective function $J$ as long as the constraint is inactive.
We consider two cases: (1) the control input is saturated, }\DIFdelend and \DIFdelbegin \DIFdel{(2) the control input is not saturated}\DIFdelend \DIFaddbegin \DIFadd{$\sigma_k$, and the Lagrange multipliers associated with the input constraints $\lambda_j,\ \forall j\in\mathcal{I}$.
This step provides the key boundedness result for the tracking error and the controller output.
In the following step, we then extend the analysis to the remaining adaptive variables of inner-layer}\DIFaddend .

\DIFdelbegin \subsection*{\DIFdel{Case 1: Control input is saturated.}}
%DIFAUXCMD
\DIFdelend \DIFaddbegin \hfill
\DIFaddend 

\DIFdelbegin \DIFdel{As a result of the }\DIFdelend \DIFaddbegin \textit{\DIFadd{Step~1: Boundedness of the Tracking Error, Output-Layer Adaptive Variables, and Input-Constraint Multipliers.}}

%DIF >  For notational simplicity in this subsection, $c_j^u(\estwth_k)$ denotes the dependence of the input constraint $c_j^u(\mm P\estNN)$ on the output-layer weight $\estwth_k$.
%DIF >  The remaining dependence on $\NNin$ and the preceding-layer weights enters through the activation vector $\act_k$, whose entries are bounded due to the use of $\tanh(\cdot)$ activation functions.

\DIFadd{Consider the following Lyapunov function 
}\begin{equation} \label{eq:Vk:def}
    V_{k}(\mv{z}_k)
    :=
    V_c
    +
    \tfrac{1}{2\alpha}\estwth_k^\top\estwth_k
    +
    \tfrac{1}{\beta_k^\theta}(\sigma_k-\sigma_k^o)^2
    +
    \textstyle\sum_{j\in\mathcal{I}}
    \tfrac{1}{2\beta_j^{u}}\lambda_j^2
    ,
\end{equation}
\DIFadd{where 
$
    \mv{z}_k
    :=
    \left(
        \mv{e}^\top,
        \estwth_k^\top,
        \sigma_k,
        \left[
            \lambda_j
        \right]_{j\in\mathcal{I}}^\top
    \right)^\top
$ and $\sigma_k^o\in\R_{>0}$ will be determined later.
The Lyapunov function $V_k:=V_k(\mv{z}_k)$ is lower bounded and radially unbounded with respect to $\mv{z}_k$.
In the following analysis, we consider the case where the output-layer weight constraint is violated, }\ie \DIFadd{$c_k^{\theta}(\estwth_k)>0$.
If this condition does not hold, then $c_k^{\theta}(\estwth_k)\le0$ implies $\norm{\estwth_k}\le\overline{\theta}_k$, and the boundedness of $\estwth_k$ follows directly.
}

\DIFadd{Invoking Lemma~\ref{lem:convexity:input:constraint}, }\DIFaddend time derivative of \DIFdelbegin \DIFdel{$V_1$ in \eqref{eq:lya1:dot2}, the invariant set $\Theta_{\fe}^{1}$ of $\fe$, }%DIFDELCMD < \ie %%%
\DIFdel{if $\fe$ leaves $\Theta_{\fe}^{1}$, $\ddtt V_1$ is negative, can be obtained using Lemma \ref{lem:stable:set} }\DIFdelend \DIFaddbegin \DIFadd{$V_{k}$ can be analyzed }\DIFaddend as follows:
\DIFdelbegin %DIFDELCMD < \begin{equation}
%DIFDELCMD <     \Theta_{\fe}^{1} 
%DIFDELCMD <     = 
%DIFDELCMD <     \left\{ 
%DIFDELCMD <         \fe \in \R^n 
%DIFDELCMD <         \mid 
%DIFDELCMD <         \norm{\fe} 
%DIFDELCMD <         \le 
%DIFDELCMD <         \tfrac{
%DIFDELCMD <             \overline\tau_d
%DIFDELCMD <             +
%DIFDELCMD <             \overline\tau
%DIFDELCMD <             +
%DIFDELCMD <             \overline{\theta}_k^*\sqrt{l_{k}+1}
%DIFDELCMD <             +
%DIFDELCMD <             \overline{\epsilon}
%DIFDELCMD <         }{
%DIFDELCMD <             \lambda_{\min}(\mm K)
%DIFDELCMD <         }
%DIFDELCMD <     \right\}
%DIFDELCMD <     .
%DIFDELCMD <     \label{eq:invariant:c1:fe}
%DIFDELCMD < \end{equation}%%%
\DIFdelend \DIFaddbegin \begin{align} \label{eq:Vk:dot:1} \nonumber
  \ddtt V_{k}
  \le&
  -\underline{p}
  \norm{\mv{e}}^2
  +
  \mv{e}^\top\mm{P}
  \left(
    \estNN - \idealNN
    - \mv{\epsilon}
  \right)
  +
  \mv{e}^\top
  \Delta \mv{u}
  \\ \nonumber
  &
  -
  \underbrace{
    \estwth_k^\top
    \pptfrac{\estNN}{\estwth_k}^\top
    \mm{P}
    \mv{e}
  }_{
    = 
    \estNN^\top\mm{P}\mv{e}
    ,\ \text{see, \eqref{eq:NN:linear}}
  }
  -
  \sigma_k
  \estwth_k^\top
  \estwth_k
  -    
  \textstyle\sum_{j\in\mathcal{I}}
  \lambda_j
  \estwth_k^\top
  \pptfrac{c^{u}_j}{\estwth_k}
  \\ % \nonumber
  &
  +
  2(\sigma_k-\sigma_k^o) c_k^{\theta}(\estwth_k)
  +
  \textstyle\sum_{j\in\mathcal{I}}
  \lambda_j c^{u}_j(\estwth)
  \\ \nonumber
  \le&
  -\underline{p}
  \norm{\mv{e}}^2
  +
  \mv{e}^\top
  \mm{P}
  \left(
    - \idealNN
    - \mv{\epsilon}
  \right)
  +
  \mv{e}^\top
  \Delta \mv{u}
  -
  \sigma_k \norm{\estwth_k}^2
  \\ \nonumber
  &
  +
  (\sigma_k-\sigma_k^o) 
  \left(
    \norm{\estwth_k}^2
    -
    \overline{\theta}_k^{2}
  \right)
  +
  \underbrace{
    \textstyle\sum_{j\in\mathcal{I}}
    \lambda_j
    c^{u}_j(\estwth)\big\vert_{\estwth_k=\mv{0}}
  }_{\text{see, Lemma~\ref{lem:convexity:input:constraint}}}
  \\ \nonumber
  =&
  -\underline{p}
  \norm{\mv{e}}^2
  +
  \mv{e}^\top
  \mm{P}
  \left(
    - \idealNN
    - \mv{\epsilon}
  \right)
  +
  \mv{e}^\top
  \Delta \mv{u}
  \\ \nonumber
  &
  -
  \sigma_k^o
  \norm{\estwth_k}^2
  -
  \sigma_k
  \overline{\theta}_k^2
  +
  \sigma_k^o
  \overline{\theta}_k^2
  +
  \textstyle\sum_{j\in\mathcal{I}}
  \lambda_j
  c^{u}_j(\estwth)\big\vert_{\estwth_k=\mv{0}}
  .
\end{align}\DIFaddend 
\DIFdelbegin \DIFdel{It is notable that the ideal DNN $\idealNN_k={\idealwV_k}^\top\idealact_k$ is bounded by $\overline{\theta}_k^*\sqrt{l_{k}+1}$.
This is because, the ideal weight }\DIFdelend \DIFaddbegin \DIFadd{The projection operators in \eqref{eq:adap:sig} and \eqref{eq:adap:lbd} are omitted, according to Lemma~\ref{lem:proj:ineq}.
}

\DIFadd{For the boundedness of }\DIFaddend $\idealwth_k$ \DIFdelbegin \DIFdel{is bounded by Assumption \ref{assum:feasible} such that $\norm{\idealwth_k}\le\overline{\theta}_k^*$ for some positive constant }\DIFdelend \DIFaddbegin \DIFadd{and $\mv{\epsilon}$, we assume that $\NNin$ remains in the approximation region $\Omega_{\NNin}$, which will be ensured later in the proof.
Hence, $\idealwth_k$ and $\mv{\epsilon}$ are bounded by }\DIFaddend $\overline{\theta}_k^*$ \DIFdelbegin \DIFdel{.
Hence, invoking $\norm{\idealwV_k}_F = \norm{\idealwth_k} \le \overline{\theta}_k^*$ and $\norm{\idealact_k}\le \sqrt{l_{k}+1}$, }\DIFdelend \DIFaddbegin \DIFadd{and $\overline{\epsilon}$, respectively.
In addition, $\idealNN$ is bounded by $\overline{\theta}_k^* \sqrt{l_k+1}$, similarly to the bound of $\estNN$.
Using Lemma~\ref{lem:delta:u:bound}, we can further analyze $\ddtt V_{k}$ as follows:
}\begin{align} \label{eq:Vk:dot:2} 
  \ddtt V_{k}
  \le&
    -\underline{p}
    \norm{\mv{e}}^2
    -
    \sigma_k^o
    \norm{\estwth_k}^2
  \\ \nonumber
  &
    +
    \overline{p}
    \norm{\mv{e}}
    \left(
      \overline{\theta}^*_k
      \sqrt{l_{k}+1}
      +
      \overline{\epsilon}
      +
      \sqrt{l_k+1}\norm{\estwth_k}
    \right)        
  \\ \nonumber
  &
    -
    \sigma_k 
    \overline{\theta}_k^2
    +
    \sigma_k^o 
    \overline{\theta}_k^2
    -
    \textstyle\sum_{j\in\mathcal{I}}
    \lambda_j
    \underline{c}_j^u
    .
\end{align}

\DIFadd{For further analysis, we apply Young's inequality which states 
%DIF >  \begin{equation}
$
  ab
  \le
  \eta a^2
  +
  \frac{1}{4\eta}b^2,
  \ 
  \forall a,b\in\R, \eta\in\R_{>0}
  .
$
%DIF >  \end{equation}
By applying the Young's inequality to third term in the right-hand side of \eqref{eq:Vk:dot:2}, }\DIFaddend we have
\DIFdelbegin \DIFdel{$\norm{\idealNN}\le \overline{\theta}_k^*\sqrt{l_{k}+1}$.
}\DIFdelend \DIFaddbegin \begin{align} \label{eq:Vk:dot:3} \nonumber
  \ddtt V_{k}
  \le &
    -(\underline{p}-\eta_1-\eta_2)
    \norm{\mv{e}}^2
    -
    \left(
      \sigma_k^o
      \!-\!
      \tfrac{1}{4\eta_2}
      \overline{p}^2
      (l_k+1)
    \right)
    \norm{\estwth_k}^2
  \\ 
  &
    -
    \sigma_k \overline{\theta}_k^2
    -
    \textstyle\sum_{j\in\mathcal{I}}
    \lambda_j
    \underline{c}_j^u
  \\ \nonumber
  &
      +
  \underbrace{
    \tfrac{1}{4\eta_1}
    \left[
    \overline{p}
      \left(
        \overline{\theta}_k^*
        \sqrt{l_{k}+1}
        +
        \overline{\epsilon}
      \right)
    \right]^2
    +
    \sigma_k^o
    \overline{\theta}_k^2
  }_{=: \delta_k \in \R_{>0}}
  ,
\end{align}
\DIFadd{where $\eta_1,\eta_2\in\R_{>0}$ and $\sigma_k^o\in\R_{>0}$ are arbitrarily chosen such that $\underline{p}-\eta_1-\eta_2>0$ for given $\underline{p}$, and $\sigma_k^o - \frac{1}{4\eta_1} \overline{p}^2 (l_k+1) > 0$, respectively.
}\DIFaddend 

\DIFdelbegin \DIFdel{To investigate the boundedness of $\estwth_k$, consider the Lyapunov function candidate $V_2:=\tfrac{1}{2\alpha}\estwth_k^\top \estwth_k$.
Taking the time derivative of $V_2$ yields:
}%DIFDELCMD < \begin{equation}
%DIFDELCMD <     \begin{aligned}
%DIFDELCMD <         \ddtt V_2 
%DIFDELCMD <         =& 
%DIFDELCMD <         -\estwth_k^\top 
%DIFDELCMD <         \left(
%DIFDELCMD <             (\mm I_{l_{k+1}}\otimes \estact_k^\top)^\top
%DIFDELCMD <             \fe
%DIFDELCMD <             +
%DIFDELCMD <             \textstyle\sum_{j\in\mathcal{I}}
%DIFDELCMD <             \lambda_j 
%DIFDELCMD <             \pptfrac{c_j}{\estwth_k}
%DIFDELCMD <         \right)
%DIFDELCMD <         \\
%DIFDELCMD <         =&
%DIFDELCMD <         -\estwth_k^\top 
%DIFDELCMD <         \Big(
%DIFDELCMD <             (\mm I_{l_{k+1}}\otimes \estact_k^\top)^\top
%DIFDELCMD <             \fe
%DIFDELCMD <             +
%DIFDELCMD <             \lambda_{\theta_k}\estwth_k
%DIFDELCMD <         \\
%DIFDELCMD <         &
%DIFDELCMD <             +
%DIFDELCMD <             \textstyle\sum_{
%DIFDELCMD <                 j\in\mathcal{I}\setminus \{\theta_i\}_{i\in\{0,\cdots,k\}}
%DIFDELCMD <             }
%DIFDELCMD <             \lambda_j 
%DIFDELCMD <             \pptfrac{c_j}{\estwth_k}
%DIFDELCMD <         \Big)
%DIFDELCMD <         \\
%DIFDELCMD <         \le&
%DIFDELCMD <         -
%DIFDELCMD <         \lambda_{\theta_k}\norm{\estwth_k}^2
%DIFDELCMD <         +
%DIFDELCMD <         \underbrace{
%DIFDELCMD <             \norm{(\mm I_{l_{k+1}}\otimes {\estact}_k^\top)}
%DIFDELCMD <         }_{=:{\gamma}_1\in\R_{\ge0}}
%DIFDELCMD <         \norm{\fe}
%DIFDELCMD <         \norm{\estwth_k}
%DIFDELCMD <         \\
%DIFDELCMD <         &
%DIFDELCMD <         -
%DIFDELCMD <         \underbrace{
%DIFDELCMD <             \textstyle\sum_{
%DIFDELCMD <                 j\in\mathcal{I}\setminus \{\theta_i\}_{i\in\{0,\cdots,k\}}
%DIFDELCMD <             }
%DIFDELCMD <             \lambda_j 
%DIFDELCMD <             \estwth_k^\top \pptfrac{c_j}{\estwth_k}
%DIFDELCMD <         }_{=:{\gamma}_2\in\R_{\ge0},\ \text{by Lemma \ref{lem:convex:angle} and Assumption \ref{assum:LICQ}}}
%DIFDELCMD <         \\
%DIFDELCMD <         \le&
%DIFDELCMD <         -
%DIFDELCMD <         \lambda_{\theta_k}\norm{\estwth_k}^2
%DIFDELCMD <         +
%DIFDELCMD <         {\gamma}_1     
%DIFDELCMD <         \norm{\fe}
%DIFDELCMD <         \norm{\estwth_k}
%DIFDELCMD <         .
%DIFDELCMD <     \end{aligned}
%DIFDELCMD <     \label{eq:lya2:dot}
%DIFDELCMD < \end{equation}%%%
\DIFdelend \DIFaddbegin \DIFadd{This result implies that $\ddtt V_{k}$ is negative outside the compact set defined by 
}\begin{equation} 
    \Omega_k
    \!
    =
    \!
    \left\{
        \!
        \mv{z}_k
        \!
        :
        \!
        a
        \norm{
            \left(
                \!
                \mv{e}^\top
                \!
                ,
                \!
                \estwth_k^\top
                \!
            \right)^\top
        }^2
        \!
        +
        \!
        b\norm{
            \left(
                \sigma_k
                % \!
                ,
                % \!
                [\lambda_j]_{j\in\mathcal{I}}^\top
            \right)^\top
        }_1
        \!
        \le
        \!
        \delta_k
        \!
    \right\}
    ,
\end{equation}\DIFaddend 
\DIFdelbegin \DIFdel{According to \eqref{eq:lya2:dot} and Lemma\ref{lem:stable:set}, the invariant set of $\estwth_k$ is defined as }%DIFDELCMD < \begin{equation}
%DIFDELCMD <     \Theta_{\theta_k}^{1} 
%DIFDELCMD <     = 
%DIFDELCMD <     \left\{ 
%DIFDELCMD <         \estwth_k \in \R^{\Xi_{k}} 
%DIFDELCMD <         \mid 
%DIFDELCMD <         \norm{\estwth_k} 
%DIFDELCMD <         \le 
%DIFDELCMD <         \tfrac{
%DIFDELCMD <             {\gamma}_1
%DIFDELCMD <             (
%DIFDELCMD <                 \overline\tau_d
%DIFDELCMD <                 +
%DIFDELCMD <                 \overline\tau
%DIFDELCMD <                 +
%DIFDELCMD <                 \overline{\theta}^*_k\sqrt{l_{k}+1}
%DIFDELCMD <                 +
%DIFDELCMD <                 \overline{\epsilon}
%DIFDELCMD <             )
%DIFDELCMD <         }{
%DIFDELCMD <             \lambda_{\theta_k}\lambda_{\min}(\mm K)
%DIFDELCMD <         }
%DIFDELCMD <     \right\}
%DIFDELCMD <     .
%DIFDELCMD <     \label{eq:invariant:c1:wth}
%DIFDELCMD < \end{equation}%%%
\DIFdelend \DIFaddbegin \DIFadd{where
$
    a
    :=
    \min
    \left\{
        \underline{p}
        - \eta_1
        - \eta_2
        ,
        \sigma_k^o
        -
        \tfrac{1}{4\eta_1}
        \overline{p}^2
        (l_k+1)
    \right\}
    \in\R_{>0}
    ,
    b
    :=
    \min
    \left\{
        \overline{\theta}_k^2
        ,
        \min_{j\in\mathcal{I}} 
        \left\{
            \underline{c}_j^u
        \right\}
    \right\}
    \in\R_{>0}
    .
$
%DIF >  and
%DIF >  $
%DIF >      \mv{z}_1 
%DIF >      := 
%DIF >      \left(
%DIF >          \mv{e}^\top
%DIF >          ,
%DIF >          \estwth_k^\top
%DIF >      \right)^\top
%DIF >      ,
%DIF >      \mv{z}_2
%DIF >      :=
%DIF >      \left(
%DIF >          \sigma_k
%DIF >          ,
%DIF >          [\lambda_j]_{j\in\mathcal{I}}^\top
%DIF >      \right)^\top
%DIF >      .
%DIF >  $
Moreover, since $V_k$ is lower bounded and radially unbounded, there exists a compact sublevel set
}\begin{equation} \label{eq:Sk:def}
    \mathcal S_k(\overline V_k)
    :=
    \left\{ 
        \mv{z}_k 
        : 
        V_k(\mv{z}_k)
        \le 
        \overline{V}_k 
    \right\}
    ,
\end{equation}\DIFaddend 
\DIFdelbegin %DIFDELCMD < 

%DIFDELCMD < %%%
\DIFdel{The satisfaction of the constraints, }%DIFDELCMD < \ie %%%
\DIFdel{$c_j,\ \forall j\in\mathcal{I}\setminus \{\theta_i\}_{i\in\{0,\cdots,k-1\}}$, can be verified through the boundedness of $\estwth_k$.
For the output layer's weight norm constraint $c_{\theta_k}$, if the corresponding Lagrange multiplier $\lambda_{\theta_k}$ increases sufficiently large due to the constraint violation, $\estwth_k$ approaches toward the origin as the invariant set $\Theta_{\theta_k}^{1}$ shrinks to a point.
The remaining input constraints can be validated implicitly through the term ${\gamma}_2$ in \eqref{eq:lya2:dot}.
Similarly to $c_{\theta_k}$, when an input constraint $c_j,\ \forall j\in\mathcal{I}\setminus \{\theta_i\}_{i\in\{0,\cdots,k\}}$, is violated, the corresponding Lagrange multiplier }\DIFdelend \DIFaddbegin \DIFadd{where $\overline{V}_k\in\R_{>0}$ is arbitrarily chosen such that $\Omega_k \subset \mathcal{S}_k(\overline{V}_k)$.
Because $\ddtt V_k<0$ for all $\mv{z}_k\notin\Omega_k$, the system initialized in $\mathcal{S}_k(\overline{V}_k)$ will never escape $\mathcal{S}_k(\overline{V}_k)$ and eventually enter $\Omega_k$ in finite time.
Therefore, $\mv e$, $\estwth_k$, $\sigma_k$, and }\DIFaddend $\lambda_j$\DIFdelbegin \DIFdel{increases sufficiently to a large value, leading to an increase in ${\gamma}_2$.
It makes ${\gamma}_2$ dominate the right-hand side of \eqref{eq:lya2:dot}, which makes $\ddtt V_2$ negative definite.
This drives $\estwth_k$ toward the origin until the constraint is satisfied.
%DIF <  By Assumption \ref{assum:convex} and Lemma \ref{lem:convex:angle}, the input constraints can be considered as convex constraint in the $\estwth_k$-space.
}\DIFdelend \DIFaddbegin \DIFadd{, $\forall j\in\mathcal I$, are uniformly ultimately bounded in $\mathcal{S}_k(\overline{V}_k)$.
}\DIFaddend 

\DIFdelbegin \DIFdel{Therefore, }\DIFdelend \DIFaddbegin \DIFadd{In addition, the boundedness of $\mv{z}_k$ implies that $\mv{x}$ is bounded, since $\mv{e}=\mv{x}-\mv{x}_d$ and $\mv{x}_d$ is bounded by Assumption~\ref{assum:xd}.
Then, by invoking }\DIFaddend the boundedness of \DIFdelbegin \DIFdel{$\fe$ and $\estwth_k$ is guaranteed by the invariant sets $\Theta_{\fe}^{1}$ }\DIFdelend \DIFaddbegin \DIFadd{$\mv{x}$, $\mv{x}_d$, and $\ddtt{\mv{x}}_d$, we can ensure that $\NNin$ remains in the approximation region $\Omega_{\NNin}$, which justifies the boundedness of $\idealwth_k$ }\DIFaddend and \DIFdelbegin \DIFdel{$\Theta_{\theta_k}^{1}$, and the satisfaction of $c_j,\ \forall j\in\mathcal{I}\setminus \{\theta_i\}_{i\in\{0,\cdots,k-1\}}$, is ensured}\DIFdelend \DIFaddbegin \DIFadd{$\mv{\epsilon}$ used in the above analysis}\DIFaddend .

\DIFdelbegin \subsection*{\DIFdel{Case 2: Control input is not saturated.}}
%DIFAUXCMD
\DIFdelend \DIFaddbegin \hfill
\DIFaddend 

\DIFdelbegin \DIFdel{Since the input constraints are inactive, the saturation function $\mysat(\cdot)$ in \eqref{eq:lya1:dot2} and the input constraint functions $c_j,\ \forall j\in\mathcal I \setminus \{\theta_i\}_{i\in\{0,\cdots,k\}}$, in \eqref{eq:lya2:dot} can be neglected. 
Consider the Lyapunov function candidate $V_3:=V_1+V_2$, whose time derivative is given by:
}%DIFDELCMD < \begin{equation}
%DIFDELCMD <     \begin{aligned}
%DIFDELCMD <         \ddtt V_3
%DIFDELCMD <         =&
%DIFDELCMD <         -\lambda_{\min}(\mm K)\norm{\fe}^2
%DIFDELCMD <         +\overline\tau_d\norm{\fe}
%DIFDELCMD <         +\fe^\top (\estNN - \idealNN - \mv\epsilon)
%DIFDELCMD <         \\
%DIFDELCMD <         &
%DIFDELCMD <         -\estwth_k^\top 
%DIFDELCMD <         \left(
%DIFDELCMD <             (\mm I_{l_{k+1}}\otimes \estact_k^\top)^\top
%DIFDELCMD <             \fe
%DIFDELCMD <             +
%DIFDELCMD <             \lambda_{\theta_k}\estwth_k
%DIFDELCMD <         \right)
%DIFDELCMD <         \\
%DIFDELCMD <         \le&
%DIFDELCMD <         -\lambda_{\min}(\mm K)\norm{\fe}^2
%DIFDELCMD <         +\fe^\top \estNN
%DIFDELCMD <         +
%DIFDELCMD <         \left(
%DIFDELCMD <             \overline\tau_d+\norm{
%DIFDELCMD <                 \idealNN+\mv\epsilon
%DIFDELCMD <             }
%DIFDELCMD <         \right)
%DIFDELCMD <         \norm{\fe}
%DIFDELCMD <         \\
%DIFDELCMD <         &
%DIFDELCMD <         -\estNN^\top\fe 
%DIFDELCMD <         -\lambda_{\theta_k}\norm{\estwth_k}^2
%DIFDELCMD <         \\
%DIFDELCMD <         \le&
%DIFDELCMD <         -
%DIFDELCMD <         \lambda_{\min}(\mm K)
%DIFDELCMD <         \norm{\fe}^2
%DIFDELCMD <         +
%DIFDELCMD <         \left(
%DIFDELCMD <             \overline\tau_d
%DIFDELCMD <             +
%DIFDELCMD <             \overline{\theta}^*_k\sqrt{l_{k}+1}
%DIFDELCMD <             +
%DIFDELCMD <             \overline{\epsilon}
%DIFDELCMD <         \right)
%DIFDELCMD <         \norm{\fe}
%DIFDELCMD <         \\
%DIFDELCMD <         &
%DIFDELCMD <         -
%DIFDELCMD <         \lambda_{\theta_k}\norm{\estwth_k}^2
%DIFDELCMD <         .
%DIFDELCMD <     \end{aligned}
%DIFDELCMD <     \label{eq:lya3:dot}
%DIFDELCMD < \end{equation}
%DIFDELCMD < %%%
\DIFdel{Based on the same reasoning as in Case 1, }\DIFdelend \DIFaddbegin \textit{\DIFadd{Step~2: Boundedness of the Inner-Layer Adaptive Variables.}}

\DIFadd{The remaining weights $[\estwth_i]_{i\in\{0,\cdots,k-1\}}$ and the Lagrange multipliers $[\sigma_i]_{i\in\{0,\cdots,k-1\}}$ are analyzed recursively from $(k-1)$th layer to the input layer ($i=0$). 
Specifically, for an arbitrary layer index $i\in\{0,\cdots,k-1\}$, we assume that the subsequent adaptive variables, $\estwth_{i+1},\sigma_{i+1},\cdots,\estwth_k,\sigma_k,$ have already been shown to be bounded, and then prove }\DIFaddend the boundedness of \DIFdelbegin \DIFdel{$\fe$ is ensured by the invariant set:
}%DIFDELCMD < \begin{equation}
%DIFDELCMD <     \Theta_{\fe}^{2}
%DIFDELCMD <     =
%DIFDELCMD <     \left\{ 
%DIFDELCMD <         \fe \in \R^n 
%DIFDELCMD <         \mid 
%DIFDELCMD <         \norm{\fe} 
%DIFDELCMD <         \le 
%DIFDELCMD <         \tfrac{
%DIFDELCMD <             \overline\tau_d
%DIFDELCMD <             +
%DIFDELCMD <             \overline{\theta}^*_k\sqrt{l_{k}+1}
%DIFDELCMD <             +
%DIFDELCMD <             \overline{\epsilon}
%DIFDELCMD <         }{
%DIFDELCMD <             \lambda_{\min}(\mm K)
%DIFDELCMD <         }
%DIFDELCMD <     \right\}
%DIFDELCMD <     .
%DIFDELCMD < \end{equation}%%%
\DIFdelend \DIFaddbegin \DIFadd{$\estwth_i$ and $\sigma_i$.
}

\DIFadd{Consider the Lyapunov function
}\begin{equation} \label{eq:Vi:def}
    V_{i}(\mv{z}_i)
    :=
    \tfrac{1}{2\alpha}\estwth_i^\top\estwth_i
    +
    \tfrac{1}{\beta_i^\theta}(\sigma_i-\sigma_i^o)^2
    ,
\end{equation}\DIFaddend 
\DIFdelbegin \DIFdel{The boundedness of }\DIFdelend \DIFaddbegin \DIFadd{where 
$
    \mv{z}_i:=
    \left(
        \estwth_i^\top,
        \sigma_i
    \right)^\top
    ,
$
and $\sigma_i^o\in\R_{>0}$ will be determined in the subsequent analysis.
Similar to $V_k$, the weight constraint $c_i^{\theta}(\estwth_i)$ is assumed to be violated, and the projection operator in \eqref{eq:adap:sig} is omitted according to Lemma~\ref{lem:proj:ineq}.
Taking time derivative of $V_i$ yields
}\begin{align} \label{eq:Vi:dot:1} 
    \ddtt V_i
    \le&
        -\sigma_i \norm{\estwth_i}^2
        -
        \estwth_i^\top
        {\pptfrac{\estNN}{\estwth_i}}^\top
        \mm{P}
        \mv{e}
    \\ \nonumber
    &
        -
        \textstyle\sum_{j\in\mathcal{I}}
        \lambda_j 
        \estwth_i^\top
        {\pptfrac{\estNN}{\estwth_i}}^\top
        \mm{P}
        \pptfrac{c^{u}_{j}}{\mv{u}}
        +
        2(\sigma_i-\sigma_i^o) c_i^{\theta}(\estwth_i)
    \\ \nonumber
    \le&
        -\sigma_i \norm{\estwth_i}^2
        +
        (\sigma_i-\sigma_i^o) 
        \left(
            \norm{\estwth_i}^2
            -
            \overline{\theta}_i^2
        \right)
    \\ \nonumber
    &
        -
        \estwth_i^\top
        \underbrace{
        \left(
            {\pptfrac{\estNN}{\estwth_i}}^\top
            \mm{P}
            \left(
                \mv{e}
                +
                \textstyle\sum_{j\in\mathcal{I}}
                \lambda_j 
                \pptfrac{c^{u}_{j}}{\mv{u}}
            \right)
        \right)
        }_{=: \gamma_i}
    \\ \nonumber
    \le&
        -\sigma_i^o \norm{\estwth_i}^2
        -\sigma_i \overline{\theta}_i^2
        +
        \norm{\estwth_i}
        \norm{\gamma_i}
        +
        \sigma_i^o \overline{\theta}_i^2
    .
\end{align}

\DIFadd{For further analysis, the boundedness of $\gamma_i$ is analyzed.
According to the DNN gradient in \eqref{eq:gradient:phi}, $\pptfrac{\estNN}{\estwth_i}$ is bounded, since $\estwth_{i+1},\cdots,\estwth_k$ are bounded by assumption.
The remaining terms $\mv e$, }\DIFaddend $\estwth_k$\DIFdelbegin \DIFdel{can also be ensured by the invariant set $\Theta_{\theta_k}^{2}=\{ \estwth_k \in \R^{\Xi_{k}} \mid \norm{\estwth_k} \le \overline\theta_k \}$, as $\lambda_{\theta_k}$ remains positive unless the weight norm constraint $c_{\theta_k}$ is satisfied.
}\DIFdelend \DIFaddbegin \DIFadd{, and $[\lambda_j]_{j\in\mathcal I}$ are bounded by the result of Step~1.
Furthermore, $\mv u$ is bounded because $\mv u$ is bounded by $\estwth_k$ in \eqref{eq:NN:bound}.
In addition, since $c_j^u(\cdot)$ is smooth by Assumption~\ref{assum:sat} and $\mv u$ is bounded, $\pptfrac{c_j^u}{\mv u}$ is bounded.
Therefore, $\gamma_i$ is bounded by a unknown positive constant $\overline{\gamma}_i\in\R_{>0}$.
}\DIFaddend 

\DIFdelbegin %DIFDELCMD < \hfill 
%DIFDELCMD < 

%DIFDELCMD < %%%
\DIFdel{The boundedness of the inner layer weights and the satisfaction of their corresponding constraints $c_j,\ \forall j\in\{\theta_i\}_{i\in\{k-1,\cdots,0\}}$, can be established recursively using \mbox{%DIFAUXCMD
\cite[Chap.~4, Thm.~1.9]{Desoer:2009aa}}\hskip0pt%DIFAUXCMD
.
The dynamics of $\estwth_i,\ \forall i\in\{k-1,\cdots,0\}$, are represented as
}%DIFDELCMD < \begin{equation}
%DIFDELCMD <     \ddtt \estwth_i 
%DIFDELCMD <     =
%DIFDELCMD <     -\alpha
%DIFDELCMD <     \left(
%DIFDELCMD <         {\pptfrac{\estNN}{\estwth_i}}^\top
%DIFDELCMD <         \fe
%DIFDELCMD <         +
%DIFDELCMD <         \lambda_{\theta_i}\estwth_i
%DIFDELCMD <         +
%DIFDELCMD <         \textstyle\sum_{
%DIFDELCMD <             j\in\mathcal{I}\setminus \{\theta_i\}_{i\in\{0,\cdots,k\}}
%DIFDELCMD <         }
%DIFDELCMD <         \lambda_j
%DIFDELCMD <         \pptfrac{c_j}{\estwth_i}
%DIFDELCMD <     \right)
%DIFDELCMD <     .
%DIFDELCMD <     \label{eq:estwth:dyn}
%DIFDELCMD < \end{equation}%%%
\DIFdelend \DIFaddbegin \DIFadd{Then, by applying the Young's inequality to the coupled term $\norm{\estwth_i}\norm{\gamma_i}$ in \eqref{eq:Vi:dot:1}, we have
}\begin{align} \label{eq:Vi:dot:2}
  \ddtt V_i
  \le&
    -
    \left(
      \sigma_i^o
      -
      1
    \right)
    \norm{\estwth_i}^2
    -
    \sigma_i \overline{\theta}_i^2
    +
    \delta_i
  ,
\end{align}\DIFaddend 
\DIFdelbegin \DIFdel{By Lemma \ref{lem:cstr:grad:bound}, $\estwth_i$ is bounded , provided that $\estwth_i,\ \forall i\in\{k,\cdots,i+1\}$, are bounded.
This holds because the system matrix $-\lambda_{\theta_i} \mm I_{\Xi_i}$ is stable, and the residual terms—such as $\pptfrac{\estNN}{\estwth_i}, \fe$, and $\pptfrac{c_j}{\estwth_i}$ —are bounded}\DIFdelend \DIFaddbegin \DIFadd{where $\delta_i:=\frac{1}{4}\overline{\gamma}_i^2 + \sigma_i^o \overline{\theta}_i^2$ is a positive constant, and $\sigma_i^o\in\R_{>0}$ is arbitrarily chosen such that $\sigma_i^o - 1 > 0$.
Similarly to Step~1, $\ddtt V_i$ is negative outside the compact set defined by 
}\begin{equation}
  \Omega_i
  =
  \left\{
    \mv{z}_i
    :
    \left(
      \sigma_i^o
      -
      \tfrac{\overline{\gamma}_i^2}{2}
    \right)
    \norm{\estwth_i}^2
    +
    \overline{\theta}_i^2
    \norm{\sigma_i}_1
    \le
    \delta_i
  \right\}
  , 
\end{equation}
\DIFadd{which implies the boundedness of $\estwth_i$ and $\sigma_i$}\DIFaddend .
Moreover, \DIFdelbegin \DIFdel{the Lagrange multiplier $\lambda_j$ does not diverge since the input constraints $c_j,\ \forall j\in\mathcal{I}\setminus \{\theta_i\}_{i\in\{0,\cdots,k\}}$ are satisfied before divergence can occur.
Therefore, starting from the $(k-1)$th layer}\DIFdelend \DIFaddbegin \DIFadd{a compact sublevel set $\mathcal{S}_i(\overline{V}_i)$ of $V_i$ can be constructed such that 
}\begin{equation} \label{eq:Si:def}
  S_i(\overline{V}_i)
  :=
  \left\{
    \mv{z}_i
    :
    V_i(\mv{z}_i)
    \le
    \overline{V}_i
  \right\}
  ,
\end{equation}
\DIFadd{where $\overline{V}_i\in\R_{>0}$ is arbitrarily chosen such that $\Omega_i \subset \mathcal{S}_i(\overline{V}_i)$.
Since $\ddtt V_i<0$ for all $\mv{z}_i\notin\Omega_i$, the system initialized in $\mathcal{S}_i(\overline{V}_i)$ will remain in $\mathcal{S}_i(\overline{V}_i)$.
}

\DIFadd{By recursively applying the above analysis from $i=k-1$ to $i=0$}\DIFaddend , the boundedness of $\estwth_i$ \DIFaddbegin \DIFadd{and $\sigma_i$ for all $i\in\{0,1,\cdots,k-1\}$ }\DIFaddend can be established\DIFdelbegin \DIFdel{recursively down to the input layer ($i=0$).
Moreover the satisfaction of the weight norm constraints for the inner layers can be verified, as the Lagrange multipliers $\lambda_{\theta_i}$, $\forall i\in\{k-1,\cdots,0\}$, increase sufficiently large to stabilize the $\estwth_i$ dynamics in \eqref{eq:estwth:dyn}, leading to its convergence toward the origin.
  }\DIFdelend \DIFaddbegin \DIFadd{, which concludes the proof.
}\DIFaddend 

\DIFdelbegin %DIFDELCMD < \hfill
%DIFDELCMD < 

%DIFDELCMD < %%%
\DIFdel{In conclusion, the filtered tracking error $\fe$ and the weight estimate $\estwth$ are bounded, }\DIFdelend \DIFaddbegin \end{proof}


%DIF >  Consider the following Lyapunov function $V_{k}:=V_c+\tfrac{1}{2\alpha}\estwth_k^\top\estwth_k$.
%DIF >  In the following analysis, we consider the case where the output-layer weight constraint is violated, \ie $c_k^{\theta}(\estwth_k)>0$.
%DIF >  If this condition does not hold, then $c_k^{\theta}(\estwth_k)\le0$ implies $\norm{\estwth_k}\le\overline{\theta}_k$, and the boundedness of $\estwth_k$ follows directly.
%DIF >  Taking time derivative of $V_{k}$, we have
%DIF >  \begin{align} \label{eq:Vk:dot:1} 
%DIF >      \ddtt V_{k}
%DIF >      \le&
%DIF >      -\underline{p}
%DIF >      \norm{\mv{e}}^2
%DIF >      +
%DIF >      \mv{e}^\top\mm{P}
%DIF >      \left(
%DIF >          \estNN - \idealNN
%DIF >          - \mv{\epsilon}
%DIF >      \right)
%DIF >      +
%DIF >      \mv{e}^\top
%DIF >      \Delta \mv{u}
%DIF >      \\ \nonumber
%DIF >      &
%DIF >      -
%DIF >      \underbrace{
%DIF >          \estwth_k^\top
%DIF >          \pptfrac{\estNN}{\estwth_k}^\top
%DIF >          \mm{P}
%DIF >          \mv{e}
%DIF >      }_{
%DIF >          = \estNN^\top\mm{P}\mv{e}
%DIF >      }
%DIF >      -
%DIF >      \sigma_k
%DIF >      \estwth_k^\top
%DIF >      \estwth_k
%DIF >      -
%DIF >      \underbrace{
%DIF >          \textstyle\sum_{j\in\mathcal{A}}
%DIF >          \lambda_j
%DIF >          \estwth_k^\top
%DIF >          \pptfrac{c^{u}_j}{\estwth_k}
%DIF >      }_{\ge 0, \text{convexity of $c^{u}_j$ w.r.t. $\estwth_k$}}
%DIF >      \\ \nonumber
%DIF >      \le&
%DIF >      -\underline{p}
%DIF >      \norm{\mv{e}}^2
%DIF >      +
%DIF >      \mv{e}^\top
%DIF >      \mm{P}
%DIF >      \left(
%DIF >          - \idealNN
%DIF >          - \mv{\epsilon}
%DIF >      \right)
%DIF >      +
%DIF >      \mv{e}^\top
%DIF >      \Delta \mv{u}
%DIF >      -
%DIF >      \sigma_k \norm{\estwth_k}^2
%DIF >      \\ \label{eq:Vk:dot:1:1}
%DIF >      \le&
%DIF >      -
%DIF >      \underline{p}
%DIF >      \norm{\mv{e}}^2
%DIF >      -
%DIF >      \sigma_k \norm{\estwth_k}^2
%DIF >      \\ \nonumber
%DIF >      &
%DIF >      +
%DIF >      \norm{\mv{e}}
%DIF >      \left[
%DIF >          \overline{p}
%DIF >          \left(
%DIF >              \overline{\theta}^*_k
%DIF >              \sqrt{l_{k}+1}
%DIF >              +
%DIF >              \overline{\epsilon}
%DIF >          \right)
%DIF >          +
%DIF >          \norm{\Delta \mv{u}}
%DIF >      \right]
%DIF >      \\ \nonumber
%DIF >      \le&
%DIF >  -\underline{p}
%DIF >      \norm{\mv{e}}^2
%DIF >      -\sigma_k \norm{\estwth_k}^2
%DIF >      \\ \nonumber
%DIF >      &
%DIF >      +
%DIF >      \norm{\mv{e}}
%DIF >      \left[
%DIF >          \overline{p}
%DIF >          \left(
%DIF >              \overline{\theta}_k^*
%DIF >              \sqrt{l_{k}+1}
%DIF >              +
%DIF >              \overline{\epsilon}
%DIF >          \right)
%DIF >          +
%DIF >          \overline{p}
%DIF >          \sqrt{l_k+1}\norm{\estwth_k}
%DIF >          -
%DIF >          \underline{u}
%DIF >      \right]
%DIF >      \\ \label{eq:Vk:dot:1:2}
%DIF >      \le&
%DIF >      -\left(\underline{p}-1\right)
%DIF >      \norm{\mv{e}}^2
%DIF >      -\left(\sigma_k - \tfrac{\overline{p}^2(l_k+1)}{2} \right)
%DIF >      \norm{\estwth_k}^2
%DIF >      \\ \nonumber
%DIF >      &
%DIF >      +
%DIF >      \left[
%DIF >      \overline{p}
%DIF >          \left(
%DIF >              \overline{\theta}_k^*
%DIF >              \sqrt{l_{k}+1}
%DIF >              +
%DIF >              \overline{\epsilon}
%DIF >          \right)
%DIF >          -
%DIF >          \underline{u}
%DIF >      \right]^2
%DIF >      ,
%DIF >  \end{align}
%DIF >  where $\mathcal{A}:=\{j\in\mathcal{I}: c_j^u(\estwth) \ge 0\}$ denotes the set of active input constraints.
%DIF >  \cred{
%DIF >  From \eqref{eq:Vk:dot:1:1}, it follows that the tracking error $\mv{e}$ ultimately converges to a compact set 
%DIF >  $
%DIF >      \Omega^{e}_1 
%DIF >      := 
%DIF >      \left\{ 
%DIF >          \mv{e} 
%DIF >          : 
%DIF >          \norm{\mv{e}} \le \tfrac{\overline{p}}{\underline{p}}
%DIF >          \left(
%DIF >              \overline{\theta}_k^*
%DIF >              \sqrt{l_{k}+1}
%DIF >              +
%DIF >              \overline{\epsilon}
%DIF >              +
%DIF >              \norm{\Delta \mv{u}}
%DIF >          \right)
%DIF >      \right\}
%DIF >      .
%DIF >  $
%DIF >  This set $\Omega_{1}^{e}$ implies that, if input saturation does not occur, $\Delta\mv{u}=\mv{0}$, the tracking error $\mv{e}$ is ultimately bounded within a compact set with a radius of $\tfrac{\overline{p}}{\underline{p}}\left(\overline{\theta}_k^* \sqrt{l_{k}+1} + \overline{\epsilon}\right)$, otherwise, the radius is proportionally enlarged by $\norm{\Delta \mv{u}}$.
%DIF >  However, the UUB property of $\estwth_k$ cannot be guaranteed by \eqref{eq:Vk:dot:1:1}, due to $\norm{\Delta\mv{u}} \le \max\{0, \norm{\mv{u}}-\underline{u}\} \le \max\{0, \overline{p}\sqrt{l_k+1}\norm{\estwth_k}-\underline{u}\}$.
%DIF >  }

%DIF >  According to \eqref{eq:Vk:dot:1:2}, the UBB of $\mv{e}$ and $\estwth_k$ is guaranteed only if $\underline{p} > 1$ and $\sigma_k > \tfrac{\overline{p}^2(l_k+1)}{2}$, respectively.
%DIF >  In other words, when input saturation occurs, $\sigma_k$ should increase to a value larger than $\tfrac{\overline{p}^2(l_k+1)}{2}$ triggered by the violation of weight constraint $c_{k}^{\theta}$, to guarantee the UBB of $\mv{e}$ and $\estwth_k$.
%DIF >  The finite reaching time of $\sigma_k$ to $\tfrac{\overline{p}^2(l_k+1)}{2},$ along with the boundedness of $\norm{\mv{e}}$ and $\norm{\estwth_k}$, can be ensured using the Gronwall-Bellman inequality, see \cite[pp. 651-652, Lemma A.1]{Khalil:2002aa}, owing to locally Lipschitz continuity of the system dynamics \eqref{eq:sys} and the adaptation law \eqref{eq:adap}.

%DIF >  \hfill

%DIF >  \cred{
%DIF >  In conclusion, the tracking error $\mv{e}$ is ultimately bounded in $\Omega_1^{e}$, which is proportionally enlarged by $\norm{\Delta\mv{u}}$.
%DIF >  The weight $\estwth_k$ is also bounded by weight constraint in absence of input saturation.
%DIF >  When input saturation occurs, $\estwth_k$ may exceed this bound for a finite time, but it subsequently converges to under the weight constraint $\overline{\theta}_k$, after $\sigma_k$ reaches $\tfrac{\overline{p}^2(l_k+1)}{2}$.
%DIF >  }

%DIF >  % \begin{figure}[!t]
%DIF >  %     \centering
%DIF >  %     \includegraphics[width=0.99\linewidth]{fig/stability_set.png}
%DIF >  %     \caption{
%DIF >  %         Illustration of example of the variable ultimate bound of the tracking error $\mv{e}$ and $\estwth_k$.
%DIF >  %         At $t=t_0$, input saturation occurs, leading to larger ultimate bound of $\mv{e}$, $\Omega_2^{e}$.
%DIF >  %         From $t_1$ to $t_2$, $\sigma_k$ increases to $\tfrac{\overline{p}^2(l_k+1)}{2}$, and thus $\estwth_k$ violates the weight constraint, $c_k^{\theta}(\estwth_k)$.
%DIF >  %         After $t_2$, $\sigma_k$ is larger than $\tfrac{\overline{p}^2(l_k+1)}{2}$, thus $\norm{\estwth_k}$ converges to $\Omega_1^{\theta}$.
%DIF >  %         Finally, the ultimate bound of $\mv{e}$ converges to $\Omega_1^{e}$, when input saturation is resolved at $t_4$.
%DIF >  %     }
%DIF >  %     \label{fig:stability:set}
%DIF >  % \end{figure}

%DIF >  \subsection{Stability of Lagrange Multipliers and Remaining Inner Layers' Weights} \label{sec:stability:analysis:sub:step2}

%DIF >  % \begin{lem}
%DIF >  % Consider the projected multiplier dynamics
%DIF >  % \[
%DIF >  % \dot\lambda_j =
%DIF >  % \begin{cases}
%DIF >  % \beta_j c_j^u(u), & \lambda_j>0 \ \text{or}\ c_j^u(u)>0,\\
%DIF >  % 0, & \lambda_j=0,\ c_j^u(u)\le0,
%DIF >  % \end{cases}
%DIF >  % \]
%DIF >  % where \(\beta_j>0\).
%DIF >  % Suppose that all closed-loop signals except \(\lambda_j\) are bounded and that \(c_j^u(u(t))\) is bounded from above.
%DIF >  % Furthermore, suppose that there exist constants \(d_j>0\) and \(\rho_j>0\) such that, whenever \(c_j^u(u)>0\),
%DIF >  % \[
%DIF >  % \dot c_j^u(u) \le d_j-\rho_j\lambda_j.
%DIF >  % \]
%DIF >  % Then \(\lambda_j(t)\) is bounded for all \(t\ge0\).
%DIF >  % \end{lem}

%DIF >  % First, consider active constraint vector $\mv{c} := \left\{ c_j^\theta :  \right\}



%DIF >  \hfill


%DIF >  However, the above UBB does not guarantee the boundedness of Lagrange multipliers, $\sigma_k$ and $[\lambda_j]_{j\in\mathcal{I}}$.
%DIF >  Consider Lyapunov function $V_3=\tfrac{1}{2\alpha}\estwth_k^\top\estwth_k + \tfrac{1}{2\beta_k}\sigma_k^2 + \sum_{j\in\mathcal{I}}\tfrac{1}{2\beta_j}\lambda_j^2$.
%DIF >  The time derivative of $V_3$ is given as
%DIF >  \begin{align} \label{eq:V3:dot}
%DIF >      \ddtt V_3
%DIF >      =&
%DIF >      -
%DIF >      \underbrace{
%DIF >          \estwth_k^\top
%DIF >          \pptfrac{\estNN}{\estwth_k}^\top
%DIF >          \mm{P}
%DIF >          \mv{e}
%DIF >      }_{
%DIF >          =:b, \text{ UBB of $\mv{z}$}
%DIF >      }
%DIF >      -
%DIF >      \sigma_k
%DIF >      \left(
%DIF >          \estwth_k^\top
%DIF >          \pptfrac{c^{\theta}_k}{\estwth_k}
%DIF >          -
%DIF >          c^{\theta}_k(\estwth_k)
%DIF >      \right)
%DIF >      \\ \nonumber
%DIF >      &
%DIF >      -
%DIF >      \textstyle\sum_{
%DIF >          j\in\mathcal{I}
%DIF >      }
%DIF >      \lambda_j
%DIF >      \left(
%DIF >          \estwth_k^\top
%DIF >          \pptfrac{c^{u}_j}{\estwth_k}
%DIF >          -
%DIF >          c^{u}_j(\estwth_k)
%DIF >      \right)
%DIF >      \\ \nonumber
%DIF >      \le&
%DIF >      -
%DIF >      \sigma_k
%DIF >      \overline{\theta}_k^2
%DIF >      +
%DIF >      \textstyle\sum_{
%DIF >          j\in\mathcal{I}
%DIF >      }
%DIF >      \lambda_j c^{u}_j(\mv{0})
%DIF >      +
%DIF >      b
%DIF >      .
%DIF >  \end{align}
%DIF >  The convexity of $c^{\theta}_k(\estwth_k)$ and $c^{u}_j(\estwth_k),\ \forall j\in\mathcal{I}$ with respect to $\estwth_k$ is utilized, \ie 
%DIF >  $
%DIF >      f(\mv{x})-f(\mv{y})
%DIF >      \le
%DIF >      \pptfrac{f}{\mv{x}}^\top(\mv{x}-\mv{y})
%DIF >      ,
%DIF >      \ 
%DIF >      \forall \mv{x},\mv{y}\in\R^m
%DIF >  $ for a convex function $f:\R^m\to\R$.
%DIF >  From \eqref{eq:V3:dot}, if any constraint is violated leading positive Lagrange multipliers, \ie $\sigma_k > 0$ or $\lambda_j > 0$, then its ultimate boundedness is guaranteed, such that $\sigma_k > \tfrac{b}{c^{\theta}_k(\mv{0})}$ or $\lambda_j > \tfrac{b}{c^{u}_j(\mv{0})}$, respectively, since constraints are satisfied at the origin, \ie $c^{\theta}_k(\mv{0})=-\overline{\theta}_k^2 < 0$ and $c^{u}_j(\mv{0}) < 0$. 

%DIF >  \hfill

%DIF >  The boundedness of the remaining inner layers' weights, $\estwth_{i},\ \forall i\in\{k-1,\cdots,0\}$, and Lagrange multipliers $\sigma_i,\forall i\in\{k-1,\cdots,0\},$ can be guaranteed recursively from $k-1$ to $0$ layers, by analyzing the time derivative of the Lyapunov function $V_{3+i} := \tfrac{1}{2\alpha}\estwth_{k-i}^\top\estwth_{k-i},\ \forall i\in\{1,\cdots,k\}$.

%DIF >  This recursive analysis utilizes the boundedness of the previously evaluated Lagrange multipliers $\sigma_m,\ m\in\{i+1,\dots,k\},$ and $\lambda_j,\ j\in\mathcal{I}$. 
%DIF >  Furthermore, it relies on the fact that the gradient $\pptfrac{\estNN}{\estwth_i}$ remains bounded provided that the subsequent layers' weights $\estwth_{m},\ m\in\{i+1,\dots,k\}$ are bounded, as shown in \eqref{eq:gradient:phi}.
%DIF >  Once the boundedness of $\estwth_i$ is guaranteed, that of $\sigma_i$ naturally follows: any persistent constraint violation causes the multiplier $\sigma_i$ to increase continuously, which in turn imposes a dominating penalty that ultimately resolves the violation and halts the growth of $\sigma_i$.
%DIF >  Due to the space limit, the detailed analysis is omitted here for brevity.

%DIF >  ——————————————————————————————————————————————
%DIF >  REMARKS
%DIF >  ——————————————————————————————————————————————

\DIFadd{Additionally, we introduce the following remarks regarding the initialization and the parameter tuning of CONAC.
}

\begin{remark}[Initialization of CONAC] \label{rem:init}
  %DIF >  - 정확한 Invariant set 모름
  %DIF >  - 초기 오차를 0에 가깝게 하는게 좋음
  %DIF >  - DNN weight 초기값도 작은게 좋음
  %DIF >  - 너무 크면 과도한 제어로 불안정해지는 경향도 있음
  %DIF >  - 
  \DIFadd{The invariant sublevel set 
  $
    \bigcap_{i\in\mathcal{K}} 
    \mathcal{S}_i(\overline{V}_i)
  $ 
  used in the stability analysis depends on unknown analysis constants.
  Therefore, an exact admissible initial domain is unavailable to compute explicitly.
  In practice, we recommend initializing; the system state $\mv{x}(0)$ close to the initial reference state $\mv{x}_d(0)$ to reduce the initial tracking error $\mv{e}(0)$; DNN weights with small random values; and zero for the Lagrange multipliers, $\sigma_i(0)=0$ and $\lambda_j(0)=0$ for all $i\in\mathcal{K}$ }\DIFaddend and \DIFdelbegin \DIFdel{all imposed constraints $c_j,\ \forall j\in\mathcal I$, are satisfied}\DIFdelend \DIFaddbegin \DIFadd{$j\in\mathcal{I}$, respectively, to help ensure that the system starts within $\mathcal{S}_k(\overline{V}_k)$}\DIFaddend .

  \DIFdelbegin %DIFDELCMD < \end{proof}
%DIFDELCMD < %%%
\DIFdelend \DIFaddbegin \DIFadd{Specifically, this initialization helps prevent initial unstable behavior, which generally occurs by poor initial performance of NAC due to the random initialized DNN weights.
  With small random weights, the initial control input is expected to be small, thereby avoiding abrupt control actions.
  However, if the system starts near an unstable or highly nonlinear operating region, small random weights may result in unstable behavior or slow adaptation.
  Hence, it is preferable to design the reference trajectory such that the system starts in a sufficiently stable operating region, and use a conservative reference trajectory during an initial warm-up phase, as introduced in \mbox{%DIFAUXCMD
\cite{Ryu:2025aa}}\hskip0pt%DIFAUXCMD
.
  This strategy is demonstrated across valuations in Section~\ref{sec:val}.
}\end{remark}
\DIFaddend 

%DIF <   SECTION SIMULATION ======================================
\DIFdelbegin \section{\DIFdel{Real-Time Implementation and Validation}}%DIFAUXCMD
\addtocounter{section}{-1}%DIFAUXCMD
%DIFDELCMD < \label{sec:sim}
%DIFDELCMD < %%%
\DIFdelend \DIFaddbegin \begin{remark}[Tuning of Adaptation Parameters] \label{rem:tuning}
  \DIFadd{The adaptation gain $\alpha$ and the multiplier update rates $\beta_i^\theta$ and $\beta_j^u$ affect the transient behavior of CONAC.
  The gain $\alpha$ determines the update speed of the DNN weights.
  A large $\alpha$ may improve the initial learning response but can also amplify discretization effects and induce oscillatory behavior in sampled-data implementation.
  Conversely, an excessively small $\alpha$ may lead to slow adaptation and poor transient tracking performance.
}

  \DIFadd{The update rates $\beta_i^\theta$ and $\beta_j^u$ determine how rapidly the Lagrange multipliers respond to weight and input constraint violations, respectively.
  Larger values accelerate the enforcement of the corresponding constraints by increasing the multiplier response when a violation occurs.
  However, excessively large update rates may introduce abrupt changes in the adaptation law, leading to chattering or high-frequency oscillations in the control input.
  Therefore, $\alpha$, $\beta_i^\theta$, and $\beta_j^u$ should be selected to balance learning speed, prompt constraint enforcement, and smooth closed-loop behavior.
  The effect of the input-constraint multiplier update rate $\beta_j^u$ is evaluated in the simulation study in Section~\ref{sec:val:sub:sim}.
}\end{remark}
\DIFaddend 

\DIFdelbegin \subsection{\DIFdel{Validation Setup}}
%DIFAUXCMD
\addtocounter{subsection}{-1}%DIFAUXCMD
%DIFDELCMD < 

%DIFDELCMD < %%%
\DIFdelend % ------------------------------------
%DIF <  Real-time Implementation Set-up
%DIF >        SECTION VALIDATIONS
% ------------------------------------
\DIFdelbegin %DIFDELCMD < \begin{figure}[t]
%DIFDELCMD <     \centering
%DIFDELCMD <         \includegraphics[width=.99\linewidth]
%DIFDELCMD <         %%%
%DIF <  \includegraphics[width=\figSizeOneCol\linewidth]
        %DIFDELCMD < {
%DIFDELCMD <             %%%
\DIFdelFL{\src/figures/expSet.
  drawio.
  png
        }%DIFDELCMD < }%%%
%DIF < 
    %DIFDELCMD < \caption{%
{%DIFAUXCMD
\DIFdelFL{Experimental setup of the two-link robotic manipulator used for real-time validation.
    }}
    %DIFAUXCMD
%DIFDELCMD < \label{fig:ctrl:exp:set}
%DIFDELCMD < \end{figure}
%DIFDELCMD < %%%
\DIFdelend \DIFaddbegin \section{\DIFadd{Validations}} \label{sec:val}
\DIFaddend 

\DIFdelbegin \DIFdel{To validate the effectiveness of the proposed CONAC, a }\DIFdelend \DIFaddbegin \DIFadd{In this section, the performance and feasibility of the proposed CONAC are validated via numerical simulations and }\DIFaddend real-time \DIFdelbegin \DIFdel{experiment was conducted on a }\DIFdelend \DIFaddbegin \DIFadd{experiments on a }\DIFaddend two-link \DIFdelbegin \DIFdel{robotic manipulator, as shown in Fig.
  ~\ref{fig:ctrl:exp:set}.
The OpenCR1.0 board \mbox{%DIFAUXCMD
\cite{opencr} }\hskip0pt%DIFAUXCMD
with a 32-bit ARM Cortex and a 216 MHz clock frequency was used for the real-time implementation.
  The two-link robotic manipulator was actuated by two Cubemars AK10-90 servo motors.
  The servos were controlled using the OpenCR1.0 board, which transmits the torque reference to the corresponding servo via controller area network (CAN) communication.
The system dynamics of the two-link robotic manipulator can follow the Euler-Lagrange equation \eqref{eq:sys1}, where $\q$ denotes the joint angles and $\rbu$ denotes the control input (joint torques), and is illustrated in Fig.~\ref{fig:robot:model}.
The measured and /or estimated system parameters are provided in Table~\ref{table:system:params}}\DIFdelend \DIFaddbegin \DIFadd{robotic manipulator in Section~\ref{sec:val:sub:sim} and \ref{sec:val:sub:exp}, respectively}\DIFaddend .

%DIF <  ------------------------------------
%DIF <  Robot Model Figure and Table
%DIF <  ------------------------------------
\DIFaddbegin \subsection{\DIFadd{Validation Setup}}

\DIFaddend \begin{figure}[t]
  \centering
  \DIFdelbeginFL %DIFDELCMD < \subfloat[Two-link robotic manipulator.]{
%DIFDELCMD <         \includegraphics[width=0.499\linewidth]{\src/figures/RobotModel.drawio.pdf}
%DIFDELCMD <         \label{fig:robot:model}
%DIFDELCMD <     }
%DIFDELCMD <     %%%
%DIF <  \hfill
    %DIFDELCMD < \subfloat[Control input saturation function.]{
%DIFDELCMD <         \includegraphics[width=0.499\linewidth]{\src/figures/saturation.drawio.pdf}
%DIFDELCMD <         \label{fig:robot:sat}
%DIFDELCMD <     }
%DIFDELCMD <     %%%
\DIFdelendFL \DIFaddbeginFL \includegraphics[width=0.5\linewidth]{\src/figures/RobotModel.drawio.pdf}
  \DIFaddendFL \caption{Two-link robotic manipulator\DIFdelbeginFL \DIFdelFL{and control saturation function}\DIFdelendFL .}
  \label{fig:robot}
\end{figure}

%DIF <  ------------------------------------
%DIF <  Robot Model Parameters
%DIF <  -----------------------------------
\begin{table}[t]
  \renewcommand{\arraystretch}{1.3}
  \caption{Two-link robotic manipulator's parameters.}
  \centering
  \DIFdelbeginFL %DIFDELCMD < \begin{tabular}{c m{9.5em} c c c }
%DIFDELCMD <     %%%
\DIFdelendFL \DIFaddbeginFL \begin{tabular}{c m{11.5em} c c c }
  \DIFaddendFL % \begin{tabular}{c c c c c}
  \hline
  \textbf{Symbol} & \textbf{Description} & \textbf{Value} \\
  \hline
  \hline 
    $m_1,m_2$ & Mass \DIFaddbeginFL \DIFaddFL{of the links }\DIFaddendFL & \DIFdelbeginFL \DIFdelFL{2.465 $\mathrm{kg}$ }\DIFdelendFL \DIFaddbeginFL \qty{2.465}{\kilo\gram} \DIFaddendFL \\
  \hline
    $l_1,l_2$  & Length \DIFaddbeginFL \DIFaddFL{of the links }\DIFaddendFL & \DIFdelbeginFL \DIFdelFL{0.2 $\mathrm{m}$ }\DIFdelendFL \DIFaddbeginFL \qty{0.2}{\meter} \DIFaddendFL \\  
  \hline
    ${l_c}_1,{l_c}_2$ & \DIFdelbeginFL \DIFdelFL{Center }\DIFdelendFL \DIFaddbeginFL \DIFaddFL{Length to the center }\DIFaddendFL of mass & \DIFdelbeginFL \DIFdelFL{0.139 $\mathrm{m}$ }\DIFdelendFL \DIFaddbeginFL \qty{0.139}{\meter} \DIFaddendFL \\
  \hline
    $I_1,I_2$  & Inertia \DIFaddbeginFL \DIFaddFL{of the links }\DIFaddendFL & \DIFdelbeginFL \DIFdelFL{0.069 $\mathrm{kg\cdot m^{2}}$ }\DIFdelendFL \DIFaddbeginFL \qty{0.069}{\kilo\gram\meter\squared} \DIFaddendFL \\
  \hline
  \end{tabular}
  \label{table:system:params}
\end{table}

%DIF >  -------------------------------------
%DIF >  Robot System
%DIF >  -------------------------------------
\DIFaddbegin \DIFadd{For validations, the two-link robot manipulator (Fig.~\ref{fig:robot}) under control input saturation (Fig.~\ref{fig:saturation:prop:multi}) was considered. 
The model parameters of the manipulator is listed in Table.~\ref{table:system:params}.
To apply CONAC, the two-link robotic manipulator's dynamics is reformulated into the form of \eqref{eq:sys}, by defining the state as $\mv{x} := \mv{q}$, and the control input as $\mv{u} := \mv{\tau}$, where $q_i$ and $\tau_i$ are the $i$-th joint angle and torque, respectively.
}

\DIFadd{In the input saturation, the torque of the second joint, $\tau_2$, is saturated in $\underline{\tau}_2:=\underline{u}_2 = \qty{-3.5}{\newton\meter}$ to $\overline{\tau}_2:=\overline{u}_2 = \qty{3.5}{\newton\meter}$, and the norm of the control input, $\norm{\mv{\tau}}$, is bounded by $\overline{\tau} := \overline{u} = \qty{11}{\newton\meter}$.
These saturation can be physically interpreted as the torque limit of the second joint actuator and the power source, respectively, and are denoted by
}\begin{equation} \label{eq:val:cstr:input}
  c^u_{\overline{\tau}_2}(\mv{\tau}) 
  := 
  \tau_2 - \overline{\tau}_2
  ,
  \ 
  c^u_{\underline{\tau}_2}(\mv{\tau}) 
  := 
  -\tau_2 + \underline{\tau}_2
  ,
  \ 
  c^u_{\overline{\tau}}(\mv{\tau})
  := 
  \norm{\mv{\tau}}^2 - \overline{\tau}^2
  .
\end{equation}
\DIFadd{Accordingly, the set of input constraints is defined as $\mathcal{I} = \{\overline{\tau}_2, \underline{\tau}_2, \overline{\tau}\}$.
}

\DIFaddend % ------------------------------------
% desired trajectory
% ------------------------------------
The desired trajectory is designed with \DIFdelbegin \DIFdel{$4$ segments as follows:
}%DIFDELCMD < \begin{equation}
%DIFDELCMD <     \qd(t) 
%DIFDELCMD <     =
%DIFDELCMD <     \begin{pmatrix}
%DIFDELCMD <         {q_d}_1
%DIFDELCMD <         \\
%DIFDELCMD <         {q_d}_2
%DIFDELCMD <     \end{pmatrix}
%DIFDELCMD <     =
%DIFDELCMD <     \begin{cases}
%DIFDELCMD <     \mathcal{P}^5(t;\mv{q}_d^{(0)},\mv{q}_d^{(1)},t_f),
%DIFDELCMD <         &
%DIFDELCMD <         \text{if } 0 \le t < t_f 
%DIFDELCMD <         ,
%DIFDELCMD <         \\
%DIFDELCMD <     \mathcal{P}^5(t;\mv{q}_d^{(1)},\mv{q}_d^{(2)},t_f),
%DIFDELCMD <         &
%DIFDELCMD <         \text{if } t_f \le t < 2t_f 
%DIFDELCMD <         ,
%DIFDELCMD <         \\
%DIFDELCMD <     \mathcal{P}^5(t;\mv{q}_d^{(2)},\mv{q}_d^{(1)},t_f),
%DIFDELCMD <         & 
%DIFDELCMD <         \text{if } 2t_f \le t < 3t_f 
%DIFDELCMD <         ,
%DIFDELCMD <         \\
%DIFDELCMD <     \mathcal{P}^5(t;\mv{q}_d^{(1)},\mv{q}_d^{(0)},t_f),
%DIFDELCMD <         &    
%DIFDELCMD <         \text{if } 3t_f \le t < 4t_f 
%DIFDELCMD <         ,
%DIFDELCMD <         \\
%DIFDELCMD <     \end{cases} 
%DIFDELCMD <     \label{eq:desired trajectory}
%DIFDELCMD < \end{equation}
%DIFDELCMD < %%%
\DIFdel{where $\mathcal{P}^5(t;\mv{x}_1,\mv{x}_2,t_f)$ denotes }\DIFdelend \DIFaddbegin \DIFadd{4 segments which move the system from 
$
    \mv{q}_d^{(0)}
    :=
    \left(-\tfrac{1}{3}\pi,\tfrac{1}{3}\pi\right)^\top    
$
to 
$
    \mv{q}_d^{(1)}
    :=
    \left(
        \tfrac{1}{4}\pi,-\tfrac{1}{2}\pi
    \right)
$, then to 
$
    \mv{q}_d^{(2)}
    :=
    \left(-\tfrac{1}{4}\pi,\tfrac{1}{4}\pi\right)^\top
$, as illustrated in Fig.~\ref{fig:robot}.
Each segment is designed using }\DIFaddend a quintic ($5$th-order) polynomial trajectory \DIFdelbegin \DIFdel{that moves the system from $\mv{x}_1\in\R^2$ to $\mv{x}_2\in\R^2$ in $t_f$ seconds}\DIFdelend \DIFaddbegin \DIFadd{to ensure the smoothness of the desired trajectory}\DIFaddend , with zero initial and final velocities.
The \DIFdelbegin \DIFdel{initial and final joint angles of the segments are defined as follows:
$
    \mv{q}_d^{(0)}=\left(-\tfrac{1}{3}\pi,\tfrac{1}{3}\pi\right)^\top
$, 
$
    \mv{q}_d^{(1)}=\left(\tfrac{1}{4}\pi,-\tfrac{1}{2}\pi\right)^\top
$, 
$
    \mv{q}_d^{(2)}=\left(-\tfrac{1}{4}\pi,\tfrac{1}{4}\pi\right)^\top
$, as illustrated in Fig.~\ref{fig:robot:model}.
The }\DIFdelend time duration of the segments is defined as \DIFdelbegin \DIFdel{$t_f=3$ s}\DIFdelend \DIFaddbegin \DIFadd{$\qty{3}{\second}$}\DIFaddend .
To demonstrate the learning capability of the DNN, the desired trajectory was applied twice. 
The first and second repetitions of the desired trajectory are referred to as Episode $1$ and Episode $2$, respectively.
%DIF <  ------------------------------------

\DIFdelbegin \DIFdel{However, due to the random initialization of the DNN weights, CONAC's initial performance is expected to be poor, leading to unstable behavior.
Therefore, to }\DIFdelend \DIFaddbegin \DIFadd{Additionally, as mentioned in Remark~\ref{rem:init}, the system is initialized at $\mv{q}(0) = \mv{q}_d^{(0)}$ to ensure a small initial tracking error, and the DNN weights are initialized with small random values.
To }\DIFaddend avoid the instability caused by the initial poor performance of \DIFdelbegin \DIFdel{CONAC}\DIFdelend \DIFaddbegin \DIFadd{DNN}\DIFaddend , a warm-up phase was introduced prior to applying the desired trajectory.
In the warm-up phase, the manipulator was initially positioned at a stable equilibrium point, denoted as ${\qd}^{\rm{origin}} = \left(-\tfrac{1}{2}\pi, 0\right)^\top$\DIFaddbegin \DIFadd{, with zero initial error, }\ie \DIFadd{$\mv{e}(0) = \mv{0}$}\DIFaddend .
Subsequently, a \DIFdelbegin \DIFdel{trajectory $\mathcal{P}^5(t; {\qd}^{\rm{origin}}, \mv{q}_d^{(0)},3)$ was used to move the manipulator from ${\qd}^{\rm{origin}}$ to }\DIFdelend \DIFaddbegin \DIFadd{conservative reference trajectory was applied to gently excite the system, which is designed as a quintic polynomial trajectory that moves the system from ${\qd}^{\rm{origin}}$ to }\DIFaddend $\mv{q}_d^{(0)}$ \DIFdelbegin \DIFdel{over a duration of 3 seconds.
Once the manipulator reached $\mv{q}_d^{(0)}$, a constant desired trajectory $\qd = \mv{q}_d^{(0)}$ was maintained for 8 seconds to allow the system to settle and eliminate the influence of transient dynamics.
}%DIFDELCMD < 

%DIFDELCMD < %%%
\DIFdelend \DIFaddbegin \DIFadd{in $\qty{3}{\second}$, and remains at $\mv{q}_d^{(0)}$ for the next $\qty{8}{\second}$, before applying the desired trajectory.
}\DIFaddend % ------------------------------------
%DIF <  Control Input Saturation
%DIF <  ------------------------------------
\DIFdelbegin \DIFdel{In the validation, we assume that the actuators are subject to a convex saturation function that projects the control input onto the convex set 
$
    \Theta_{\mysat} :=\left\{ \Theta_{\mysat, {\tau_2}} \cap \Theta_{\mysat, {\rbu}} \right\}
$, as illustrated in Fig.~\ref{fig:robot:sat}, where
}%DIFDELCMD < \begin{subequations}
%DIFDELCMD <     \begin{align}
%DIFDELCMD <         \Theta_{\mysat, {\tau_2}}
%DIFDELCMD <         :=&
%DIFDELCMD <         \left\{
%DIFDELCMD <             \tau_2
%DIFDELCMD <             \mid
%DIFDELCMD <             \vert{\tau_2}\vert \le 
%DIFDELCMD <             3.5
%DIFDELCMD <             % \overline\tau_2
%DIFDELCMD <         \right\}
%DIFDELCMD <         ,
%DIFDELCMD <         \label{eq:sat:tau2}
%DIFDELCMD <         \\    
%DIFDELCMD <         \Theta_{\mysat, {\rbu}}
%DIFDELCMD <         :=&
%DIFDELCMD <         \left\{
%DIFDELCMD <             \rbu
%DIFDELCMD <             \mid
%DIFDELCMD <             \norm{\rbu} \le 
%DIFDELCMD <             11
%DIFDELCMD <             % \overline\tau
%DIFDELCMD <         \right\}
%DIFDELCMD <         .
%DIFDELCMD <         \label{eq:sat:norm}
%DIFDELCMD <     \end{align}
%DIFDELCMD <     \label{eq:sat}%%%
%DIF <  %%
%DIFDELCMD < \end{subequations}
%DIFDELCMD < %%%
\DIFdel{The sets defined in \eqref{eq:sat:tau2} and \eqref{eq:sat:norm} can be physically interpreted as the current limit imposed by the torque limit of the second joint actuator and the power source, respectively.
The corresponding input constraints are mathematically described in Appendix~\ref{sec:appen:cstr} as the input bound constraint \eqref{eq:cstr:input:bound} and the input norm constraint \eqref{eq:cstr:input:norm}, respectively}\DIFdelend \DIFaddbegin 

\hfill

%DIF >  -------------------------------------
%DIF >  Comparative Controllers
%DIF >  -------------------------------------
\DIFadd{For comparative analysis, four controllers were implemented, with their properties summarized in Table~\ref{table:controllers}.
}

%DIF >  왜 optimization-based가 배제되었나
%DIF >  - 우리는 structural methods를 기반으로힘
%DIF >  - point-wise constraint handling이 아닌 adaptation 

\DIFadd{In the comparative study, optimization-based methods, such as MPC and CBF-QP, are not included because the focus is on NAC-based controllers.
This choice isolates the proposed contribution, namely the integration of multiple convex input constraints into the online adaptation law of NAC while retaining a Lyapunov-based stability analysis}\DIFaddend .

%DIF <  ------------------------------------
%DIF <  Controller Table
%DIF <  ------------------------------------
\begin{table}[t]
  \renewcommand{\arraystretch}{1.3}
  \caption{\DIFdelbeginFL \DIFdelFL{Properties of the controllers}\DIFdelendFL \DIFaddbeginFL \DIFaddFL{Controllers used for numerical simulations.}\DIFaddendFL }
  \label{table:controller}
  \centering
  \DIFdelbeginFL %DIFDELCMD < \begin{tabular}{m{.25cm} m{4cm} l}
%DIFDELCMD <     %%%
\DIFdelendFL \DIFaddbeginFL \begin{tabular}{m{.7cm} m{3.8cm} l}
  \DIFaddendFL % \begin{tabular}{c l l}
  \hline
    &\bf{Description}&\DIFdelbeginFL %DIFDELCMD < \bf{Input Constraint Handling}%%%
\DIFdelendFL \DIFaddbeginFL \bf{Input Constraint}\DIFaddendFL \\
  \hline
  \hline
    \DIFdelbeginFL %DIFDELCMD < \multirow{3}{*}{(C$_1$)}%%%
\DIFdelendFL \DIFaddbeginFL \multirow{3}{*}{
      \shortstack{
        {(C$_1$)}
        \\
        \protect\colorLine{blue}{solid}
      }
    }\DIFaddendFL & CONAC with & \DIFdelbeginFL \DIFdelFL{\eqref{eq:cstr:input:bound} $c_{\overline\tau_2}$ and $c_{\underline\tau_2}$ }\DIFdelendFL \DIFaddbeginFL \DIFaddFL{$c^u_{\overline{\tau}_2}$ with $\overline{\tau}_2 = 3.5$, 
    }\DIFaddendFL \\
    & large update rate\DIFdelbeginFL \DIFdelFL{$\beta_j$}\DIFdelendFL . & \DIFdelbeginFL \DIFdelFL{with $\overline{\tau}_2=-\underline{\tau}_2=3.5$, 
    and }\DIFdelendFL \DIFaddbeginFL \DIFaddFL{$c^u_{\underline{\tau}_2}$ with $\underline{\tau}_2 = -3.5$, 
    }\DIFaddendFL \\
    & (\DIFdelbeginFL \DIFdelFL{$\beta_{\overline{\tau}_2}=\beta_{\underline{\tau}_2}=100$ and }\DIFdelendFL \DIFaddbeginFL \DIFaddFL{$\beta_{\overline{\tau}_2},\beta_{\underline{\tau}_2}=100$, }\DIFaddendFL $\beta_{\rbu}=10$) & \DIFdelbeginFL \DIFdelFL{\eqref{eq:cstr:input:norm} $c_{\rbu}$ }\DIFdelendFL \DIFaddbeginFL \DIFaddFL{$c^u_{\overline{\tau}}$ }\DIFaddendFL with $\overline{\tau}=11$\DIFaddbeginFL \DIFaddFL{, \eqref{eq:val:cstr:input}}\DIFaddendFL . 
    \\
  \hline
    \DIFdelbeginFL %DIFDELCMD < \multirow{3}{*}{(C$_2$)}%%%
\DIFdelendFL \DIFaddbeginFL \multirow{3}{*}{
      \shortstack{
        {(C$_2$)}
        \\
        \protect\colorLine{cyan}{solid}
      }
    }\DIFaddendFL & CONAC with & \DIFdelbeginFL \DIFdelFL{\eqref{eq:cstr:input:bound} $c_{\overline\tau_2}$ and $c_{\underline\tau_2}$ }\DIFdelendFL \DIFaddbeginFL \DIFaddFL{$c^u_{\overline{\tau}_2}$ with $\overline{\tau}_2 = 3.5$, 
    }\DIFaddendFL \\
    & small update rate\DIFdelbeginFL \DIFdelFL{$\beta_j$}\DIFdelendFL \DIFaddbeginFL \DIFaddFL{, $\beta_j, \forall j\in\mathcal{I}$}\DIFaddendFL . & \DIFdelbeginFL \DIFdelFL{with $\overline{\tau}_2=-\underline{\tau}_2=3.5$, 
    and }\DIFdelendFL \DIFaddbeginFL \DIFaddFL{$c^u_{\underline{\tau}_2}$ with $\underline{\tau}_2 = -3.5$, 
    }\DIFaddendFL \\
    & ($\beta_{\overline{\tau}_2}=\beta_{\underline{\tau}_2}=1$\DIFdelbeginFL \DIFdelFL{and }\DIFdelendFL \DIFaddbeginFL \DIFaddFL{, }\DIFaddendFL $\beta_{\rbu}=0.1$) & \DIFdelbeginFL \DIFdelFL{\eqref{eq:cstr:input:norm} $c_{\rbu}$ }\DIFdelendFL \DIFaddbeginFL \DIFaddFL{$c^u_{\overline{\tau}}$ }\DIFaddendFL with $\overline{\tau}=11$\DIFaddbeginFL \DIFaddFL{, \eqref{eq:val:cstr:input}}\DIFaddendFL . 
    \\
  \hline
    \DIFdelbeginFL %DIFDELCMD < \multirow{2}{*}{(C$_3$)}%%%
\DIFdelendFL \DIFaddbeginFL \multirow{2}{*}{
    \shortstack{  
      (C$_3$) 
      \\
      \protect\colorLine{gray}{solid}
    }
    }\DIFaddendFL & \DIFdelbeginFL \DIFdelFL{CONAC without input constraints.}\DIFdelendFL \DIFaddbeginFL \DIFaddFL{NAC without }\DIFaddendFL & \DIFdelbeginFL %DIFDELCMD < \multirow{2}{*}{ \shortstack[l]{No input constraints\\are imposed.} } %%%
\DIFdelendFL \DIFaddbeginFL \DIFaddFL{No input constraints }\DIFaddendFL \\
    & \DIFdelbeginFL \DIFdelFL{($\beta_{\overline{\tau}_2}=\beta_{\underline{\tau}_2}=\beta_{\rbu}=0$) }\DIFdelendFL \DIFaddbeginFL \DIFaddFL{input constraints \mbox{%DIFAUXCMD
\cite{Patil:2022aa}}\hskip0pt%DIFAUXCMD
. }\DIFaddendFL &  \DIFaddbeginFL \DIFaddFL{is imposed.
    }\DIFaddendFL \\
\hline
    \DIFdelbeginFL %DIFDELCMD < \multirow{2}{*}{(C$_4$)}%%%
\DIFdelendFL \DIFaddbeginFL \multirow{2}{*}{
      \shortstack{  
      (C$_4$)
      \\
      \protect\colorLine{magenta}{solid}
      }
    }\DIFaddendFL & NAC with auxiliary system \DIFdelbeginFL \DIFdelFL{. }\DIFdelendFL & $\overline{\tau}_1=-\underline{\tau}_1=10.428$\DIFaddbeginFL \DIFaddFL{, }\DIFaddendFL \\
    & \DIFdelbeginFL \DIFdelFL{(conventional method) }\DIFdelendFL \DIFaddbeginFL \DIFaddFL{\mbox{%DIFAUXCMD
\cite{Chen:2011aa,Gao:2026aa,Ma:2025aa}}\hskip0pt%DIFAUXCMD
. }\DIFaddendFL & $\overline{\tau_2}=-\underline{\tau_2}=3.5$. \\    
  \hline
  \end{tabular}
  \label{table:controllers}
\end{table}

% ------------------------------------
%DIF <  Comparative Study
%DIF <  ------------------------------------
\DIFdelbegin %DIFDELCMD < 

%DIFDELCMD < \hfill
%DIFDELCMD < 

%DIFDELCMD < %%%
\DIFdel{For comparative analysis, four controllers were implemented, with their properties summarized in Table~\ref{table:controllers}.
}%DIFDELCMD < 

%DIFDELCMD < %%%
%DIF <  ------------------------------------
\DIFdelend % Controllers: CONAC with large beta 
% ------------------------------------
\subsubsection*{Controller 1 \normalfont{(C$_1$)}—\DIFdelbegin \textit{\DIFdel{CONAC with a large update rate $\beta_j,\ \forall j\in\mathcal{I}$, for the Lagrange Multipliers}}%DIFAUXCMD
\DIFdelend \DIFaddbegin \textit{\DIFadd{CONAC with a large update rate $\beta^u_j$ for $\lambda_j,\ \forall j\in\mathcal{I}$}}\DIFaddend }

To handle \DIFdelbegin \DIFdel{input saturation}\DIFdelend \DIFaddbegin \DIFadd{all imposed input saturations}\DIFaddend , the proposed CONAC was implemented with the input bound constraints \DIFdelbegin \DIFdel{$c_{\overline\tau_2}$ and $c_{\underline\tau_2}$ \eqref{eq:cstr:input:bound} }\DIFdelend \DIFaddbegin \DIFadd{$c^u_{\overline\tau_2}$ and $c^u_{\underline\tau_2}$ }\DIFaddend for the second link, and the input norm constraint \DIFdelbegin \DIFdel{$c_{\rbu}$ \eqref{eq:cstr:input:norm}, as follows:
}%DIFDELCMD < \begin{equation}
%DIFDELCMD <     c_{\overline\tau_2}     
%DIFDELCMD <     =
%DIFDELCMD <     \tau_2-\overline\tau_2
%DIFDELCMD <     % \le 0
%DIFDELCMD <     ,
%DIFDELCMD <     \;
%DIFDELCMD <     c_{\underline\tau_2} 
%DIFDELCMD <     =
%DIFDELCMD <     \underline\tau_2-\tau_2
%DIFDELCMD <     % \le 0
%DIFDELCMD <     ,
%DIFDELCMD <     \;
%DIFDELCMD <     c_{\rbu}
%DIFDELCMD <     =
%DIFDELCMD <     \tfrac{1}{2}
%DIFDELCMD <     \left(
%DIFDELCMD <         \norm{\rbu}-\overline\tau
%DIFDELCMD <     \right)^2 
%DIFDELCMD <     % <0
%DIFDELCMD <     ,
%DIFDELCMD <     \label{eq:sim:cstr:CONAC}
%DIFDELCMD < \end{equation}
%DIFDELCMD < %%%
\DIFdel{where $\overline\tau_2=-\underline{\tau}_2=3.5$, and $\overline\tau=11$.
Therefore, controller (C$_1$) utilizes the full capability }\DIFdelend \DIFaddbegin \DIFadd{$c^u_{\rbu}$ defined in \eqref{eq:val:cstr:input}, which allows to fully utilize the authority }\DIFaddend of the actuator\DIFdelbegin \DIFdel{, see \eqref{eq:sat:tau2} and \eqref{eq:sat:norm}.
The }\DIFdelend \DIFaddbegin \DIFadd{.
Additionally, the }\DIFaddend weight norm constraints \DIFdelbegin \DIFdel{$c_{\theta_i}(\cdot),\ \forall i\in\{0,\cdots,k\}$}\DIFdelend \DIFaddbegin \DIFadd{$c^\theta_{i},\ \forall i\in \mathcal{K}$}\DIFaddend , defined in \DIFdelbegin \DIFdel{\eqref{eq:cstr:weight:ball}}\DIFdelend \DIFaddbegin \DIFadd{\eqref{eq:opt:prob:const1}}\DIFaddend , were also incorporated into the CONAC framework with $\overline\theta_0=\overline\theta_1=\overline\theta_2=6$.
As described in \DIFdelbegin \DIFdel{\eqref{eq:control:est}, the control }\DIFdelend \DIFaddbegin \DIFadd{\eqref{eq:ctrl:NN:actual}, the }\DIFaddend input (torque) corresponds to the output of the DNN, and the adaptation law is defined in \DIFdelbegin \DIFdel{\eqref{eq:adaptation law}}\DIFdelend \DIFaddbegin \DIFadd{\eqref{eq:adap}}\DIFaddend .

The update rates of the Lagrange multipliers for the weight norm constraints were set to \DIFdelbegin \DIFdel{$\beta_{\theta_i}=1,\ \forall i\in\{0,\cdots,k\}$}\DIFdelend \DIFaddbegin \DIFadd{$\beta_{\theta_i}=1,\ \forall i\in\mathcal{K}$}\DIFaddend .
For the input constraints, large update rates were used: $\beta_{\overline{\tau}_2} = \beta_{\underline{\tau}_2} = 100$ and $\beta_{\rbu} = 10$, to investigate the effect of large update rates.

% ------------------------------------
% Controllers: CONAC with small beta
% ------------------------------------
\subsubsection*{Controller 2 \normalfont{(C$_2$)}—\DIFdelbegin \textit{\DIFdel{CONAC with a small update rate $\beta_j,\ \forall j\in\mathcal{I}$, for the Lagrange Multipliers}}%DIFAUXCMD
\DIFdelend \DIFaddbegin \textit{\DIFadd{CONAC with a small update rate $\beta^u_j$ for $\lambda_j,\ \forall j\in\mathcal{I}$}}\DIFaddend }

The proposed CONAC was implemented with the same input constraints as in (C$_1$), but with update rates for the input constraints reduced by a factor of $100$: $\beta_{\overline{\tau}_2}=\beta_{\underline{\tau}_2}=1$ and $\beta_{\rbu}=0.1$.
All other properties and parameters were kept identical to those of (C$_1$).

% ------------------------------------
%DIF <  Controllers: CONAC without ctrl cstr
%DIF >  Controllers: NAC without ctrl cstr
% ------------------------------------
\subsubsection*{Controller 3 \normalfont{(C$_3$)}—\DIFdelbegin \textit{\DIFdel{CONAC without input constraints}}%DIFAUXCMD
\DIFdelend \DIFaddbegin \textit{\DIFadd{NAC without input constraints \mbox{%DIFAUXCMD
\cite{Patil:2022aa}}\hskip0pt%DIFAUXCMD
}}\DIFaddend }

The third controller was \DIFaddbegin \DIFadd{employed from \mbox{%DIFAUXCMD
\cite{Patil:2022aa} }\hskip0pt%DIFAUXCMD
which does not handle input constraints.
The controller structure is }\DIFaddend identical to (C$_1$) and (C$_2$), except that the input constraints were removed \DIFdelbegin \DIFdel{to highlight the constraint-handling capability of the proposed CONAC.
In this case, the update rates of the Lagrange multipliers associated with the input constraints, }%DIFDELCMD < \ie %%%
\DIFdel{$\beta_{\overline{\tau}_2},\beta_{\underline{\tau}_2}$ and $\beta_{\rbu}$, were set to zero, such that the corresponding Lagrange multipliers $[\lambda_j]_{j\in\{\overline{\tau}_2,\underline{\tau}_2,\rbu\}}$ were not updated, see \eqref{eq:adap:lbd}}\DIFdelend \DIFaddbegin \DIFadd{and weight boundedness was handled using the projection operator}\DIFaddend .
All other properties and parameters\DIFaddbegin \DIFadd{, $\overline{\theta}_i,\ \forall i\in\mathcal{K}$, }\DIFaddend were kept the same as in (C$_1$) and (C$_2$).
%DIF <  \MSRY{
%DIF <      weight constraint관련된 파라미터는 C1-3까지 모두 같습니다.
%DIF <      C2의 control input 관련 beta만 100배 작아졌다는 말 보충하였습니다.
%DIF <  }

% ------------------------------------
% Controllers: Aux. system in CONAC-AUX
% ------------------------------------
\subsubsection*{Controller 4 \normalfont{(C$_4$)}—\textit{NAC with auxiliary system (conventional method)}}

As a baseline for comparison \DIFaddbegin \DIFadd{in handling input constraints}\DIFaddend , the auxiliary system-based approach \DIFdelbegin \DIFdel{from \mbox{%DIFAUXCMD
\cite{Esfandiari:2014aa, Karason:1994aa, Esfandiari:2015aa} }\hskip0pt%DIFAUXCMD
}\DIFdelend \DIFaddbegin \DIFadd{with nonlinear auxiliary dynamics from \mbox{%DIFAUXCMD
\cite{Chen:2011aa,Gao:2026aa,Ma:2025aa} }\hskip0pt%DIFAUXCMD
}\DIFaddend was implemented.
%DIF <  in place of the proposed CONAC for handling input constraints.
The auxiliary system is defined by 
\DIFdelbegin \DIFdel{$\ddtt{\mv\zeta} = \mm{A}_\zeta \mv{\zeta} + \mm{B}_\zeta \Delta\rbu$, with initial condition $\mv{\zeta}\vert_{t=0} = \mv{0}_{2\times1}$, }\DIFdelend \DIFaddbegin \begin{equation}
  \ddtt \mv{\zeta} =
  \begin{cases}
    -\mm{K}\mv{\zeta}
    -
    \tfrac{
      \vert 
        \mv{\zeta}^\top 
        \Delta^\zeta \mv{u} 
      \vert
      + 
      \tfrac{1}{2} 
      \Delta^\zeta \mv{u}^\top 
      \Delta^\zeta \mv{u}
    }{\norm{\mv{\zeta}}^2}
    \mv{\zeta}
    +
    \Delta^\zeta \mv{u}
    , 
    & 
    \norm{\mv{\zeta}} \ge \mu
    ,
    \\
    \mv{0}
    , 
    & 
    \norm{\mv{\zeta}} < \mu
    ,
  \end{cases}
\end{equation}
\DIFaddend where $\mv{\zeta}\in\R^2$ \DIFdelbegin \DIFdel{represents }\DIFdelend \DIFaddbegin \DIFadd{denotes }\DIFaddend the auxiliary state \DIFdelbegin \DIFdel{variables, and $\mm{A}_\zeta\in\R_{<0}^{2\times 2}$ and $\mm{B}_\zeta\in\R^{2\times 2}$ are user-designed matrices.
The auxiliary state $\mv{\zeta}$ is driven by }\DIFdelend \DIFaddbegin \DIFadd{which is initialized to zero, $\mm{K}\in\R_{>0}^{2\times 2}$ is a positive definite gain matrix, and $\Delta^\zeta$ represents }\DIFaddend the overflow amount of the \DIFdelbegin \DIFdel{control input $\Delta\rbu$, where each element is defined as $\Delta\tau_{i} = \tau_{i}-\max(\underline\tau_i,\min(\overline\tau_i,\tau_i)),\ \forall i\in\{1,2\}$.
}%DIFDELCMD < 

%DIFDELCMD < %%%
\DIFdel{This auxiliary state was incorporated into the adaptation law by replacing the filtered tracking error $\fe$ with $\fe+\mv{\zeta}$.
As a result, the weights were adapted to minimize both the filtered tracking error and the auxiliary state, thereby mitigating saturation in each control input.
}%DIFDELCMD < 

%DIFDELCMD < %%%
\DIFdel{The auxiliary system matrices were chosen as $\mm{A}_\zeta=\mydiag(-10,-10)$ }\DIFdelend \DIFaddbegin \DIFadd{each control input, }\DIFaddend and \DIFdelbegin \DIFdel{$\mm{B}_\zeta=\mydiag(10^3,10^3)$, 
to balance responsiveness to input overflowwith the convergence speed of the auxiliary dynamics.
Since the auxiliary system }\DIFdelend \DIFaddbegin \DIFadd{$\mu > 0$ is a small positive constant.
It should be noted that $\Delta^\zeta\mv{u}$ denotes element-wise input overflow, which is defined as $\Delta^\zeta u_i = u_i - \max(\underline{u}_i, \min(\overline{u}_i, u_i)),\ \forall i\in\{1,2\}$, and thus }\DIFaddend only accounts for \DIFaddbegin \DIFadd{the }\DIFaddend element-wise input saturation\DIFdelbegin \DIFdel{—corresponding to the input bound constraint\eqref{eq:cstr:input:bound}—the upper bound of the first actuator torque $\overline{\tau}_1=-\underline{\tau}_1$ was set to $10.428$ to }\DIFdelend \DIFaddbegin \DIFadd{, without considering the input norm constraint.
The limits of the control input are set as $\underline{u}_1=-\overline{u}_1=10.428$ and $\underline{u}_2=-\overline{u}_2=3.5$, to }\DIFaddend approximate the input norm constraint $\overline{\tau} = 11$\DIFdelbegin \DIFdel{, given $\overline{\tau}_1^2 + \overline{\tau}_2^2 = \overline{\tau}^2$.
}%DIFDELCMD < 

%DIFDELCMD < %%%
\DIFdelend \DIFaddbegin \DIFadd{.
This auxiliary state was incorporated into the control input, $\mv{u}\leftarrow \mv{u} + \mv{\zeta}$.
}\cred{
The parameters of auxiliary system were chosen as $\mu = 0.1$ and $\mm{K} = \mydiag(10,10)$, to balance the responsiveness to input overflow with the convergence speed of the auxiliary dynamics.
}
\DIFaddend All other properties and parameters were kept the same as in (C$_1$) \DIFdelbegin \DIFdel{, }\DIFdelend \DIFaddbegin \DIFadd{and }\DIFaddend (C$_2$)\DIFdelbegin \DIFdel{, and (C$_3$), except that the weight boundedness was handled using the projection operator used in \mbox{%DIFAUXCMD
\cite{Patil:2022aa}}\hskip0pt%DIFAUXCMD
}\DIFdelend .

\hfill

% ------------------------------------
% Controller Parameters
% ------------------------------------
For all controllers, \ie (C$_1$), (C$_2$), (C$_3$), and (C$_4$), the DNNs shared the same architecture, consisting of two hidden layers with four nodes each, \ie $k=2$, and $l_0=6, l_1=4, l_2=4$, and $l_3=2$ in \DIFdelbegin \DIFdel{\eqref{eq:DNN:architecture}}\DIFdelend \DIFaddbegin \DIFadd{\eqref{eq:dnn:def}}\DIFaddend .
The input vector to the DNNs was defined as 
\DIFdelbegin \DIFdel{${\q}_n=\left(\q^\top,\dq^\top,\fe^\top,1\right)^\top\in\R^{7}$, where the scalar $1$ was included to account for the bias term in the weight }\DIFdelend \DIFaddbegin \DIFadd{$
  \NNin
  =
  \left(
    \q^\top,
    \dq^\top,
    \fe^\top,
    1
  \right)^\top
  \in\R^{7}
  ,
$ 
where $\fe$ denotes the filtered tracking error defined as 
$
  \fe 
  = 
  \ddtt\mv{e} + \mm{\Lambda} \mv{e}  
  ,
$
where $\mv{e} = \mv{q} - \mv{q}_d$ is the tracking error and $\mm{\Lambda}=\mydiag(5,15)$ is the filtering }\DIFaddend matrix.
The adaptation gain was set to $\alpha = 0.2$\DIFdelbegin \DIFdel{, and the filtering matrix in \eqref{eq:err:filtered} was chosen as $\mm{\Lambda}=\mydiag(5,15)$.
To support real-time implementation, the controller operated at a sampling rate of $250$ Hz, with control results collected via CAN communication using a PCAN-USB interface and recorded on a computer at the same rate.
}\DIFdelend \DIFaddbegin \DIFadd{.
}\DIFaddend 

The tracking performance was quantitatively evaluated using the $L_2$-norm of the each joint angle's tracking error, defined as $\sqrt{\int_0^T\norm{e_i(t)}\,\der t},\ \forall i\in\{1,2\}$, where $T\in\R_{>0}$ denotes the duration of each episode.

\DIFdelbegin %DIFDELCMD < \begin{figure*}[t]
%DIFDELCMD <     \centering
%DIFDELCMD <     \subfloat[Tracking performance of $1$st joint angle $q_1$.]{
%DIFDELCMD <     %   \includegraphics
%DIFDELCMD <         \includegraphics[width=\figSizeTwoCol\linewidth]
%DIFDELCMD <         {
%DIFDELCMD <         % fig/Fig1.pdf
%DIFDELCMD <         \src/measurement/figures/compare/Fig1.eps
%DIFDELCMD <         }%
%DIFDELCMD <         \label{fig:ctrl:result:q1}}
%DIFDELCMD <     \hfill
%DIFDELCMD <     %%%
%DIF <  \vfill
    %DIFDELCMD < \subfloat[Tracking performance of $2$nd joint angle $q_2$.]{
%DIFDELCMD <         \includegraphics[width=\figSizeTwoCol\linewidth]
%DIFDELCMD <         {
%DIFDELCMD <         % \src/measurement/figures/compare/Fig11.eps
%DIFDELCMD <         \src/measurement/figures/compare/Fig2.eps
%DIFDELCMD <         }%
%DIFDELCMD <         \label{fig:ctrl:result:q2}}
%DIFDELCMD <     \vfill
%DIFDELCMD <     \subfloat[Control input of the $1$st joint $\tau_1$.]{
%DIFDELCMD <     \includegraphics[width=\figSizeTwoCol\linewidth]
%DIFDELCMD <         {
%DIFDELCMD <             \src/measurement/figures/compare/Fig3.eps
%DIFDELCMD <         }%
%DIFDELCMD <         \label{fig:ctrl:result:tau:1}}
%DIFDELCMD <     \hfill
%DIFDELCMD <     %%%
%DIF <  \vfill
        %DIFDELCMD < \subfloat[Control input of the 2nd joint $\tau_2$.]{
%DIFDELCMD <         \includegraphics[width=\figSizeTwoCol\linewidth]
%DIFDELCMD <         {
%DIFDELCMD <             \src/measurement/figures/compare/Fig4.eps
%DIFDELCMD <         }%
%DIFDELCMD <         \label{fig:ctrl:result:tau:2}}
%DIFDELCMD <     \vfill
%DIFDELCMD <     \subfloat[Norm of the joint torque input $\rbu$.]{
%DIFDELCMD <         \includegraphics[width=\figSizeTwoCol\linewidth]
%DIFDELCMD <         {
%DIFDELCMD <             \src/measurement/figures/compare/Fig5.eps
%DIFDELCMD <         }%
%DIFDELCMD <         \label{fig:ctrl:result:tau:norm}}
%DIFDELCMD <     \hfill
%DIFDELCMD <     %%%
%DIF <  \vfill
    %DIFDELCMD < \subfloat[Norm of DNN weights ${\estwth}_i$.]{
%DIFDELCMD <         % \includegraphics
%DIFDELCMD <         \includegraphics[width=\figSizeTwoCol\linewidth]
%DIFDELCMD <         {
%DIFDELCMD <             % fig/Fig7.pdf
%DIFDELCMD <             \src/measurement/figures/compare/Fig7.eps
%DIFDELCMD <         }%
%DIFDELCMD <         \label{fig:ctrl:result:weight}}        
%DIFDELCMD <     %%%
%DIFDELCMD < \caption{%
{%DIFAUXCMD
\DIFdelFL{Real-time implementation results of (C$_1$) }%DIFDELCMD < \protect\colorLine{blue}{solid}%%%
\DIFdelFL{, (C$_2$) }%DIFDELCMD < \protect\colorLine{cyan}{solid}%%%
\DIFdelFL{, (C$_3$) }%DIFDELCMD < \protect\colorLine{gray}{solid}%%%
\DIFdelFL{, (C$_4$) }%DIFDELCMD < \protect\colorLine{magenta}{solid}%%%
\DIFdelFL{, and desired trajectory $\qd$ }%DIFDELCMD < \protect\colorLine{red}{dashed}%%%
\DIFdelFL{.
        %DIF <  \vspace{-5mm}
    }}
%DIFAUXCMD
%DIFDELCMD < \label{fig:ctrl:result}
%DIFDELCMD < \end{figure*}
%DIFDELCMD < %%%
\DIFdelend \DIFaddbegin \subsection{\DIFadd{Comparative Study via Numerical Simulations}} \label{sec:val:sub:sim}
\DIFaddend 

%DIF <  ------------------------------------
%DIF <  Norm of error in Episode 1: 
%DIF <  C1 e1 ep1: 25.130
%DIF <  C1 e2 ep1: 25.016
%DIF <  C2 e1 ep1: 21.719
%DIF <  C2 e2 ep1: 10.530
%DIF <  C3 e1 ep1: 23.168
%DIF <  C3 e2 ep1: 13.061
%DIF <  C4 e1 ep1: 24.195
%DIF <  C4 e2 ep1: 12.927
%DIF <  Norm of error in Episode 2: 
%DIF <  C1 e1 ep2: 7.066
%DIF <  C1 e2 ep2: 2.198
%DIF <  C2 e1 ep2: 8.021
%DIF <  C2 e2 ep2: 2.435
%DIF <  C3 e1 ep2: 8.408
%DIF <  C3 e2 ep2: 3.554
%DIF <  C4 e1 ep2: 19.721
%DIF <  C4 e2 ep2: 4.256
%DIF <  Improvement in Episode 2: 
%DIF <  C1 e1: 0.719
%DIF <  C1 e2: 0.912
%DIF <  C2 e1: 0.631
%DIF <  C2 e2: 0.769
%DIF <  C3 e1: 0.637
%DIF <  C3 e2: 0.728
%DIF <  C4 e1: 0.185
%DIF <  C4 e2: 0.671
%DIF <  ------------------------------------
\DIFaddbegin \subsubsection{\DIFadd{Tracking Performance}}
\DIFaddend 

\DIFdelbegin %DIFDELCMD < \begin{table}[!t]
%DIFDELCMD <     \renewcommand{\arraystretch}{1.3}
%DIFDELCMD <     %%%
%DIFDELCMD < \caption{%
{%DIFAUXCMD
\DIFdelFL{Quantitative\,comparison\,of\,performances'\,$L_2$\,norm.}}
    %DIFAUXCMD
%DIFDELCMD < \centering
%DIFDELCMD <     \begin{tabular}{c c c c c}
%DIFDELCMD <     \hline
%DIFDELCMD < 		& %%%
%DIFDELCMD < \multicolumn{4}{c}{%
{%DIFAUXCMD
\textbf{\DIFdelFL{Episode 1}}%DIFAUXCMD
} %DIFAUXCMD
%DIFDELCMD < \\
%DIFDELCMD <     \hline
%DIFDELCMD < 	\hline 
%DIFDELCMD < 		& %%%
\DIFdelFL{(C$_1$) }%DIFDELCMD < & %%%
\DIFdelFL{(C$_2$) }%DIFDELCMD < & %%%
\DIFdelFL{(C$_3$) }%DIFDELCMD < & %%%
\DIFdelFL{(C$_4$) }%DIFDELCMD < \\
%DIFDELCMD < 	\hline
%DIFDELCMD < 		%%%
\DIFdelFL{$e_1\times10^{3}$ / rad }%DIFDELCMD < & %%%
\DIFdelFL{$25.130$ }%DIFDELCMD < & %%%
\DIFdelFL{$21.719$ }%DIFDELCMD < & %%%
\DIFdelFL{$23.168$ }%DIFDELCMD < & %%%
\DIFdelFL{$24.195$ }%DIFDELCMD < \\ 
%DIFDELCMD < 	\hline
%DIFDELCMD <         %%%
\DIFdelFL{$e_2\times10^{3}$ / rad }%DIFDELCMD < & %%%
\DIFdelFL{$25.016$ }%DIFDELCMD < & %%%
\DIFdelFL{$10.530$ }%DIFDELCMD < & %%%
\DIFdelFL{$13.061$ }%DIFDELCMD < & %%%
\DIFdelFL{$12.927$ }%DIFDELCMD < \\
%DIFDELCMD < 	\hline
%DIFDELCMD <         & %%%
%DIFDELCMD < \multicolumn{4}{c}{%
{%DIFAUXCMD
\textbf{\DIFdelFL{Episode 2}}%DIFAUXCMD
} %DIFAUXCMD
%DIFDELCMD < \\
%DIFDELCMD <     \hline
%DIFDELCMD <     \hline
%DIFDELCMD <         & %%%
\DIFdelFL{(C$_1$) }%DIFDELCMD < & %%%
\DIFdelFL{(C$_2$) }%DIFDELCMD < & %%%
\DIFdelFL{(C$_3$) }%DIFDELCMD < & %%%
\DIFdelFL{(C$_4$) }%DIFDELCMD < \\
%DIFDELCMD < 	\hline
%DIFDELCMD <     \multirow{2}{*}{$e_1\times10^{3}$ / rad} 
%DIFDELCMD <         & %%%
\DIFdelFL{$7.066$ }%DIFDELCMD < & %%%
\DIFdelFL{$8.021$ }%DIFDELCMD < & %%%
\DIFdelFL{$8.408$ }%DIFDELCMD < & %%%
\DIFdelFL{$19.721$ }%DIFDELCMD < \\
%DIFDELCMD <         & %%%
\DIFdelFL{($-71.6\%$) }%DIFDELCMD < & %%%
\DIFdelFL{($-63.1\%$) }%DIFDELCMD < & %%%
\DIFdelFL{($-63.7\%$) }%DIFDELCMD < & %%%
\DIFdelFL{($-18.5\%$) }%DIFDELCMD < \\
%DIFDELCMD <     \hline
%DIFDELCMD <     \multirow{2}{*}{$e_2\times10^{3}$ / rad} 
%DIFDELCMD <         & %%%
\DIFdelFL{$2.198$ }%DIFDELCMD < & %%%
\DIFdelFL{$2.435$ }%DIFDELCMD < & %%%
\DIFdelFL{$3.554$ }%DIFDELCMD < & %%%
\DIFdelFL{$4.256$ }%DIFDELCMD < \\
%DIFDELCMD <         & %%%
\DIFdelFL{($-88.4\%$) }%DIFDELCMD < & %%%
\DIFdelFL{($-76.9\%$) }%DIFDELCMD < & %%%
\DIFdelFL{($-72.8\%$) }%DIFDELCMD < & %%%
\DIFdelFL{($-67.1\%$) }%DIFDELCMD < \\
%DIFDELCMD <     \hline
%DIFDELCMD <     \end{tabular}
%DIFDELCMD <     \label{tab:sim:L2}
%DIFDELCMD < \end{table}
%DIFDELCMD < %%%
\DIFdelend \DIFaddbegin \cred{초기화 + 웜업 페이즈 + 일반적인 학습의 성능}
\DIFaddend 

\DIFdelbegin \subsection{\DIFdel{Validation Results}}
%DIFAUXCMD
\addtocounter{subsection}{-1}%DIFAUXCMD
\DIFdelend \DIFaddbegin \subsubsection{\DIFadd{Constraint Handling Performance}}
\DIFaddend 

%DIF <  ------------------------------------
%DIF <  Tracking Performance
%DIF <  ------------------------------------
\DIFdelbegin \subsubsection{\DIFdel{Tracking Performance}}
%DIFAUXCMD
\addtocounter{subsubsection}{-1}%DIFAUXCMD
\DIFdelend \DIFaddbegin \cred{여러 제약조건을 다루는 것에 대한 유연성}
\DIFaddend 

\DIFdelbegin \DIFdel{The real-time implementation results are presented in Fig.~\ref{fig:ctrl:result}, and the quantitative comparison of tracking performance is summarized in Table~\ref{tab:sim:L2}.
Across the two episodes, all controllers improved their tracking performance. 
Specifically, the $L_2$-norm of $e_1$ was reduced by $71.6 \%$ for (C$_1$), $63.1 \%$ for (C$_2$), $63.7 \%$ for (C$_3$), and $18.5 \%$ for (C$_4$); while the $L_2$-norm of $e_2$ was reduced by $88.4 \%$, $76.9 \%$, $72.8 \%$, and $67.1 \%$, respectively.
Qualitatively, as illustrated in Fig.~\ref{fig:ctrl:result:q1} and Fig.~\ref{fig:ctrl:result:q2}, oscillations were suppressed during the second episode.
These results indicate that the DNNs in all controllers successfully learned the ideal control input $\rbu^*$ using the filtered tracking error $\fe$, despite $\rbu^*$ being completely unknown a priori.
Furthermore, the stable learning observed under random weight initialization and without prior knowledge of the system demonstrated the effectiveness of the proposed CONACin enabling online learning capability.
}\DIFdelend \DIFaddbegin \subsubsection{\DIFadd{KKT-like Optimality of CONAC}}
\DIFaddend 

\DIFdelbegin \DIFdel{Notably, in both episodes, all controllers failed to track the desired trajectory of $q_1$ during certain time intervals: from $17.5$ s to $19.9$ s in Episode $1$, and from $23.5$ s to $25.9$ s and from $29.5$ s to $31.9$ s in Episode $2$, as shown in Fig.~\ref{fig:ctrl:result:q1}. 
This degradation in tracking performance was attributed to actuator saturation during those intervals, as observed in Fig.~\ref{fig:ctrl:result:tau:2} and Fig.~\ref{fig:ctrl:result:tau:norm}.
}%DIFDELCMD < 

%DIFDELCMD < %%%
\DIFdel{The tracking performance of CONAC with input constraints—}%DIFDELCMD < \ie %%%
\DIFdel{(C$_1$) and (C$_2$)—outperformed the auxiliary system-based NAC (C$_4$) in the second episode, as shown in Table~\ref{tab:sim:L2}.
This improvement was attributed to the ability of (C$_1$) and (C$_2$) to fully utilize the actuator's capacity, whereas (C$_4$) was limited by the design of the auxiliary system, which constrained the first joint actuator.
While (C$_3$) also showed better performance than (C$_4$), this was primarily because it did not consider input saturation and thus usedmore of the actuator's capacity.
However, since (C$_3$) did not handle the control input saturation, it generated excessively large control efforts to minimize the tracking error, even though these efforts could not be realized due to actuator limits, as shown in Fig.~\ref{fig:ctrl:result:tau:1}–\ref{fig:ctrl:result:tau:norm}.
In addition, the severity of saturation violations in (C$_3$) increased during periods when the input constraints were active. 
These large violations pose a significant risk to the physical system, potentially causing actuator damage or failure. 
Furthermore, once the control saturation was lifted—typically due to changes in the desired trajectory—this resulted in oscillatory behavior, as CONAC's weights had been adapted to generate excessively high control inputs, see Fig. 
~\ref{fig:ctrl:result:tau:norm}, where the input constraints were lifted at $25.9$ s and $31.9$ s.
}\DIFdelend \DIFaddbegin \cred{제약조건을 지속적으로 위반되더라도 최적점을 찾아가는 것}
\DIFaddend 

%DIF <  ------------------------------------
%DIF <  Control Input Saturation
%DIF <  ------------------------------------
\begin{figure}[t]
    \centering
    \DIFdelbeginFL %DIFDELCMD < \subfloat[Control input $\rbu$ locus during the time interval from $29.4$ s to $31.9$ s.]{
%DIFDELCMD <     % \subfloat[Auxiliary states $\mv{\zeta}$ of CONAC-AUX.]{
%DIFDELCMD <         \includegraphics[width=\figSizeOneCol\linewidth]
%DIFDELCMD <         {
%DIFDELCMD <             \src/measurement/figures/compare/Fig6.eps
%DIFDELCMD <         }%
%DIFDELCMD <         \label{fig:ctrl:result:scope:control}}
%DIFDELCMD <         %%%
\DIFdelendFL \DIFaddbeginFL \subfloat[Control results of Comparative Controllers.]{
    %   \includegraphics
        \includegraphics[width=\figSizeOneCol\linewidth]
        {
        % fig/Fig1.pdf
        \src/script_simulation/figures/compare/Fig1.eps
        }%
        \label{fig:val:sim:general}}
    \DIFaddendFL \vfill
    \DIFdelbeginFL %DIFDELCMD < \subfloat[Lagrange multipliers $\lambda_j$ during the time interval from $29.4$ s to $31.9$ s.]{
%DIFDELCMD <         \includegraphics[width=\figSizeOneCol\linewidth]
%DIFDELCMD <         {
%DIFDELCMD <             \src/measurement/figures/compare/Fig8.eps
%DIFDELCMD <         }%
%DIFDELCMD <         \label{fig:ctrl:result:scope:beta}}
%DIFDELCMD <         \vfill
%DIFDELCMD <     \subfloat[Auxiliary states $\mv{\zeta}$ during the time interval from $29.4$ s to $31.9$ s.]{
%DIFDELCMD <         \includegraphics[width=\figSizeOneCol\linewidth]
%DIFDELCMD <         {
%DIFDELCMD <             \src/measurement/figures/compare/Fig9.eps
%DIFDELCMD <         }%
%DIFDELCMD <         \label{fig:ctrl:result:scope:aux}}
%DIFDELCMD <     %%%
%DIFDELCMD < \caption{%
{%DIFAUXCMD
\DIFdelFL{Constraint handling behavior of (C$_1$) }%DIFDELCMD < \protect\colorLine{blue}{solid}%%%
\DIFdelFL{, (C$_2$) }%DIFDELCMD < \protect\colorLine{cyan}{solid}%%%
\DIFdelFL{, (C$_3$) }%DIFDELCMD < \protect\colorLine{gray}{solid}%%%
\DIFdelFL{, and (C$_4$) }%DIFDELCMD < \protect\colorLine{magenta}{solid} %%%
\DIFdelFL{during the time interval from $29.4$ s to $31.9$ s.
        %DIF <  \vspace{-5mm}
    }}
%DIFAUXCMD
%DIFDELCMD < \label{fig:ctrl:result:scope}
%DIFDELCMD < %%%
\DIFdelendFL \DIFaddbeginFL \subfloat[Control inputs in time interval from 29.5 to 31.5 seconds.]{
        \includegraphics[width=\figSizeOneCol\linewidth]
        {
        % \src/measurement/figures/compare/Fig11.eps
        \src/script_simulation/figures/compare/Fig2.eps
        }%
        \label{fig:val:sim:interested}}     
    \caption{
        \DIFaddFL{Real-time implementation results of (C$_1$) }\protect\colorLine{blue}{solid}\DIFaddFL{, (C$_2$) }\protect\colorLine{cyan}{solid}\DIFaddFL{, (C$_3$) }\protect\colorLine{gray}{solid}\DIFaddFL{, (C$_4$) }\protect\colorLine{magenta}{solid}\DIFaddFL{, and desired trajectory $\qd$ }\protect\colorLine{red}{dashed}\DIFaddFL{.
        %DIF >  \vspace{-5mm}
    }}
\label{fig:val:sim}
\DIFaddendFL \end{figure}

\DIFdelbegin %DIFDELCMD < \begin{figure}[!t]
%DIFDELCMD < 	\centering
%DIFDELCMD < 	\includegraphics[width=0.65\linewidth]{
%DIFDELCMD < 		\src/figures/lambda_effect.drawio.pdf
%DIFDELCMD < 		}
%DIFDELCMD < 	%%%
%DIFDELCMD < \caption{%
{%DIFAUXCMD
\DIFdelFL{The effect of the Lagrange multiplier $\lambda_{j}$ on the adaptation direction of $\widehat{\mv{\theta}}_2$ is illustrated. 
        The notations $(\cdot^{\prime})$ and $(\cdot^{\prime\prime})$ denote two distinct cases corresponding to large and small values of the Lagrange multiplier, respectively, }%DIFDELCMD < \ie %%%
\DIFdelFL{$\lambda_{j}^{\prime} > \lambda_{j}^{\prime\prime}$.
        %DIF <  \vspace{-5mm}
    }}
	%DIFAUXCMD
%DIFDELCMD < \label{fig:lambda_effect}
%DIFDELCMD < \end{figure}
%DIFDELCMD < %%%
\DIFdelend \DIFaddbegin \subsection{\DIFadd{Feasibility Study via Real-Time Experiments}} \label{sec:val:sub:exp}
\DIFaddend 


\DIFdelbegin %DIFDELCMD < \hfill
%DIFDELCMD < 

%DIFDELCMD < %%%
\subsubsection{\DIFdel{Constraint Handling}} %DIFAUXCMD
\addtocounter{subsubsection}{-1}%DIFAUXCMD
%DIFDELCMD < \label{sec:sim:constraint}
%DIFDELCMD < 

%DIFDELCMD < %%%
\DIFdel{In this section, the constraint-handling capability of all controllers is investigated.
As shown in Fig.~\ref{fig:ctrl:result:tau:1}–\ref{fig:ctrl:result:tau:norm}, all controllers except (C$_3$), which did not consider the control input saturation, satisfied the control input saturation illustrated in Fig.~\ref{fig:robot:sat}.
In (C$_1$) and (C$_2$), control input saturation was addressed through corresponding input constraints within the constrained optimization framework, whereas (C$_4$) employed an auxiliary system, as summarized in Table~\ref{table:controllers}.
}%DIFDELCMD < 

%DIFDELCMD < %%%
\DIFdel{To investigate the constraint handling process, the time interval from $29.4$ s to $31.9$ s was selected, during which the control input saturation was active for all controllers.
In Fig.~\ref{fig:ctrl:result:scope}, the control input locus, the Lagrange multipliers, and the auxiliary states are illustrated over this time interval.
The properties of the controllers are illustrated in Fig.~\ref{fig:ctrl:result:scope:control}, where the control input loci of (C$_1$) and (C$_2$) fully explore the feasible input domain shown in Fig.~\ref{fig:robot:sat}, whereas (C$_3$) exceeds the input saturation limits, and (C$_4$) is restricted in its first actuator input limit due to the auxiliary system design.
}%DIFDELCMD < 

%DIFDELCMD < %%%
%DIF <  beta
\DIFdel{The input constraint handling from $29.4$ s to $31.9$ s is examined in detail. 
For (C$_1$) and (C$_2$), activation of input constraints $c_j$ dynamically generated the corresponding Lagrange multipliers $\lambda_j$, as shown in Fig.~\ref{fig:ctrl:result:scope:beta}. 
These multipliers guided the adaptation of the weights $\estwth_i,\ \forall i\in\{0,1,2\}$, to reduce constraint violations in accordance with \eqref{eq:adap:th}.
}%DIFDELCMD < 

%DIFDELCMD < %%%
%DIF <  comparison; big, small beta
\DIFdel{To further illustrate this effect, Fig.~\ref{fig:lambda_effect} compares the adaptation direction of the output layer's weight $\estwth_2$ under two different magnitudes of $\lambda_j$. 
Under Assumptions \ref{assum:convex} and \ref{assum:LICQ}, the convexity of $c_j$ in the $\estwth_2$-space (see Lemma~\ref{lem:convex:angle}) is illustrated in Fig.~\ref{fig:lambda_effect}.
The figure highlights the influence of the update rate $\beta_j$: a larger $\beta_j$ (as in (C$_1$)) yields a larger $\lambda_j$, resulting in a more aggressive adjustment of $\estwth_2$ via the term $-\lambda_j\pptfrac{c_j}{\estwth_2}$ in \eqref{eq:adap:th}, thereby accelerating constraint satisfaction—though at the cost of oscillatory learning behavior.
In contrast, the smaller update rate in (C$_2$) leads to more conservative weight updates andsmoother, but slower, convergence.
These behaviors are clearly observed in Fig.~\ref{fig:ctrl:result:scope:control} and Fig.~\ref{fig:ctrl:result:scope:beta}, where (C$_1$) satisfies the input constraint $c_{\rbu}$ more quickly than (C$_2$), which shows repeated growth and reset of $\lambda_{\rbu}$}\DIFdelend \DIFaddbegin \DIFadd{To validate the feasibility of the real-time implementation of the proposed CONAC, a real-time experiment was conducted on a two-link robotic manipulator, as shown in Fig}\DIFaddend .\DIFdelbegin %DIFDELCMD < 

%DIFDELCMD < %%%
%DIF <  zeta
\DIFdel{For (C$_4$), the auxiliary state $\zeta_1$ was generated in response to the control input saturation of the first link actuator $\tau_1$, as shown in Fig.~\ref{fig:ctrl:result:scope:aux}.
Since the weights were adapted to reduce both the filtered error $\fe$ and the auxiliary state $\mv{\zeta}$, }\DIFdelend \DIFaddbegin \DIFadd{~\ref{fig:ctrl:exp:set}. 
The OpenCR1.0 board \mbox{%DIFAUXCMD
\cite{opencr} }\hskip0pt%DIFAUXCMD
with a 32-bit ARM Cortex and a $\qty{216}{\mega\hertz}$ clock frequency was used.
The two-link robotic manipulator was actuated by two Cubemars AK10-90 servo motors. 
The servos were controlled using the OpenCR1.0 board, which transmits the torque reference to }\DIFaddend the \DIFdelbegin \DIFdel{control input saturation was subsequently reduced, as evident in Fig.~\ref{fig:ctrl:result:scope:control}.
Once the saturation was resolved at around $31.2$ s (see Fig.~\ref{fig:ctrl:result:tau:1}), the auxiliary state $\zeta_1$ converged to zero due to the stable state matrix $\mm{A}_\zeta$ in the auxiliary system dynamics, given by $\ddtt\mv{\zeta} = \mm{A}_\zeta \mv{\zeta} + \mm{B}_\zeta \Delta\rbu$ with $\Delta\rbu=\mv{0}_{2\times1}$, as shown again in Fig.
~\ref{fig:ctrl:result:scope:aux}.
}\DIFdelend \DIFaddbegin \DIFadd{corresponding servo via controller area network (CAN) communication at a sampling rate of $\qty{250}{\hertz}$ using a PCAN-USB interface.
Since the measured and/or estimated system parameters are used for the numerical validation in Section~\ref{sec:val:sub:sim}, system parameters for the real-time experiment can be found in Table~\ref{table:system:params}, as well.
}\DIFaddend 

\DIFdelbegin %DIFDELCMD < \hfill 
%DIFDELCMD < 

%DIFDELCMD < %%%
\DIFdel{Finally, the handling of the weight norm constraints $c_{\theta_i},\ \forall i\in\{0,1,2\}$, is investigated. 
Among all controllers, the weight norm constraint was only activated for the output layer weights $\estwth_2$ in (C$_3$), as shown in Fig.~\ref{fig:ctrl:result:weight}. 
This is because the weights in (C}\DIFdelend \DIFaddbegin \DIFadd{Same scenario including desired trajectories and constraints were used and (C}\DIFaddend $_1$) \DIFdelbegin \DIFdel{, }\DIFdelend \DIFaddbegin \DIFadd{and }\DIFaddend (C$_2$) \DIFdelbegin \DIFdel{, and (C$_4$) were implicitly suppressed by the input constraints, which were not applied in (C3).
%DIF <  , \ie by limiting the control input, the weight norms were also limited.
%DIF <  \MSRY{필요하다면 왜 제어입력 제약하는게 가중치 제약도 되는지 추가 가능}
Furthermore, the constraints on the inner and hidden layer weights were never violated; that is, both $\norm{\estwth_0}$ and $\norm{\estwth_1}$ consistently remained below their limits.
This was attributed to the fact that the backpropagated error was scaled by the gradients of the activation function, which are bounded in $(0,1]$, }%DIFDELCMD < \ie %%%
\DIFdel{$\sigma(\cdot)=\tanh(\cdot)$.
}\DIFdelend \DIFaddbegin \DIFadd{were implemented.
}\DIFaddend 

\DIFdelbegin \DIFdel{Additional details on the constraint handling process are provided in Fig.~\ref{fig:ctrl:result:scope:beta}.Similar to the case of the input constraints, the Lagrange multipliers $\lambda_{\theta_2}$ increased at around $30.4$ s, as shown in Fig.~\ref{fig:ctrl:result:scope:beta}, when $\estwth_2$ exceeded the prescribed norm constraint $\overline\theta_2$ (see Fig.~\ref{fig:ctrl:result:weight}), thereby adjusting the adaptation direction of $\estwth_2$ to mitigate the violation of the constraint $c_{\theta_2}$.
}\DIFdelend \DIFaddbegin \begin{figure}[t]
    \centering
        \includegraphics[width=.99\linewidth]
        %DIF >  \includegraphics[width=\figSizeOneCol\linewidth]
        {
            \src/figures/expSet.drawio.png
        }%DIF > 
    \caption{
        \DIFaddFL{Experimental setup of the two-link robotic manipulator used for real-time validation.
    }}
    \label{fig:ctrl:exp:set}
\end{figure}
\DIFaddend 

%DIF <  In addition, it is worth noting that the input constraints were repeatedly violated even after the constraint handling mechanisms were applied, as shown in Fig.~\ref{fig:ctrl:result:scope:control}. 
%DIF <  This occurred because the controllers were implemented in discrete time with a relatively low sampling rate of $250$ Hz.
\DIFdelbegin %DIFDELCMD < 

%DIFDELCMD <     %%%
%DIF <  % ------------------------------------
    %DIF <  % Karush-Kuhn-Tucker (KKT) conditions}
    %DIF <  % ------------------------------------
    %DIF <  \hfill
%DIFDELCMD < 

%DIFDELCMD <     %%%
%DIF <  \subsubsection{Satisfaction of Karush-Kuhn-Tucker (KKT) conditions}
%DIFDELCMD < 

%DIFDELCMD <     %%%
%DIF <  Since the proposed CONAC was implemented in discrete time, a direct demonstration of the KKT conditions at steady state is not straightforward. 
    %DIF <  % % \MSRY{
    %DIF <  % %     그림을 잘 확인해보니 제약조건이 위반되는 구간동안 정상상태에 도달하지 못하더라구요... 이전 논문 결과에서와 마찬가지로 그러할것으로 생각했는데 다른 요인이 있는 것 같습니다. 삭제하던지 아니면 보여질수없는 이유를 더 설명해야할 것으로 보입니다.
    %DIF <  % % }
    %DIF <  % % \ctxt{red}
    %DIF <  % % {
    %DIF <  % %     However, the near-constant values of the weights and Lagrange multipliers during constraint violation suggest an implicit satisfaction of the KKT conditions, as shown in Fig.~\ref{fig:ctrl:result:weight} and Fig.~\ref{fig:ctrl:result:scope:beta}
    %DIF <  % %     Specifically, the Lagrange multipliers corresponding to active constraints were positive, while those associated with inactive constraints remained zero.
    %DIF <  % % }
    %DIF <  % % \ctxt{blue}
    %DIF <  % % {
    %DIF <      This is because, the KKT conditions are typically defined for continuous-time systems, and their direct application to discrete-time systems requires careful consideration of the sampling rate.
    %DIF <      Although exact satisfaction of the KKT conditions was not observed—due to the lack of a clear steady state in the discrete-time implementation—the control results remain sufficiently close to the optimal behavior implying that the KKT conditions were approximately satisfied.
    %DIF <  % % }
%DIFDELCMD < 

%DIFDELCMD < %%%
%DIF <  ------------------------------------
%DIF <  Computational Time
%DIF <  ------------------------------------
\DIFdelend \hfill

\DIFdelbegin \subsubsection{\DIFdel{Effectiveness of CONAC in Real-time Applications}}
%DIFAUXCMD
\addtocounter{subsubsection}{-1}%DIFAUXCMD
%DIFDELCMD < 

%DIFDELCMD < %%%
\DIFdel{The effectiveness of the proposed CONAC in }\DIFdelend \DIFaddbegin \DIFadd{The feasibility of }\DIFaddend real-time \DIFdelbegin \DIFdel{applications was evaluated by measuring the computational time using CPU clock ticks on the OpenCR1.0 board. 
As shown in Fig.~\ref{fig:ctrl:real:result:cmp:time}, the computational times of all controllers remained below $4$ ms, demonstrating that }\DIFdelend \DIFaddbegin \DIFadd{implementation of }\DIFaddend the proposed CONAC is \DIFdelbegin \DIFdel{suitable for real-time implementation at a sampling rate of $250$ Hz}\DIFdelend \DIFaddbegin \DIFadd{analyzed}\DIFaddend .


\DIFdelbegin %DIFDELCMD < \begin{figure}[t]
%DIFDELCMD <     %%%
\DIFdelendFL \DIFaddbeginFL \begin{figure*}[t]
    \DIFaddendFL \centering
    \DIFdelbeginFL %DIFDELCMD < \includegraphics[width=\figSizeOneCol\linewidth]
%DIFDELCMD <         {
%DIFDELCMD <             %%%
\DIFdelFL{\src/measurement/figures/compare/Fig10.eps
        }%DIFDELCMD < }%%%
%DIF < 
    \DIFdelendFL \DIFaddbeginFL \subfloat[Tracking performance of $1$st joint angle $q_1$.]{
    %   \includegraphics
        \includegraphics[width=\figSizeTwoCol\linewidth]
        {
        % fig/Fig1.pdf
        \src/measurement/figures/compare/Fig1.eps
        }%
        \label{fig:ctrl:result:q1}}
    \hfill
    %DIF >  \vfill
    \subfloat[Tracking performance of $2$nd joint angle $q_2$.]{
        \includegraphics[width=\figSizeTwoCol\linewidth]
        {
        % \src/measurement/figures/compare/Fig11.eps
        \src/measurement/figures/compare/Fig2.eps
        }%
        \label{fig:ctrl:result:q2}}
    \vfill
    \subfloat[Control input of the $1$st joint $\tau_1$.]{
    \includegraphics[width=\figSizeTwoCol\linewidth]
        {
            \src/measurement/figures/compare/Fig3.eps
        }%
        \label{fig:ctrl:result:tau:1}}
    \hfill
    %DIF >  \vfill
        \subfloat[Control input of the 2nd joint $\tau_2$.]{
        \includegraphics[width=\figSizeTwoCol\linewidth]
        {
            \src/measurement/figures/compare/Fig4.eps
        }%
        \label{fig:ctrl:result:tau:2}}
    \vfill
    \subfloat[Norm of the joint torque input $\rbu$.]{
        \includegraphics[width=\figSizeTwoCol\linewidth]
        {
            \src/measurement/figures/compare/Fig5.eps
        }%
        \label{fig:ctrl:result:tau:norm}}
    \hfill
    %DIF >  \vfill
    \subfloat[Norm of DNN weights ${\estwth}_i$.]{
        % \includegraphics
        \includegraphics[width=\figSizeTwoCol\linewidth]
        {
            % fig/Fig7.pdf
            \src/measurement/figures/compare/Fig7.eps
        }%
        \label{fig:ctrl:result:weight}}        
    \DIFaddendFL \caption{
        \DIFdelbeginFL \DIFdelFL{Computational time $T_{\rm{comp}}$ }\DIFdelendFL \DIFaddbeginFL \DIFaddFL{Real-time implementation results }\DIFaddendFL of (C$_1$) \protect\colorLine{blue}{solid}, (C$_2$) \protect\colorLine{cyan}{solid}, (C$_3$) \protect\colorLine{gray}{solid}, \DIFdelbeginFL \DIFdelFL{and }\DIFdelendFL (C$_4$) \protect\colorLine{magenta}{solid}\DIFaddbeginFL \DIFaddFL{, and desired trajectory $\qd$ }\protect\colorLine{red}{dashed}\DIFaddendFL .
        % \vspace{-5mm}
    }
\DIFdelbeginFL %DIFDELCMD < \label{fig:ctrl:real:result:cmp:time}
%DIFDELCMD < \end{figure}
%DIFDELCMD < %%%
\DIFdelend \DIFaddbegin \label{fig:ctrl:result}
\end{figure*}
\DIFaddend 


%DIF <   SECTION CONCLUSION ======================================
%DIF >  ------------------------------------
%DIF >        SECTION CONCLUSION
%DIF >  ------------------------------------
\section{Conclusion} \label{sec:conclusion}

This paper proposed a constrained optimization-based neuro-adaptive controller (CONAC) for unknown \DIFdelbegin \DIFdel{Euler-Lagrange }\DIFdelend systems, addressing \DIFdelbegin \DIFdel{both weight norm and }\DIFdelend \DIFaddbegin \DIFadd{multiple }\DIFaddend convex input constraints within a unified \DIFaddbegin \DIFadd{constrained }\DIFaddend optimization framework. 
The stability of CONAC was analyzed using Lyapunov theory, demonstrating bounded tracking errors and \DIFdelbegin \DIFdel{weight estimation under real-time adaptation.
}%DIFDELCMD < 

%DIFDELCMD < %%%
\DIFdel{CONAC effectively incorporated both convex input constraints and weight norm constraints, ensuring that actuator limitations and neural network weights remained within predefined bounds}\DIFdelend \DIFaddbegin \DIFadd{the adaptive variables, including the deep neural network (DNN) weights and Lagrange multipliers, under the proposed adaptation law}\DIFaddend .
By embedding the\DIFdelbegin \DIFdel{se}\DIFdelend  constraints into the optimization process, CONAC ensured weight convergence in accordance with the Karush-Kuhn-Tucker (KKT)\DIFaddbegin \DIFadd{-like optimality }\DIFaddend conditions, thereby ensuring both stability and optimality.

\DIFdelbegin \DIFdel{Real-time implementation validated the superior performance of CONAC compared to a conventional method. 
CONAC successfully handled complex input constraints while rigorously managing weight norms, resulting in improved tracking accuracy and stable behavior without significant oscillations.
}\DIFdelend %DIF >  CONAC effectively incorporated multiple convex input constraints, ensuring that actuator limitations and closed-loop stability.

\DIFaddbegin \DIFadd{Numerical simulations demonstrated the flexibility of CONAC in handling various convex input constraints simultaneously, and the effect of the update rates of the Lagrange multipliers on enforcement of input constraints.
Additionally, real-time implementation validated the feasibility of CONAC in real-world applications.
}

\DIFaddend Future work may extend this approach to include constraints on both control inputs and system states, \DIFaddbegin \DIFadd{and minimize the infinite horizon error cost, thereby }\DIFaddend further enhancing the flexibility and \DIFdelbegin \DIFdel{robustness }\DIFdelend \DIFaddbegin \DIFadd{optimality }\DIFaddend of neuro-adaptive control systems within constrained optimization frameworks.

%DIF < % The Appendices part is started with the command \appendix;
%DIF < % appendix sections are then done as normal sections
\DIFdelbegin %DIFDELCMD < \appendix
%DIFDELCMD < %%%
\DIFdelend %DIF >  -------------------------------------
%DIF >  Appendix
%DIF >  -------------------------------------
%DIF >  \appendix

\DIFdelbegin \subsection{\DIFdel{Input Constraint Candidates}}%DIFAUXCMD
\addtocounter{subsection}{-1}%DIFAUXCMD
%DIFDELCMD < \label{sec:appen:cstr} 
%DIFDELCMD < 

%DIFDELCMD < %%%
\DIFdel{This section introduces potential input constraints that can be incorporated into the proposed CONAC framework.
}%DIFDELCMD < 

%DIFDELCMD < %%%
\subsubsection{\DIFdel{Input Bound Constraint}}%DIFAUXCMD
\addtocounter{subsubsection}{-1}%DIFAUXCMD
%DIFDELCMD < \label{sec:appen:cstr:input:bound}
%DIFDELCMD < 

%DIFDELCMD < %%%
\DIFdel{Most physical systems are subject to control input limitations due to inherent electrical and mechanical constraints.
These are expressed as $\mv{c}_{\overline \tau}:= [c_{\overline \tau_i}]_{i\in\{1,\cdots,n\}}$ and $\mv{c}_{\underline\tau}:= [c_{\underline\tau_i}]_{i\in\{1,\cdots,n\}}$, where
}%DIFDELCMD < \begin{equation}
%DIFDELCMD <     \begin{aligned}
%DIFDELCMD <         c_{\overline \tau_i}=\tau_{(i)} - {\tau_{\overline \tau_i}}
%DIFDELCMD <         ,
%DIFDELCMD <         \quad
%DIFDELCMD <         c_{\underline\tau_i}={\tau_{\underline\tau_i}}-\tau_{(i)}
%DIFDELCMD <         ,
%DIFDELCMD <     \end{aligned}
%DIFDELCMD <     \label{eq:cstr:input:bound}
%DIFDELCMD < \end{equation}
%DIFDELCMD < %%%
\DIFdel{with $\tau_{\overline \tau_i}\in\R$ and $\tau_{\underline\tau_i}\in\R$ representing the maximum and minimum control input bounds, respectively.
The gradients of $\mv{c}_{\overline \tau}$ and $\mv{c}_{\underline\tau}$ with respect to $\estwth$ are given by
}%DIFDELCMD < \begin{equation}
%DIFDELCMD <     \begin{aligned}
%DIFDELCMD <         \pptfrac{\mv c_{\overline \tau}}{\estwth}
%DIFDELCMD <         =& 
%DIFDELCMD <         % \left[
%DIFDELCMD <         %     \pptfrac{c_{\overline \tau_i}}{\estwth}^\top
%DIFDELCMD <         % \right]_{i\in\{1,\cdots,n\}}
%DIFDELCMD <         % \begin{bmatrix}
%DIFDELCMD <         %     \pptfrac{c_{\overline \tau_1}}{\estwth}^\top \\
%DIFDELCMD <         %     \vdots \\
%DIFDELCMD <         %     \pptfrac{c_{\overline \tau_n}}{\estwth}^\top
%DIFDELCMD <         % \end{bmatrix}
%DIFDELCMD <         = 
%DIFDELCMD <             +\pptfrac{\estNN}{\estwth}
%DIFDELCMD <             % \\
%DIFDELCMD <         % =&
%DIFDELCMD <         =
%DIFDELCMD <         +
%DIFDELCMD <         \begin{bmatrix}
%DIFDELCMD <             (\mm{I}_{l_{k+1}}\otimes \estact_{k}^\top)&
%DIFDELCMD <             \!\cdots\! &
%DIFDELCMD <             (
%DIFDELCMD <                 \cdot
%DIFDELCMD <             )
%DIFDELCMD <         \end{bmatrix} 
%DIFDELCMD <         \in
%DIFDELCMD <         \mathbb R^{n\times \Xi}
%DIFDELCMD <         , 
%DIFDELCMD <         \\
%DIFDELCMD <         \pptfrac{\mv c_{\underline\tau}}{\estwth}         
%DIFDELCMD <         =
%DIFDELCMD <         & 
%DIFDELCMD <         % \left[
%DIFDELCMD <         %     \pptfrac{c_{\underline\tau_i}}{\estwth}^\top
%DIFDELCMD <         % \right]_{i\in\{1,\cdots,n\}}
%DIFDELCMD <         % \begin{bmatrix}
%DIFDELCMD <         %     \pptfrac{c_{\underline\tau_1}}{\estwth}^\top \\
%DIFDELCMD <         %     \vdots \\
%DIFDELCMD <         %     \pptfrac{c_{\underline\tau_n}}{\estwth}^\top
%DIFDELCMD <         % \end{bmatrix}
%DIFDELCMD <         = 
%DIFDELCMD <         -\pptfrac{\estNN}{\estwth}
%DIFDELCMD <         % \\
%DIFDELCMD <         % =&
%DIFDELCMD <         =
%DIFDELCMD <         -\begin{bmatrix}
%DIFDELCMD <             (\mm{I}_{l_{k+1}}\otimes \estact_{k}^\top)&
%DIFDELCMD <             % V_k^\top \phi_{k}' (I_{l_{k}}\otimes  \phi_{k-1}^\top)&
%DIFDELCMD <             \!\cdots\! &
%DIFDELCMD <             (
%DIFDELCMD <                 \cdot
%DIFDELCMD <             )
%DIFDELCMD <         \end{bmatrix} 
%DIFDELCMD <         \in
%DIFDELCMD <         \mathbb R^{n\times \Xi}
%DIFDELCMD <         .
%DIFDELCMD <     \end{aligned}
%DIFDELCMD <     \label{eq:cstr:input:bound:grad}
%DIFDELCMD < \end{equation}
%DIFDELCMD < 

%DIFDELCMD < %%%
\subsubsection{\DIFdel{Input Norm Constraint}}%DIFAUXCMD
\addtocounter{subsubsection}{-1}%DIFAUXCMD
%DIFDELCMD < \label{sec:appen:cstr:input:norm}
%DIFDELCMD < 

%DIFDELCMD < %%%
\DIFdel{Consider the control input $\rbu$ as the torque corresponding to its generalized coordinate. 
Since torque is typically linearly proportional to current, actuators that share a common power source are often subject to total current limitations. 
This can be captured by the following inequality constraint: 
}%DIFDELCMD < \begin{equation}
%DIFDELCMD <     c_{\rbu}
%DIFDELCMD <     =
%DIFDELCMD <     \tfrac{1}{2}\
%DIFDELCMD <     \left(
%DIFDELCMD <         \norm{
%DIFDELCMD <         \mv{\tau}
%DIFDELCMD <         }^2 
%DIFDELCMD <         -
%DIFDELCMD <         \overline\tau^2
%DIFDELCMD <     \right)
%DIFDELCMD <     ,
%DIFDELCMD <     \label{eq:cstr:input:norm}
%DIFDELCMD < \end{equation}
%DIFDELCMD < %%%
\DIFdel{with the maximum allowable control input magnitude $\overline\tau\in\R_{>0}$. This input norm constraint is also commonly applied in current and torque control problems for electric motors \mbox{%DIFAUXCMD
\cite{Choi:2024aa}}\hskip0pt%DIFAUXCMD
.
%DIF <  The corresponding Lagrangian multiplier is $\lambda_{u_b}$. 
The gradient of $c_{\rbu}$ with respect to $\estwth$ is given by
}%DIFDELCMD < \begin{equation}
%DIFDELCMD <     \pptfrac{\mv c_{\rbu}}{\estwth}
%DIFDELCMD <     \!=\! 
%DIFDELCMD <     \textstyle\sum_{i=1}^n \tau_{i} 
%DIFDELCMD <     \left(
%DIFDELCMD <         \myrow_i
%DIFDELCMD <         \left(
%DIFDELCMD <             \pptfrac{\estNN}{\estwth}
%DIFDELCMD <         \right)
%DIFDELCMD <     \right)^\top  
%DIFDELCMD <     \!=\! 
%DIFDELCMD <     {\pptfrac{\estNN}{\estwth}}^\top
%DIFDELCMD <     \mv{\tau}
%DIFDELCMD <     % \tau^\top (I_{l_{k+1}}\otimes \estNN_k^\top)
%DIFDELCMD <     \in \mathbb R^{\Xi}.
%DIFDELCMD <     \label{eq:cstr:input:norm:grad}
%DIFDELCMD < \end{equation}
%DIFDELCMD < 

%DIFDELCMD < \hfill
%DIFDELCMD < 

%DIFDELCMD < %%%
\DIFdel{It should be noted that constraints \eqref{eq:cstr:input:bound} and \eqref{eq:cstr:input:norm} can be imposed simultaneously, as their gradients \eqref{eq:cstr:input:bound:grad} and \eqref{eq:cstr:input:norm:grad} are linearly independent, satisfying Assumption \ref{assum:LICQ}, }%DIFDELCMD < \ie %%%
\DIFdel{the LICQ condition.
}%DIFDELCMD < 

%DIFDELCMD < %%%
\DIFdelend % ------------------------------------
% Bibliography
% ------------------------------------
\bibliographystyle{IEEEtran}
\DIFdelbegin %DIFDELCMD < \bibliography{\template/refs}
%DIFDELCMD < %%%
%DIF <  \bibliography{\template/refs_abbr}
\DIFdelend %DIF >  \bibliographystyle{apalike}
\DIFaddbegin \bibliography{localRefs, fromLR}
\DIFaddend 

% ------------------------------------
% Biography
% ------------------------------------
% \vspace{-10mm}
\begin{IEEEbiography}[{
    \includegraphics[width=1in,height=1.25in,clip,keepaspectratio]{\template/people/msRyu.jpg}}]{Myeongseok Ryu}
    received the B.S. degree in Mechanical Engineering from Incheon National University, South Korea, in 2023, and the M.S. degree in Mechanical Engineering from GIST, South Korea, in 2025. 
    He is currently pursuing the Ph.D. degree in the Cho Chun Shik Graduate School of Mobility at Korea Advanced Institute of Science and Technology (KAIST), South Korea.
    His research interests include neuro-adaptive control, online optimization, and contraction theory.
    % \vspace{-10mm}
\end{IEEEbiography}

\DIFdelbegin %DIFDELCMD < \begin{IEEEbiography}[{
%DIFDELCMD <     \includegraphics[width=1in,height=1.25in,clip,keepaspectratio]{\template/people/dhHong.jpg}}]{Donghwa Hong}
%DIFDELCMD <     %%%
\DIFdelend \DIFaddbegin \begin{IEEEbiography}[{

\includegraphics[width=1in,height=1.25in,clip,keepaspectratio]{\template/people/dhHong.jpg}}]{Donghwa Hong}
    \DIFaddend received his B.S. degree in Physics and Engineering Physics from Yonsei University, Wonju, South Korea, in \DIFdelbegin \DIFdel{2023. He is currently pursuing }\DIFdelend \DIFaddbegin \DIFadd{2023, and }\DIFaddend the M.S. degree in Mechanical \DIFaddbegin \DIFadd{and Robotics }\DIFaddend Engineering from GIST, Gwangju, South Korea\DIFdelbegin \DIFdel{.  
}\DIFdelend \DIFaddbegin \DIFadd{, in 2026. Following his M.S., he is working in the Cho Chun Shik Graduate School of Mobility at KAIST as a researcher. }\DIFaddend His research interests include \DIFdelbegin \DIFdel{robot control, neural-networks, }\DIFdelend \DIFaddbegin \DIFadd{online learning }\DIFaddend and model identification. 
    % \vspace{-10mm}
\end{IEEEbiography}

\begin{IEEEbiography}[{
    \includegraphics[width=1in,height=1.25in,clip,keepaspectratio]{\template/people/khChoi.jpg}}]{Kyunghwan Choi}
    received his B.S., M.S., and Ph.D. degrees in Mechanical Engineering from KAIST, Daejeon, South Korea, in 2014, 2016, and 2020, respectively. Following his Ph.D., he worked at the Center for Eco-Friendly \& Smart Vehicles at KAIST as a Postdoctoral Fellow and later as a Research Assistant Professor. In 2022, he joined GIST as an Assistant Professor in the Department of Mechanical and Robotics Engineering. Currently, he serves as an Assistant Professor at the Cho Chun Shik Graduate School of Mobility at KAIST, where he is also the Director of the Mobility Intelligence and Control Laboratory. His research focuses on optimal and learning-based control for connected, automated, and electrified vehicles (CAEVs).
\end{IEEEbiography}

\end{document}